\documentclass[JEP,XML,SOM,Unicode,NoEqCountersInSection,published]{cedram}
\datereceived{2021-08-12}
\dateaccepted{2022-04-11}
\dateepreuves{2022-04-25}

\usepackage[scr=rsfs,cal=euler]{mathalfa}
\multlinegap0pt

\newcommand{\RedefinitSymbole}[1]{%
\expandafter\let\csname old\string#1\endcsname=#1
\let#1=\relax
\newcommand{#1}{\csname old\string#1\endcsname\,}%
}
\RedefinitSymbole{\forall} \RedefinitSymbole{\exists}
\newcommand{\Psfrac}[2]{(\sfrac{#1}{#2})}
\newcommand{\psfrac}[2]{\sfrac{(#1)}{#2}}
\newcommand{\spfrac}[2]{\sfrac{#1}{(#2)}}
\newcommand{\pspfrac}[2]{\sfrac{(#1)}{(#2)}}
\def\mto{\mathchoice{\longmapsto}{\mapsto}{\mapsto}{\mapsto}}

\newenvironment{enumeratei}
{\bgroup\def\theenumi{\roman{enumi}}\def\theenumii{\arabic{enumii}}\begin{enumerate}}
{\end{enumerate}\egroup}

\newcommand{\sbullet}{{\scriptscriptstyle\bullet}}

\theoremstyle{theorem}
\newtheorem{theo}{Theorem}
\newtheorem{prop}[theo]{Proposition}
\newtheorem{coro}[theo]{Corollary}
\newtheorem{lemm}[theo]{Lemma}

\theoremstyle{definition}
\newtheorem{defi}[theo]{Definition}
\newtheorem{rema}[theo]{Remark}

\newcommand\Deltag{\mathbf{\Delta}}
\newcommand\R{\mathbb{R}}
\newcommand\bbC{\mathbb{C}}\let\C\bbC
\newcommand\K{\mathbb{K}}
\newcommand\E{\mathbf{E}}
\newcommand\bfP{\mathbf{P}}\let\P\bfP
\newcommand\N{\mathbb{N}}
\newcommand\Z{\mathbb{Z}}
\newcommand\Q{\mathbb{Q}}
\newcommand\CC{\mathcal{C}}
\newcommand\bbS{\mathbb{S}}
\newcommand\T{\mathbb{T}}
\newcommand\ep{\varepsilon}
\newcommand\pun{\mathbf{1}}

\newcommand{\eq}[1]{\mathrel{\mathop{\kern0pt\sim}_{#1}}}
\newcommand\eqn{\eq{n\to +\infty}}
\newcommand{\eqs}[1]{\mathrel{\mathop{\kern0pt\simeq}_{#1}}}
\newcommand\eqsn{\eqs{n\to +\infty}}

\datepublished{2022-04-27}
\begin{document}
\frontmatter
\title[Condition for probabilistic continuity on a compact manifold]{A necessary and sufficient condition for probabilistic continuity on a boundaryless compact Riemannian manifold}

\alttitle{Condition nécessaire et suffisante pour la continuité probabiliste sur une variété riemannienne compacte sans bord}

\author[\initial{R.} \lastname{Imekraz}]{\firstname{Rafik} \lastname{Imekraz}}
\address{La Rochelle Université, MIA, EA 3165\\
F-17031 La Rochelle, France}
\email{rafik.imekraz@univ-lr.fr}
\urladdr{https://pageperso.univ-lr.fr/rafik.imekraz/}

\thanks{ Partially supported by the ANR projects ESSED ANR-18-CE40-0028 and UNIRANDOM ANR-17-CE40-0008}

\subjclass{60G50, 60G15, 58J40, 46B09}

\keywords{Paley-Zygmund theorems, Laplace-Beltrami operator, compact manifolds, Gaussian processes}

\altkeywords{Théorèmes de Paley-Zygmund, opérateur de Laplace-Beltrami, variétés compactes, processus gaussiens}

\begin{abstract}
We give a necessary and sufficient condition for the uniform convergence of random series of eigenfunctions on a boundaryless compact Riemannian manifold. As a consequence, we generalize an estimate of Burq and Lebeau about the supremum of a random eigenfunction. Finally, we prove that our results are universal with respect to the random variables (this is a Riemannian analogue of a result of Marcus and Pisier), with respect to compact submanifolds and with respect to the Riemannian structure of the manifold. Our proofs rely on several tools like the Dudley-Fernique theorem, the Slepian comparison theorem and the semi-classical functional calculus for elliptic operators on compact manifolds.
\end{abstract}

\begin{altabstract}
Nous donnons une condition nécessaire et suffisante pour la convergence uniforme de séries aléatoires de fonctions propres sur une variété riemannienne compacte sans bord. Comme conséquence, nous généralisons une estimation de Burq et Lebeau concernant les bornes d'une fonction propre aléatoire. Finalement, nous prouvons que nos résultats sont universels par rapport aux variables aléatoires utilisées (il s'agit d'un analogue riemannien d'un résultat de Marcus et Pisier), par rapport aux sous-variétés compactes et à la structure riemannienne de la variété compacte. Nos preuves reposent sur plusieurs outils dont le théorème de Dudley-Fernique, le théorème de comparaison de Slepian et un calcul fonctionnel semi-classique pour les opérateurs elliptiques d'une variété compacte.
\end{altabstract}

\maketitle
\vspace*{-\baselineskip}
\tableofcontents
\mainmatter

\vspace*{-3\baselineskip}\mbox{}
\section{Introduction and statement of the main result}
Let $\mathcal{M}$ be a boundaryless compact Riemannian manifold of dimension $d\geq 2$, the goal of Theorem \ref{sol-u} is to give a necessary and sufficient condition to ensure the convergence, in the Banach space $\CC^0(\mathcal{M})$, of suitable random linear combinations of eigenfunctions on $\mathcal{M}$. The problem studied is posed in the paper of Tzvetkov \cite[p.\,XV\nobreakdash-6]{Tzvetkov4} and, as it will be explained in this introduction, is equivalent to control expectations of supremum of linear combinations of eigenfunctions. Thus we need to generalize an optimal two-sided inequality obtained by Burq and Lebeau in \cite{burq-lebeau,burq-lebeau2}, see \eqref{burqleb} and the new inequalities \eqref{newr} and \eqref{sol-u-eq}.
The results below are Riemannian analogues of those of the classical theory, for $\mathcal{M}$ being a compact group, whose final treatment is done in the book of Marcus and Pisier \cite{pisier1981}. And we moreover prove, via the semi-classical analysis, that all our results are independent of the Riemannian structure of $\mathcal{M}$ (see Theorem \ref{indep-met}).%\enlargethispage{.5\baselineskip}%

Our result really looks like a sufficient condition obtained by Salem and Zygmund in \cite[p.\,291]{salem1954}, the main new point is that our Salem-Zygmund condition is necessary and sufficient in the manifold framework (see \eqref{pz-con30} in Theorem \ref{sol-u}) whereas it is merely sufficient (and not necessary) in the classical theory on the torus (see \cite[p.\,292]{salem1954}). As underlined in the paper \cite{burq-lebeau}, a clue suggesting that better results should exist in the Riemannian setting is that the dimensions of the eigenspaces of the Laplace-Beltrami operator on $\bbS^d$, for $d\geq 2$, tend to $+\infty$ in contrast to the classical setting $L^2(\mathbb{T})=\bigoplus_{n\in \Z} \C e^{in\sbullet}$. The present contribution gives another enlightenment (see \hbox{below} Theorem \ref{sol-u} for definitions): it turns out that the Dudley pseudo-distance of the Gaussian random wave of each eigenspace of $\bbS^d$ is equivalent to a very simple distance (see Theorems \ref{pro-s+} and \ref{pseudo-dis}).

Before going into details on compact manifolds (and giving a precise definition of our random series), it is worthwhile to recall the main results on the torus $\mathbb{T}=\R/ 2\pi \Z$ (we refer to \cite{pisier1981} for the general case with compact groups). Although the first results had been proved for Fourier series with real trigonometric functions, it is known that there is no difference to deal with complex trigonometric functions (see \cite[p.\,122]{pisier1981}). For any real sequence $(c_n)_{n\in \N}\in \ell^2(\N)$, let us consider the function $f\in L^2(\mathbb{T})$ defined by
\begin{equation*}
f(x)=\sum_{n\in \N}c_n e^{inx}.
\end{equation*}
Now fix a sequence $(\varepsilon_n)_{n\in \N}$ of i.i.d. Rademacher random variables, in other words $\P(\varepsilon_n=1)=\P(\varepsilon_n=-1)=\sfrac{1}{2}$.
With probability $1$ with respect to $\omega$ running over a reference probability space $\Omega$, Paley and Zygmund proved in \cite[p.\,347]{paley1930} that the following random series
\begin{equation}\label{pz}
f^\omega(x):=\sum_{n\in \N}\varepsilon_n(\omega) c_n e^{inx}
\end{equation}
almost surely converges in $\CC^0(\mathbb{T})$ provided that the following condition is fulfilled
\begin{equation}\label{pz-con}
\exists \gamma>1, \quad \sum_{n=2}^{+\infty} c_n^2 \ln^{\gamma}(n)<+\infty.
\end{equation}
The strong motivation of the resurgence of conditions like \eqref{pz-con} is that they can be interpreted as probabilistic Sobolev embeddings (more details are given in the introductions of \cite{Tzvetkov4,burq-lebeau} and \cite[p.\,2734]{imek-aif}).
In \cite[p.\,291]{salem1954}, Salem and Zygmund relaxed the assumption \eqref{pz-con} by the following one
\begin{equation}\label{pz-con2}
\sum_{n=2}^{+\infty} \frac{1}{n\sqrt{\ln(n)}}\biggl(\sum_{k=n}^{+\infty}c_k^2 \biggr)^{1/2} <+\infty,
\end{equation}
and also remarked that such a condition is not necessary.
An important step is made by Marcus in the papers \cite{marcus1973,marcus1975} in which we learn that \eqref{pz-con2} becomes a necessary and sufficient condition if the sequence $(|c_n|)_{n\in \N}$ in \eqref{pz} is non-increasing and if the Rademacher random variables $\ep_n$ are replaced with a sequence $(g_n^{\C})_{n\in \N}$ of independent complex standard Gaussian random variables\footnote{In particular we have $\E[g_n^\C]=0$ and $\E[|g_n^\C|^2]=1$.} $\mathcal{N}_\mathcal{C}(0,1)$ (see \cite{marcus1975} for general complex symmetric random variables satisfying suitable normalization conditions).

For the problem of finding a necessary and sufficient condition, the solution had finally come from a drastically different point of view. More precisely, for any sequence $(c_n)\in \ell^2(\Z)$, the general Gaussian Fourier series
\begin{equation}\label{ran-fou}
f^{G,\omega}(x)=\sum_{n\in \Z} g_n^{\C}(\omega) c_n e^{inx}
\end{equation}
is now seen as a stationary Gaussian random process on the torus $\mathbb{T}$.
Here the word ``stationary'' means that the random process $(f^{G}(x))_{x\in \mathbb{T}}$ is invariant under the group action of $\mathbb{T}$: in particular, for any angle $\alpha \in \R$, $f^{G,\omega}(x+\alpha)$ has the same distribution as that of $f^{G,\omega}(x)$ due to the complex symmetry of the complex Gaussian variables~$g_n^\C$.
It is worthwhile to underline that such a stationary assumption has no sense for a general compact manifold $\mathcal{M}$.
Then the important result, now called the Dudley-Fernique theorem, allows for proving that the almost sure continuity of the Gaussian functions
$x\mto f^{G,\omega}(x)$ is equivalent to the so-called entropy condition:
\begin{equation}\label{entrop}
\int_{0}^{+\infty} \sqrt{\ln N_{\delta}(\T,\ep)} d\ep<+\infty,
\end{equation}
\enlargethispage{-2\baselineskip}%
%\pagebreak[2]
where
\begin{itemize}
\item the function $\delta:\mathbb{T}\times \mathbb{T}\to \R^+$ is the Dudley pseudo-distance of $\mathbb{T}$ defined by\vspace*{-3pt}
\begin{equation*}
\delta(x,y)^2 =\E\left[| f^{G,\omega}(x)-f^{G,\omega}(y)|^2\right]=4\sum_{n\in \Z} |c_n|^2 \sin^2\Bigl(\frac{n(x-y)}{2}\Bigr),
\end{equation*}
\item the number $N_{\delta}(\T,\ep)$ is the covering number of the torus $\mathbb{T}$ with respect to $\delta$, namely the minimal number of open $\delta$-balls of radius $\ep>0$ whose union covers $\mathbb{T}$.
\end{itemize}
Moreover, Marcus and Pisier proved that one may replace the sequence of Gaussian random variables in \eqref{ran-fou} with any sequence of independent real centered random variables $(X_n)_{n\in \Z}$ satisfying\vspace*{-3pt}
\begin{equation*}
0<\inf_{n\in \Z}\E[|X_n|] \quad \text{and} \quad \sup_{n\in \Z} \E[|X_n|^2]<+\infty.
\end{equation*}
Such a result is usually called a universality phenomenon and allows for coming back to the Rademacher random variables $X_n=\ep_n$ (see \cite[p.\,7--9]{pisier1981} and \cite[p.\,28, Cor.\,7.3]{pisier1977}).

One may think that the entropy condition \eqref{entrop} is a bit abstract but it really captures the essence of the almost sure continuity.
Firstly, the entropy condition allows for showing the sufficiency of the more concrete Paley-Zygmund \eqref{pz-con} and Salem-Zygmund~\eqref{pz-con2} assumptions (see \cite[Part VII.1]{pisier1981}).
Secondly, the entropy condition (which is of qualitative nature) is equivalent to the following quantitative version (see \cite[p.\,11]{pisier1981} or \cite[p.\,3--4]{pisier1977}):\vspace*{-3pt}
\begin{equation}\label{equiv-ent}
\begin{aligned}
\E_\omega\biggl[\Big\|\sum_{k=-n}^n \ep_k(\omega) c_k e^{ikx} \Big\|_{\CC_x^0(\mathcal{\mathbb{T}})}\biggr] & \simeq \E_\omega\biggl[\Big\|\sum_{k=-n}^n g_k^{\C}(\omega) c_k e^{ikx} \Big\|_{\CC_x^0(\mathcal{\mathbb{T}})}\biggr]\\ 
& \simeq |c_0|+\int_{0}^{+\infty} \sqrt{\ln N_{\delta}(\T,\ep)} d\ep.
\end{aligned}
\end{equation}
For instance, here are the so-called Salem-Zygmund inequalities (see \cite[p.\,259]{queff2}):\vspace*{-3pt}
\begin{equation}\label{salemz}
\forall n\gg 1, \qquad \E_{\omega}\biggl[\frac{1}{\sqrt{n}}\Bigl\Vert \sum_{k=n+1}^{2n}\ep_k(\omega) e^{ikx} \Bigr\Vert_{\CC_x^0(\mathbb{T})} \biggl]\simeq \sqrt{\ln(n)}.
\end{equation}
For the sequel, it is worthwhile to note that the Dudley pseudo-distance $\delta_n$ of\vspace*{-3pt}
\begin{equation*}
\frac{1}{\sqrt{n}}\sum_{k=n+1}^{2n}g_k^{\C}(\omega) e^{ikx}
\end{equation*} satisfies the following estimates\footnote{
By setting $\theta=\psfrac{x-y}{2}$, one checks that \eqref{salemz2} comes from the exact formulas $\delta_n(x,y)^2 = \Psfrac{4}{n}\sum_{k=n+1}^{2n} \sin^2\left(k\theta\right)=2-\Psfrac{2\sin(n\theta)}{n\sin(\theta)}\cos((3n+1)\theta)$ by separating the cases $0\leq \theta\leq \sfrac{2}{n}$ and $\sfrac{2}{n}\leq \theta\leq \sfrac{\pi}{2}$.}
\begin{equation}\label{salemz2}
\forall n\gg1, \quad \delta_n\simeq \min(1,n\delta_{g}), \quad \text{where $\delta_{g}$ is the Riemannian distance of }\mathbb{T},
\end{equation}
which in turn allows to show the implication $\eqref{equiv-ent}\Rightarrow\eqref{salemz}$.

It is time to recall the known literature about the generalization of the previous results to a boundaryless compact Riemannian manifold $\mathcal{M}$ of dimension $d\geq 2$.
We~denote by $\Deltag$ the non-positive Laplace-Beltrami operator of $\mathcal{M}$.
The Hilbert space $L^2(\mathcal{M})$ is considered with respect to the Riemannian measure $\operatorname{vol}_{\mathcal{M}}$ of $\mathcal{M}$.
We recall that $L^2(\mathcal{M})$ admits a Hilbert basis $(\phi_k)_{k\in \N}$ of eigenfunctions of $\Deltag$:
\begin{equation}\label{hilb-bas}
\Deltag \phi_k=-\lambda_k^2 \phi_k,\qquad 0=\lambda_0<\lambda_1\leq \lambda_2 \leq \dots \to +\infty.
\end{equation}
For simplicity, we shall assume that each $\phi_k$ is real-valued.
By fixing a sequence of coefficients $(c_k)_{k\in \N}$, a natural option is to replace the random Fourier series \eqref{pz} and~\eqref{ran-fou} with the following random series:
\begin{equation}\label{ran-fou2}
\sum_{k\in \N} \ep_k(\omega) c_k \phi_k(x) \quad \text{or}\quad \sum_{k\in \N} g_k(\omega) c_k \phi_k(x),\qquad (\omega,x)\in \Omega\times \mathcal{M},
\end{equation}
where $(g_k)_{k\in \N}$ is a sequence of i.i.d. real Gaussian random variables $\mathcal{N}_{\R}(0,1)$.
For the random series \eqref{ran-fou2}, Tzvetkov generalized in \cite{Tzvetkov4} the Paley-Zygmund condition~\eqref{pz-con}, that is a sufficient condition ensuring the almost sure convergence in $\CC^0(\mathcal{M})$. In contrast with the sequence of eigenvalues $(\lambda_k^2)_{k\geq 0}$, the sequence of eigenfunctions $(\phi_k)_{k\geq 0}$ may not be unique (actually if $\Deltag$ has multidimensional eigenspaces). For instance, on $\mathcal{M}=\bbS^d$, there are sub-sequences of eigenfunctions having very different behaviors like concentration around a point or around a geodesic (those are called zonal eigenfunctions or Gaussian beams).
Without any further information on the sequence $(\phi_k)_{k\in \N}$, it seems hopeless to expect to have a simple, necessary and sufficient condition for the almost sure convergence in $\CC^0(\mathcal{M})$ of \eqref{ran-fou2}.
By comparison with the classical theory on a compact group $G$, the adequate random series are defined by taking account of the irreducible representations of $G$ (see \cite{figa1967,pisier1981}). In the Riemannian framework, it is thus natural to make an additional restriction that counterbalances the possible concentration of sub-sequences of eigenfunctions.
The option chosen by Burq and Lebeau in \cite{burq-lebeau} will be detailed below and relies in gathering eigenfunctions whose eigenvalues are in a same suitable small interval (for instance for the same eigenvalue for $\mathcal{M}=\bbS^d$, see \eqref{sd-enn}).
We also refer to \cite[p.\,272]{imek-crusep} and also \cite{randomh,imek-aif}
for more details about that multidimensional point of view.

Before writing the slight modification of \eqref{ran-fou2},
let us underline that the situation is drastically different for the $L^p$ case with $p<+\infty$. Indeed, there is a very nice formula of Maurey that gives a necessary and sufficient condition for the almost sure convergence of unidimensional series like \eqref{pz} and \eqref{ran-fou2} in $L^p$ for $p<+\infty$ (see \cite[p.\,22, Cor.\,1]{maurey1973}, \cite[Th.\,1.d.6]{linden2} or \cite[Prop.\,2.1]{randomh} for a recent use).
For a multidimensional analogue of the Maurey formula, we refer to \cite[Th.\,2.1 \& 2.21]{imek-crusep} and \cite[Th.\,12.1]{imek-aif}.
Without going into details on the difficulties of the $L^p$ case for $p<+\infty$, one could say that the analogue problems of finding necessary and sufficient conditions on $L^p$ are quite well understood for finite $p$ (see the papers \cite{tzvetkov-ay,Tzvetkov4,Grivaux,randomh,imek-crusep} and \cite[Th.\,2.3, 4.5 \& 4.6]{imek-aif}).

It is time to properly define the random series that naturally replace \eqref{ran-fou2} in the case of the sphere $\bbS^d$ (the case of manifold is considered just after).
We recall that the sequence of eigenvalues of the Laplace-Beltrami $\Deltag$ on the sphere $\bbS^d$ is given, without counting multiplicities, by $-n(n+d-1)$ with $n\in \N$. Moreover, the dimension of the eigenspace
\begin{equation}\label{sd-enn}
E_n=\ker(\Deltag+n(n+d-1)) \subset L^2(\bbS^d)
\end{equation}
satisfies the numerical equivalence $\dim(E_n)\simeq n^{d-1}$ for $n\geq 1$.
For any function $f\in L^2(\bbS^d)$, we write
\begin{equation*}
f=\sum_{n\in \N}f_n \quad \text{with } f_n=\sum_{\lambda_k^2=n(n+d-1)} \langle f, \phi_k\rangle \phi_k \in E_n.
\end{equation*}
For any $n\in \N^\star$, we now consider a uniform random vector $U_n:\Omega\to \R^{\dim(E_n)}$, namely whose probability distribution is the
probability spherical measure on the unit sphere $\bbS^{\dim(E_n)-1}$.
We write the coordinates $U_n=(U_{n,k})_{k}$ with $k$ running over the finite set of integers such that
$\lambda_k^2=n(n+d-1)$. A natural generalization of \eqref{pz} and the first random series in \eqref{ran-fou2} is given by	
\begin{equation}\label{ran-fou4}
f^\omega(x)=\sum_{n\geq 1} f_n^\omega(x), \quad \text{with } f_n^\omega=
\left\Vert f_n\right\Vert_{L^2(\bbS^d)} \sum_{\lambda_k^2=n(n+d-1)} U_{n,k}(\omega) \phi_k \in E_n,
\end{equation}
where the uniform random vectors $U_n$ are assumed to be mutually independent.
The random function $f_n^\omega$ is thus a random eigenfunction of norm $\|f_n\|_{L^2(\bbS^d)}$.
That formalism is connected to that used by Burq-Lebeau in \cite[App.\,C]{burq-lebeau} (see \cite[p.\,274]{imek-crusep} for more details).

On a general boundaryless compact Riemannian manifold $\mathcal{M}$, the idea is to replace the eigenspace $E_n\subset L^2(\bbS^d)$ with the subspace $E_{(Kn-K,Kn]}$
of $L^2(\mathcal{M})$ defined as follows
\begin{equation} \label{defiEn}
E_{(Kn-K,Kn]}:= \operatorname{Span}\{\phi_k,\ \lambda_k\in (Kn-K,Kn]\},
\end{equation}
where the spectral parameter $K>0$ is large enough. In particular, it is known that the analogue of the asymptotic $\dim(E_n)\simeq n^{d-1}$ is given by $\dim (E_{(Kn-K,Kn]})\simeq n^{d-1}$ (see \cite[p.\,923]{burq-lebeau} or \cite[Lem.\,8.1]{imek-aif}).
In other words, for any $f=\sum_{n\geq 1}f_n \in L^2(\mathcal{M})$ with $f_n\in E_{(Kn-K,Kn]}$, we set the following random series
\begin{equation} \label{ran-fouu}
f^\omega = \sum_{n\geq 1}f_n^\omega,\quad \text{with } f_n^{\omega}:=\| f_n\|_{L^2(\mathcal{M})} \sum_{\lambda_k \in (Kn-K,Kn]} U_{n,k}(\omega) \phi_k,
\end{equation}
where $U_n:\Omega\to \R^{\dim(E_{(Kn-K,Kn]})}$ is a uniform random vector as above.
At this stage of this introduction, we must recall an optimal result proved by Burq and Lebeau.
By~using the previous notations, \cite[Th.\,5, p.\,930]{burq-lebeau} states the following optimal bound on a uniform random eigenfunction $f_n^\omega$ on $L^2(\bbS^d)$ with eigenvalue $-n(n+d-1)$:%
\begin{equation}\label{burqleb}
\forall n\gg 1, \quad \E_{\omega}\left[\left\Vert f_n^\omega \right\Vert_{\CC^0(\bbS^d)} \right]\simeq \sqrt{\ln(n)} \left\Vert f_n \right\Vert_{L^2(\bbS^d)}.
\end{equation}
We also refer to \cite[Th.\,3]{canzani2014high} for a geometric control of the upper bound in \eqref{burqleb}.
The idea we keep in mind is that the inequalities \eqref{burqleb} of Burq-Lebeau should really be seen as a Riemannian analogue of the Salem-Zygmund inequalities \eqref{salemz}.
\setlength{\unitlength}{1cm}
\begin{center}
\begin{picture}(12.5,3.3)
\put(7,2){\framebox(5.5,1){\begin{tabular}{cc} The Salem-Zygmund\\ inequalities \eqref{salemz} on $\mathbb{T}$ \end{tabular}}}
\put(7,0.5){\framebox(5.5,1){\begin{tabular}{cc}The Burq-Lebeau inequalities \eqref{burqleb} \\ on $\bbS^d$ (with $d\geq 2$) \end{tabular} }}
\put(0,2){\framebox(5.5,1){\begin{tabular}{cc}Classical theory \\on a compact group\end{tabular}}}
\put(0,0.5){\framebox(5.5,1){\begin{tabular}{cc}Problem posed on a compact \\ Riemannian manifold (of $\dim\geq 2$) \end{tabular}}}
\put(5.7,2.5){\vector(1,0){1}}
\put(5.7,1){\vector(1,0){1}}
\end{picture}
\vspace*{-\baselineskip}
\end{center}
Before stating our main result, we recall the striking equivalence of the following two statements:
\begin{enumeratei}
\item for almost every $\omega\in \Omega$, the random series $\sum f_n^\omega$ converges in $\CC^0(\bbS^d)$,
\item the random series $\sum f_n^\omega$ converges in $L^1(\Omega,\CC^0(\bbS^d))$.
\end{enumeratei}
For unidimensional Rademacher random series $\sum \ep_n(\omega)f_n$, an analogue equivalence of the previous assertions was proved by Kahane and is indeed completely independent of the Banach space $\CC^0(\bbS^d)$ (see the reference book \cite{kahane} about random series, \cite[p.\,142, Rem.\,1]{queff1} or \cite[p.\,43]{pisier1981}).
For the multidimensional case we are interested in, such an equivalence is a consequence of a result by Marcus and Pisier \cite[p.\,92]{pisier1981}.
We now understand that a general solution of our problem is equivalent to find optimal bounds of
the expectations
$\E_\omega\big[\| \sum_{n=1}^N f_n^\omega\|_{\CC^0(\bbS^d)}\big]$
generalizing the Burq-Lebeau asymptotics \eqref{burqleb}. Our work will show the following new result on~$\bbS^d$ (with $d\geq 2$):
\begin{equation}\label{newr}
\E_\omega\biggl[\Big\| \sum_{n=1}^N f_n^\omega \Big\|_{\CC^0(\bbS^d)}\biggr]\simeq \sum_{p=1}^N \frac{1}{p\sqrt{\ln(p+1)}}\, \biggl(\sum_{n=p}^N \|f_n\|_{L^2(\bbS^d)}^2 \biggr)^{\sfrac{1}{2}}.
\end{equation}
The right-hand side clearly looks like \eqref{pz-con2} but no monotonicity assumption is needed (in~contrast to \cite{marcus1975} for the torus $\mathbb{T}$).
More generally, Theorem \ref{sol-u} settles the general case of a compact manifold with a Salem-Zygmund type condition \eqref{pz-con30} which is necessary and sufficient in contrast to the classical results (for which analogue estimates of \eqref{pz-con30} are sufficient but not necessary).

\begin{theo}\label{sol-u}
There is a constant $K_0>0$ depending only on the Riemannian manifold $\mathcal{M}$ such that, for any $K\geq K_0$, if one considers
\begin{itemize}
\item a sequence $(f_n)_{n\geq 1}$ satisfying $f_n\in E_{(Kn-K,Kn]}$ (see \eqref{defiEn}) for each $n\in \N^\star$,
\item a non-zero dimensional compact submanifold $\mathcal{M}_s\subset \mathcal{M}$ with smooth (eventually empty) boundary,
\item a sequence of independent real random variables\footnote{Note that we can choose $X_n=1$ for each $n\geq 1$ in order to get \eqref{newr}.} $(X_n)_{n\geq 1}$ satisfying
\begin{equation}\label{hyp-Xn}
0<\inf_{n\geq 1}\E[|X_n|] \quad \text{and}\quad \sup_{n\geq 1}\E[|X_n|^2]<+\infty,
\end{equation}
\end{itemize}
\pagebreak[2]
then
\begin{enumerate}
\item we have the numerical equivalence for any $N\in \N^\star$:\vspace*{-3pt}
\begin{equation}\label{sol-u-eq}
\E_{\omega}\biggl[\sup_{x\in \mathcal{M}_s}\Big| \sum_{n=1}^{N} X_n(\omega) f_n^{\omega}(x) \Big| \biggr] \simeq \sum_{p=1}^N \frac{1}{p\sqrt{\ln(p+1)}}\,\biggl(\sum_{n=p}^N \left\Vert f_n\right\Vert_{L^2(\mathcal{M})}^2 \biggr)^{\sfrac{1}{2}},
\end{equation}
where the constants of that equivalence are independent of $N$ and depend on $\mathcal{M}$, $\mathcal{M}_s$, $K$, $\inf_{n\geq 1} \E[|X_n|]$ and $\sup_{n\geq 1} \E[|X_n|^2]$
and where we assume that all the random variables $\omega\mto X_n(\omega)$ and $\omega\mto f_n^\omega$ (defined in \eqref{ran-fouu}) are mutually independent.

\item Moreover, the following two statements are equivalent:
\begin{itemize}
\item the random series $\sum X_n(\omega)f_n^{\omega}$ almost surely converges in $\mathcal{C}^0(\mathcal{M}_s)$,
\item the following Salem-Zygmund condition is fulfilled:\vspace*{-3pt}
\begin{equation}\label{pz-con30}
\sum_{p=1}^{+\infty} \frac{1}{p\sqrt{\ln(p+1)}}\,\biggl(\sum_{n=p}^{+\infty} \left\Vert f_n\right\Vert_{L^2(\mathcal{M})}^2 \biggr)^{\sfrac{1}{2}}<+\infty.
\end{equation}
\end{itemize}
\item Finally, the same conclusion also holds true for the particular case $\mathcal{M}=\bbS^d$ if each $f_n$ belongs to $\ker(\Deltag+n(n+d-1))$
provided that we consider \eqref{ran-fou4} instead of~\eqref{ran-fouu} (and thus $K$ is irrelevant in the equivalence \eqref{sol-u-eq}).
\end{enumerate}
\end{theo}
Here are two comments about the previous result:
\begin{itemize}
\item as a consequence of Point 2) of Theorem \ref{sol-u}, one sees that the almost sure convergence in $\mathcal{C}^0(\mathcal{M})$ is equivalent to the almost sure convergence in $\mathcal{C}^0(\mathcal{M}_s)$ for any submanifold $\mathcal{M}_s$ (more comments in that direction are given in Section \ref{univer}),
\item Theorem \ref{sol-u} supersedes \cite[Th.\,2.1]{imek-aif} in which one may find a sufficient condition in the spirit of the original Paley-Zygmund theorem.
\end{itemize}

Similarly to the theory developed by Marcus and Pisier, the good idea is first to consider our problem with a suitable Gaussian analogue of $\sum f_n^\omega$.
A known intuition suggests that the uniform random vector $U_n:\Omega\to \bbS^{\dim(E_n)-1}$, in \eqref{ran-fou4}, is closely related to the Gaussian vector of $\R^{\dim(E_n)}$ with distribution $\mathcal{N}\bigl(0,\spfrac{1}{\dim(E_n)}\operatorname{Id}\bigr)$ (see \cite[p.\,58]{pisiervol}). We then introduce the following Gaussian analogue of \eqref{ran-fou4} for any $f\in L^2(\bbS^d)$\vspace*{-3pt}
\begin{equation}\label{ran-fouG}
f^{G,\omega}:=\sum_{n\geq 1} f_n^{G,\omega}, \quad \text{with } f_n^{G,\omega}= \frac{\left\Vert f_n\right\Vert_{L^2(\bbS^d)}}{\sqrt{\dim(E_n)}}\sum_{\lambda_k^2=n(n+d-1)} g_{n,k} (\omega)\phi_k,
\end{equation}
that must be seen as a Gaussian process on the manifold $\bbS^d$ (where $(g_{n,k})_{(n,k)}$ is a sequence of independent standard Gaussian random variables $\mathcal{N}_\R(0,1)$ with $n\in \N^\star$ and $\lambda_k^2=n(n+d-1)$).
On a general boundaryless compact Riemannian manifold\,$\mathcal{M}$, the Gaussian analogue of \eqref{ran-fouu} should be\vspace*{-3pt}
\begin{equation} \label{ran-fouG2}
f^{G,\omega} = \sum_{n\geq 1}f_n^{G,\omega}, \quad \text{with } f_n^{G,\omega}:=\frac{\left\Vert f_n\right\Vert_{L^2(\mathcal{M})}}{\sqrt{\dim(E_{(Kn-K,Kn]})}}\sum_{\lambda_k \in (Kn-K,Kn]}\hspace*{-3mm} g_{n,k}(\omega) \phi_k.
\end{equation}
The random functions $f_n^{G,\omega}$ are usually called Gaussian random waves and their study had been introduced by Zelditch in \cite{zelditch2009real}.
Let us explain another reason for replacing~$f_n^\omega$ with $f_n^{G,\omega}$. The rotational invariance of Gaussian vectors implies the following distribution equivalence (see \cite[p.\,2731]{imek-aif}):
\begin{equation}\label{chi2a}
f_n^{G,\omega} \sim \frac{1}{\sqrt{\dim(E_{(Kn-K,Kn]})}}\,\biggl(\sum_{\lambda_k\in (Kn-K,Kn]} g_{n,k}^2(\omega)\biggr)^{1/2} f_n^\omega,
\end{equation}
where all the random variables involved are assumed to be mutually independent.
As~a consequence of \eqref{chi2a}, one may directly replace $f_n^{\omega}$ with $f_n^{G,\omega}$ for getting bounds of $\E_\omega[\left\Vert f_n^\omega \right\Vert_{\CC^0(\mathcal{M})}]$ (see Appendix \ref{chi2}):
\begin{align*}\nonumber
\E_{\omega}\bigl[\Vert f_n^{G,\omega}&\Vert_{\CC^0{(\mathcal{M})}} \bigr]\\
& = \frac{1}{\sqrt{\dim(E_{(Kn-K,Kn]})}}\,\E\biggl[\biggl(\sum_{Kn-K<\lambda_k\leq Kn} g_{n,k}^2\biggr)^{1/2} \biggr]\times \E_\omega\Bigl[\left\Vert f_n^{\omega}\right\Vert_{\CC^0{(\mathcal{M})}} \Bigr]\\
& \simeq \E_{\omega}\Bigl[\left\Vert f_n^{\omega}\right\Vert_{\CC^0{(\mathcal{M})}} \Bigr].
\end{align*}
Our results will show that the previous equivalence can be generalized to linear combinations of $f_n^\omega$ and $f_n^{G,\omega}$ (that is a Riemannian analogue of a result by Marcus and Pisier, see Section \ref{univer}).

Let us now explain how to deal with \eqref{ran-fouG} or \eqref{ran-fouG2}.
In the classical theory, the main ingredient is the fact that the Gaussian random process \eqref{ran-fou} is stationary and thus one may use the Dudley-Fernique theorem giving a complete understanding of its almost sure continuity.
It seems to be reasonable, at least for the Gaussian random series \eqref{ran-fouG} on spheres $\bbS^d$, that such ideas can be combined with the transitivity of the isometry group (see an argument at the end of Section 2).
It is however clear that such arguments are no longer sufficient for \eqref{ran-fouG2} on a general compact Riemannian manifold~$\mathcal{M}$ since that case does not seem to fulfill any stationary assumption. However, without a stationary assumption, the Dudley theorem \cite{dudley1967} gives a sufficient condition for the almost sure continuity. We refer to the work of Tzvetkov \cite{Tzvetkov4} for a use of the Dudley theorem in the spirit of the Paley-Zygmund theorem (and also \cite[Annex\,6]{randomh} for an adaptation of the argument of Tzvetkov).

In order to obtain a necessary and sufficient condition, there are several angles of attack.
The first way the author took was to use a generalization of the Dudley-Fernique theorem that weakens the stationary assumption (actually a minor variant of \cite[Th.\,2.7.4, p.\,61]{dudley2014}). As pointed out by a referee, one may simplify such a strategy by using the Slepian comparison theorem
(see \cite[Cor.\,3.14]{ledoux} or \cite[p.\,73]{queff2}).
That is indeed such a proof that will be followed in the sequel: one first completely understands the Gaussian process \eqref{ran-fouG} on the sphere $\bbS^d$ for $d\geq 2$ (via the Dudley-Fernique theorem) and then the Slepian comparison theorem will show that the almost sure continuity of \eqref{ran-fouG} is equivalent to that of \eqref{ran-fouG2} for the manifold setting. In both cases (submanifolds of $\bbS^d$ or of a general manifold), we will obtain sharp estimates of so-called Dudley pseudo-distances whose definition is now recalled.

A fundamental idea of the paper \cite{dudley1967} of Dudley can be summarized as follows: the almost sure properties of the Gaussian process $(f^{G,\omega}(x))_{x\in \mathcal{M}}$ are closely related to the pseudo-distance $\delta$ given by
\begin{equation}\label{delt}
\forall (x,y)\in \mathcal{M}^2, \quad \delta(x,y):=\sqrt{\E_{\omega}[|f^{G,\omega}(x)-f^{G,\omega}(y)|^2]}.
\end{equation}
Except in very specific examples, one cannot expect to have a simple formula for $\delta(x,y)$.
In our concrete situation, the pseudo-distance $\delta$ will be equivalent to a simpler distance or pseudo-distance which is a function of the original distance $\delta_g$.
An easy computation in \eqref{delt} indeed shows the formula
\begin{equation*}
\delta(x,y)^2=\sum_{n\geq 1} \left\Vert f_n \right\Vert_{L^2(\mathcal{M})}^2 \delta_n(x,y)^2,
\end{equation*}
where the partial pseudo-distances $\delta_n$ are given by
\begin{equation}\label{intro-dn}
\delta_n(x,y)^2:=\frac{1}{\dim(E_{(Kn-K,Kn]})} \sum_{\lambda_k \in (Kn-K,Kn]} |\phi_k(x)-\phi_k(y)|^2.
\end{equation}
In another context, a study of those partial pseudo-distances is done by Canzani and Hanin in \cite{canzani2017c} with specific geometric assumptions on the Riemannian manifold~$\mathcal{M}$.
In order to dispense with any geometric assumption on $\mathcal{M}$, we shall combine the following microlocal tools and ideas:
\begin{itemize}
\item a part of the work of Canzani and Hanin (in this case Proposition \ref{can-han} below coming from \cite[Lem.\,5]{canzani2015scaling}) about the derivative of the spectral function on a boundaryless compact manifold,
\item choosing $K$ large enough (such an idea is due to Burq and Lebeau in \cite{burq-lebeau}) and using off-diagonal estimates obtained by Hörmander (see the proof of Proposition \ref{far} below).
\end{itemize}
As a consequence, we will prove the equivalence $\delta_n\simeq \min(1,n\delta_g)$ where $\delta_g$ is the Riemannian distance of $\mathcal{M}$ (that is a similar form to \eqref{salemz2}). Roughly speaking, such an equivalence is possible because concentration of individual eigenfunctions does not matter. The simplicity of such an equivalence is the reason allowing us to simplify the entropy integral in order to recover the Salem-Zygmund condition \eqref{pz-con30}.
This article is organized as follows:
\begin{itemize}
\item In Section \ref{result}, we state the two main theorems concerning the Gaussian random series \eqref{ran-fouG} on $\bbS^d$ and \eqref{ran-fouG2} on $\mathcal{M}$.
More precisely, Theorem \ref{sol} and Theorem \ref{sol-q} are respectively of qualitative and quantitative nature.
\item In Section \ref{univer}, we first discuss a trivial consequence of Theorem \ref{sol-u} about universality with respect to the choice of $X_n$ in $\sum X_n(\omega) f_n^\omega$ and with respect to the chosen submanifold $\mathcal{M}_s$.
Then we state Theorem \ref{indep-met} ensuring that our analysis is actually independent of the Riemannian metric initially chosen on $\mathcal{M}$. That result is proved in Section \ref{s-wrt} via the semi-classical analysis of the Laplace-Beltrami operator seen as an elliptic differential operator on the compact manifold $\mathcal{M}$.
\item Section \ref{met-con} is devoted to metric considerations giving a suitable reformulation of the entropy integral.
\item In Sections \ref{partial1} and \ref{partial2}, we prove that the partial pseudo-distances $\delta_n$ in \eqref{intro-dn} are, in some sense, equivalent to the explicit distance $\min(1,n\delta_g)$ (where $\delta_g$ stands for the Riemannian distance). The proofs need to make an accurate comparison with the Bessel function $J_{\sfrac{d}{2}-1}$. Actually, Theorem \ref{pro-s+} and Theorem \ref{pseudo-dis} are the main contributions of the paper that allow to understand the Gaussian processes \eqref{ran-fouG} and~\eqref{ran-fouG2}.
\item Section \ref{s-p2} is written for the sake of clarity and contains a few precisions about the Gaussian processes \eqref{ran-fouG} and \eqref{ran-fouG2}.
\item Sections \ref{preuv-solSd}, \ref{preuv-solq} and \ref{preuv-solM} are devoted to the proofs of the main results in the Gaussian case for spheres and compact manifolds.
As written above, we first deal with the sphere case $\bbS^d$ with $d\geq 2$. Few technical details are necessary because the equivalence $\delta_n\simeq \min(1,n\delta_g)$ does not hold in the whole sphere (for $n$ even, $\delta_n$ is not a distance on $\bbS^d$ and does not distinguish antipodal points). The Slepian comparison theorem is then used to deal with submanifolds of $\bbS^d$ or of a general compact manifold $\mathcal{M}$.
\item Sections \ref{sec-marc} and \ref{sec-marc2} contain the proof of Theorem \ref{sol-u} that deals with the initial random series $\sum X_n(\omega)f_n^\omega$.
We shall use a truncation argument (already present in the work of Marcus-Pisier) adapted here for the so-called $\chi$ distributions.
\item Section \ref{s-wrt} contains the proof of the invariance with respect to the Riemannian metric of $\mathcal{M}$ (proof of Theorem \ref{indep-met}). The proof makes use of the theory of semi-classical pseudo-differential operators.
\item Finally, we have gathered two appendices presenting either computations or proofs of more or less known results that we have not found in the literature in the form we need.
\end{itemize}

\subsubsection*{Acknowledgements} The author would like to thank Hervé Queffélec for discussions about majorizing measures.
Moreover, the author would like to thank the referees who studied the present paper and gave improvements of our initial proofs.
Actually, the reference \cite[Th.\,2.7.4, p.\,61]{dudley2014}, Proposition \ref{lem-ref}, Proposition \ref{casSd} and the idea to use the Slepian comparison theorem (see
Step 1 in Proposition \ref{ieii} and the proof in Section~\ref{preuv-solM}) are due to them.

\section{Statements of the main results in the Gaussian case}\label{result}

We now study a necessary and sufficient condition ensuring that the Gaussian random series $\sum f_n^{G,\omega}$ almost surely converges in $\CC^0(\mathcal{M})$.
It turns out that the condition $\sum_{n\geq 1} \|f_n\|_{L^2(\mathcal{M})}^2<+\infty$ implies that the Gaussian random series $f^{G,\omega}:=\sum_{n\geq 1} f_n^{G,\omega}$ defines an element of $L^2(\Omega)$, hence the Dudley pseudo-distance $\delta$ in \eqref{delt} is well defined on $\mathcal{M}$ (see Proposition \ref{con-nec}).

\begin{theo}\label{sol}
There is a constant $K_0>0$ depending only on the Riemannian manifold $\mathcal{M}$ such that, for any $K\geq K_0$, if one considers
\begin{itemize}
\item a sequence of functions $(f_n)_{n\geq 1}$ satisfying $f_n\in E_{(Kn-K,Kn]}$ for each $n\in \N^\star$ (see \eqref{defiEn}) and $\sum_{n\geq 1}\|f_n\|_{L^2(\mathcal{M})}^2<+\infty$,
\item a non-zero dimensional compact submanifold $\mathcal{M}_s\subset \mathcal{M}$ with smooth (eventually empty) boundary,
\end{itemize}
then the following assertions are equivalent:
\begin{enumeratei}
\item\label{soli} The Gaussian random series $\sum f_n^{G,\omega}$ is almost surely convergent in $\mathcal{C}^0(\mathcal{M}_s)$ (where $f_n^{G,\omega}$ is defined in \eqref{ran-fouG2}).
\item\label{solii} The Salem-Zygmund condition is fulfilled
\begin{equation}\label{pz-con3}
\sum_{p=1}^{+\infty} \frac{1}{p\sqrt{\ln(p+1)}}\biggl(\sum_{n=p}^{+\infty} \left\Vert f_n\right\Vert_{L^2(\mathcal{M})}^2 \biggr)^{\sfrac{1}{2}}<+\infty.
\end{equation}
\item\label{soliii} Denoting by $\ep\mto N_{\delta}(\mathcal{M}_s,\ep)$ the covering number\footnote{We recall that $N_{\delta}(\mathcal{M}_s,\ep)$ is the minimal number of open $\delta$-balls of radius $\ep$ that cover $\mathcal{M}_s$.} function of $\mathcal{M}_s$ with respect to the Dudley pseudo-distance
$\delta$ defined in \eqref{delt}, then the entropy condition is satisfied%
\begin{equation}\label{sol-ent}
\int_{0}^{+\infty} \sqrt{\ln(N_{\delta}(\mathcal{M}_s,\ep))} d\ep <+\infty.
\end{equation}
\end{enumeratei}
\item[] In the specific case $\mathcal{M}=\mathbb{S}^d$, with $d\geq 2$, the same conclusion holds if each $f_n$ belongs to the eigenspace $\ker(\Deltag+n(n+d-1))$ provided that we consider \eqref{ran-fouG} instead of \eqref{ran-fouG2}.
\end{theo}
As for Theorem \ref{sol-u}, the almost sure convergence in $\mathcal{C}^0(\mathcal{M})$ is also equivalent to the almost sure convergence in $\mathcal{C}^0(\mathcal{M}_s)$ for any submanifold $\mathcal{M}_s$.

Let us discuss a quantitative version of the last result.
To avoid any problem of measurability due to the uncountability of the submanifold $\mathcal{M}_s$ of $\mathcal{M}$, one usually sets%
\begin{equation}\label{exp-conv}
\E_{\omega}\biggl[\sup_{x\in\mathcal{M}_s}\Big|\sum_{n\geq 1} f_n^{G,\omega}(x)\Big|\biggr]:=\sup_{\substack{\mathcal{F}\subset \mathcal{M}_s\\ \mathcal{F} \text{ countable}}} \E_{\omega}\biggl[\sup_{x\in \mathcal{F}}\Big|\sum_{n\geq 1}f_n^{G,\omega}(x)\Big|\biggr].
\end{equation}
In the concrete situation where the four conditions of Theorem \ref{sol} are true, the random function $x\mto \sum_{n\geq 1}f_n^{G,\omega}(x)$ is almost surely continuous and \eqref{exp-conv} clearly recovers
the classical meaning of $\E_{\omega}\bigl[\Vert \sum_{n\geq 1}f_n^{G,\omega}\Vert_{\CC^0(\mathcal{M}_s)} \bigr]$ by choosing a dense countable subset~$\mathcal{F}$ of $\mathcal{M}_s$. We can now state the quantitative version of Theorem \ref{sol}.

\begin{theo}\label{sol-q}
Under the assumptions of Theorem \ref{sol}, the expectation \begin{equation*}
\E_{\omega}\biggl[\sup_{x\in\mathcal{M}_s}\Big| \sum_{n\geq 1} f_n^{G,\omega}(x)\Big|\biggr]
\end{equation*} and the two numbers appearing in \eqref{pz-con3} and \eqref{sol-ent} are equivalent up to a multiplicative loss merely depending on the Riemannian manifold $\mathcal{M}$, on the submanifold $\mathcal{M}_s$ and on $K$.

Finally, the conclusion also holds true for the particular case $\mathcal{M}=\bbS^d$ in a similar fashion to the last statement of Theorem \ref{sol} (and thus $K$ is irrelevant in the numerical equivalences).
\end{theo}
In the specific case $\mathcal{M}=\mathcal{M}_s=\bbS^d$ with $d\geq 2$, the equivalence of the entropy condition \eqref{sol-ent} and the expectation \eqref{exp-conv} can be deduced by considering $\bbS^d$ as a homo\-geneous space based on the orthogonal group $O_d(\R)$. Actually by fixing one point $x_0\in \bbS^d$ then we have
\begin{equation*}
\E_{\omega}\biggl[\sup_{x\in S^d}\Big| \underbrace{\sum_{n\geq 1} f_n^{G,\omega}(x)}_{=f^{G,\omega}(x)}\Big|\biggr]=
\E_{\omega}\biggl[\sup_{t\in O_d(\R)}\Big| \underbrace{\sum_{n\geq 1} f_n^{G,\omega}(t\cdot x_0)}_{=F^{G,\omega}(t)}\Big|\biggr]
\end{equation*}
where $F^{G,\omega}$ denotes a Gaussian process on $O_d(\R)$ which is stationary (we refer to Proposition \ref{casSd} for computations showing that $\E[F^{G,\omega}(s)F^{G,\omega}(t)]$ merely depends on $\langle s (x_0), t (x_0)\rangle$ and so merely on $s^{-1}t$).
Then the Dudley-Fernique theorem\footnote{Actually we need a non-Abelian Dudley-Fernique theorem that has the same proof as the Abelian case, as noticed in \cite[p.\,96]{pisier1981}.} ensures that the last expectations are equivalent to the entropy condition of $F^{G,\omega}$ on $O_d(\R)$. The Dudley pseudo-distances of $(f^{G,\omega}(x))_{x\in \bbS^d}$ and $(F^{G,\omega}(t))_{t\in O_d(\R)}$, respectively denoted by $\delta$ and $\delta'$, are linked via the formula
\begin{equation*}
\forall (s,t)\in O_d(\R)\times O_d(\R), \quad \delta(s x_0, t x_0) = \delta'(s,t).
\end{equation*}
As a consequence, for any $\ep>0$, we easily see the equality $N_\delta(\bbS^d,\ep)=N_{\delta'}(O_d(\R),\ep)$ of the covering numbers. Hence, the entropy conditions of the Gaussian processes
$(f^{G,\omega}(x))_{x\in \bbS^d}$ and $(F^{G,\omega}(t))_{t\in O_d(\R)}$ are equivalent.

In order to complete the previous remark, we stress that the main novelty of the paper is that the Dudley pseudo-distance of the a priori non-stationary Gaussian process $(f^{G,\omega}(x))_{x\in \mathcal{M}_s}$ with $\mathcal{M}_s\subset \bbS^d$ or $\mathcal{M}_s\subset \mathcal{M}$, is in some sense equivalent to a much simpler pseudo-distance of the form $\Upsilon(\delta_g)$ (see Theorem \ref{pro-s+} and Theorem \ref{pseudo-dis}) with $\Upsilon:[0,+\infty)\to [0,+\infty)$ non-decreasing near $0$ and $\Upsilon$ being quite explicit with respect to the coefficients of the random series (see \eqref{ddu}).
\section{Application to universality results for random series}\label{univer}

In the statement of Theorem \ref{sol-u}, the Salem-Zygmund condition \eqref{pz-con30} does not involve the random variables $X_n$ provided that the mutual independence and the moment assumption \eqref{hyp-Xn} are assumed. Consequently, the almost sure convergence of the random series $\sum X_n(\omega)f_n^\omega(x)$ in $\CC^0(\mathcal{M})$ is universal with respect to the random variables $(X_n)$. That is an analogue of a result by Marcus and Pisier dealing with the random Fourier series $(\omega,x)\in \Omega\times \R/ 2\pi \Z\mto \sum X_n(\omega) e^{inx}$ (see \cite[p.\,7--9]{pisier1981} and \cite[p.\,28, Cor.\,7.3]{pisier1977}).

For the same reason, Theorem \ref{sol-u} also shows the universality with respect to reasonable submanifolds $\mathcal{M}_s$ of $\mathcal{M}$
although the eigenfunctions are considered with respect to the Laplace-Beltrami operator $\Deltag$ of the whole Riemannian manifold $\mathcal{M}$. Indeed with the notations of Theorem \ref{sol-u}, we directly see the equivalence of the following two assertions:
\begin{itemize}
\item the random series $\sum X_n(\omega) f_n^{\omega}$ is almost surely convergent in $\mathcal{C}^0(\mathcal{M}_s)$,
\item the random series $\sum X_n(\omega) f_n^{\omega}$ is almost surely convergent in $\mathcal{C}^0(\mathcal{M})$.
\end{itemize}
For instance, $\mathcal{M}_s$ may be a closed geodesic of $\mathcal{M}$ despite the fact that $\mathcal{M}_s$ is negligible for the Riemannian volume.
The previous equivalence contrasts with the classical theory on a compact group.
Maybe, the nearest result would be the following one: for any compact Abelian group $G$ and any compact subset $T\subset G$ of \textit{non-empty interior}, then a classical random series almost surely converges in $\CC^0(G)$ if and only if it almost surely converges in $\CC^0(T)$ (see the argument following \cite[Th.\,13.3]{ledoux}).

We now want to study the universality with respect to the Riemannian metric of $\mathcal{M}$.
We first recall that the Pisier space on $\mathbb{T}$ is the space of functions $\sum_{n\in \Z}c_n e^{inx}\in L_x^2(\mathbb{T})$ such that the associated random Fourier series \eqref{ran-fou} almost surely converges in $\CC^0(\mathbb{T})$
(see \cite[p.\,2]{pisier1977} or \cite[Ch.\,6]{queff2}). In our context, it is clear how to transfer the notion of Pisier space to a boundaryless Riemannian compact manifold $\mathcal{M}$.
\begin{defi}
Given a function $f\in L^2(\mathcal{M})$ and a parameter $K\geq K_0$, one may decompose
\begin{equation*}
f=\biggl(\int_{\mathcal{M}}f(x)d\mu(x)\biggr)\pun_{\mathcal{M}}+\sum_{n\geq 1} f_n, \quad \text{with } f_n:=\Pi_{(Kn-K,Kn]}(f),
\end{equation*}
where $\mu$ is the Riemannian probability measure of $\mathcal{M}$ and $\Pi_{(Kn-K,Kn]}:L^2(\mathcal{M})\to L^2(\mathcal{M})$ is the spectral projector on the spectral window $(Kn-K,Kn]$ with respect to $\sqrt{-\Deltag}$. Then the ``Pisier space on $\mathcal{M}$'' is the space of functions $f\in L^2(\mathcal{M})$ such that the random series $\sum_{n\geq 1} f_n^{\omega}$ almost surely converges in $\CC^0(\mathcal{M})$.
\end{defi}
For $K\gg 1$, the Salem-Zygmund condition \eqref{pz-con30} gives an explicit semi-norm on $L^2(\mathcal{M})$ that characterizes the functions $f\in L^2(\mathcal{M})$ of the previous definition. But such a semi-norm clearly involves the spectral decomposition of the Laplace-Beltrami operator $\Deltag$ which itself is defined via the Riemannian metric of $\mathcal{M}$.
It turns out that the following universality result holds.

\begin{theo}\label{indep-met}
With the above notations, the ``Pisier space on $\mathcal{M}$'' is independent of the spectral parameter $K$ (provided that $K\gg 1$) and of the Riemannian metric of $\mathcal{M}$.
\end{theo}
Remembering that the space $L^2(\mathcal{M})$
does not depend on the Riemannian structure of $\mathcal{M}$, one may compare Theorem \ref{indep-met} to the following $L^p$ Paley-Zygmund theorem: for any $p\in [1,+\infty)$, the random series $\sum f_n^\omega$ almost surely converges in $L^p(\mathcal{M})$ if and only if $f$ belongs to $L^2(\mathcal{M})$ (see \cite[Th.\,2.3]{imek-aif} but such a result essentially appears in \cite{burq-lebeau} in a different form).

The independence with respect to $K$ in Theorem \ref{indep-met} will be a simple consequence of the Salem-Zygmund condition \eqref{pz-con30} whereas the independence with respect to the Riemannian metric is more involved and uses
\begin{itemize}
\item a semi-classical reformulation of the Salem-Zygmund condition as follows (for a suitable $\Psi\in \CC_c^\infty(\R)$):
\begin{equation*}
\int_{0}^{1} \frac{\|f-\Psi(-h^2 \Deltag)f\|_{L^2(\mathcal{M})}}{h\sqrt{-\ln(h)}}\, dh <+\infty
\end{equation*}
(note that such a reformulation strikingly looks like the condition giving the continuity of stationary Gaussian processes \cite{marcus-shepp,marcus1973} and Proposition \ref{lem-ref}),
\item a few estimates on semi-classical pseudo-differential operators (see Lemma \ref{lem-sc}),
\item an expansion of the operator $\Psi(-h^2 \Deltag)$ with respect to $h$ (see the proof of Proposition \ref{wrt2}), usually referred as a functional semi-classical calculus, as done by Burq-Gérard-Tzvetkov in \cite{burq2004}.
\end{itemize}

\section{Metric considerations}\label{met-con}

In the present section, let us consider a general compact metric space $(\mathcal{M},\delta_g)$. Although we shall settle general facts, we must have in mind that the next results will be applied in the rest of the paper for $\mathcal{M}$ being a Riemannian compact manifold (or submanifold) with its Riemannian distance $\delta_g$.

\begin{prop}\label{lem-ref}
Let $\mu$ be a Borel probability measure on $(\mathcal{M},\delta_g)$ and let $D_g$ be the diameter of $(\mathcal{M},\delta_g)$. We moreover assume that there are positive constants $H_1,H_2,\sigma$ such that for any $t\in [0,D_g]$ and $x\in \mathcal{M}$ we have\vspace*{-3pt}
\begin{equation}\label{ln-hom}
H_1 \left(\frac{t}{D_g}\right)^{\sigma} \leq \mu (B_{\delta_g}(x,t)) \leq H_2 \left(\frac{t}{D_g}\right)^{\sigma}.
\end{equation}
Let $\Upsilon: [0,+\infty)\to [0,+\infty)$ be a subadditive, non-decreasing and right-continuous function satisfying $\Upsilon(0)=0$ and a pseudo-distance $\delta$ on $\mathcal{M}$ which is equivalent to the pseudo-distance $\Upsilon(\delta_g)$ as follows:\vspace*{-3pt}
\begin{equation}\label{equi}
\exists \varrho\geq 1, \quad \frac{\delta}{\sqrt{\varrho}}\leq \Upsilon(\delta_g)\leq \sqrt{\varrho}\, \delta.
\end{equation}
Denoting by $N_\delta(\mathcal{M},\ep)$ the minimal number of $\delta$-open balls of radius $\ep$ that cover $\mathcal{M}$, then the entropy condition for $\delta$\vspace*{-3pt}
\begin{equation}\label{lem-ent}
\int_{0}^{+\infty} \sqrt{\ln(N_\delta(\mathcal{M},\ep))}d\ep<+\infty
\end{equation}
is equivalent to\vspace*{-3pt}
\begin{equation*}
\int_{0}^{1} \frac{\Upsilon(t)}{t \sqrt{\ln(1/t)}}dt <+\infty.
\end{equation*}
\end{prop}

\skpt
\begin{proof}\vspace*{-.5\baselineskip}
\subsubsection*{Topological remark}
One may assume that $\Upsilon$ dos not identically vanish, otherwise the conclusion is obvious and of no interest.
We also note that the subadditivity assumption ensures that $\Upsilon(t)>0$ for any $t>0$ and then $\Upsilon(\delta_g)$ turns out to be a distance. We also easily infer that $\delta$ and $\delta_g$ give rise to the same topology on $\mathcal{M}$.

\subsubsection*{Step 1} We now introduce a pseudo-inverse function $\Upsilon^{-1}:[0,+\infty)\to [0,+\infty]$ as follows
\begin{equation*}
\forall \ep\geq 0, \quad \Upsilon^{-1}(\ep):=\inf\{t\geq 0,\, \Upsilon(t)\geq \ep\},
\end{equation*}
with the usual convention $\inf \varnothing=+\infty$.
The right-continuity of $\Upsilon$ implies the following two equivalences
\begin{align*}
\Upsilon(t)\geq \ep &\iff t \geq \Upsilon^{-1}(\ep), \\
\Upsilon(t)<\ep &\iff t <\Upsilon^{-1}(\ep).
\end{align*}
Hence, we are able to come back to the open balls for the original distance $\delta_g$. For any $x\in \mathcal{M}$ and any $\varepsilon>0$ we have
\begin{align*}
B_{\Upsilon(\delta_g)}(x,\ep)&=
B_{\delta_g}(x,\Upsilon^{-1}(\ep)),\\
\mu(B_{\Upsilon(\delta_g)}(x,\ep))&=
\mu(B_{\delta_g}(x,\Upsilon^{-1}(\ep))).
\end{align*}
Let us set $H=\sfrac{H_2}{H_1}$. The last equality, \eqref{ln-hom} and \eqref{equi} easily imply that $\mu$ is also almost-homogeneous for $\delta$ in the following sense (in which we note the multiplicative loss $\varrho$ in the right-hand side):
\begin{equation}\label{homogg}
\forall (x,y)\in \mathcal{M}^2, \ \forall \ep>0,\quad \mu(B_{\delta}(x,\ep))\leq H \mu(B_{\delta}(y,\varrho \ep)).
\end{equation}
\subsubsection*{Step 2} The following consequence of Step 1 is well-known: for any $\ep>0$ and any $x\in\mathcal{M}$ we have $\mu(B_{\delta}(x,\ep))>0$ and
\begin{equation}\label{min-maj-nb}
\frac{1}{H\mu(B_\delta(x,\varrho\ep))} \leq N_\delta(\mathcal{M},\ep) \leq \frac{H}{\mu(B_\delta(x,\sfrac{\ep}{2\varrho}))}.
\end{equation}
For the bound from below, we consider $x_1,\dots,x_N$ with $N=N_\delta(\mathcal{M},\ep)$ so that the open ball $B_{\delta}(x,\ep)$ cover $\mathcal{M}$. By the the subadditivity of the probability measure $\mu$ (from the previous topological remark, $B_\delta(x,\ep)$ is also open for the original topology on $\mathcal{M}$), \eqref{homogg} and the equality
\begin{equation*}
\mathcal{M}:=B_{\delta}(x_1,\ep)\cup\dots \cup B_{\delta}(x_N,\ep)
\end{equation*}
we get the bound from below of \eqref{min-maj-nb}.
For the bound from above in \eqref{min-maj-nb}, let $N^\star$ be the so-called packing number defined as the largest integer such that $(\mathcal{M},\delta)$ contains $N^\star$ disjoint open balls $B_{\delta}(x_1,\sfrac{\ep}{2}),\dots,B_{\delta}(x_{N^\star},\sfrac{\ep}{2})$. It is well known that \hbox{$N_\delta(\mathcal{M},\ep)\leq N^\star$}. We then easily obtain \eqref{min-maj-nb} thanks to \eqref{homogg} and the inequalities
\begin{equation*}
\frac{1}{H}\, N_{\delta}(\mathcal{M},\ep) \mu(B_{\delta}(x,\sfrac{\ep}{2\varrho}))\leq \sum_{k=1}^{N^\star} \mu(B_{\delta}(x_k,\sfrac{\ep}{2})) \leq 1.
\end{equation*}

\subsubsection*{Step 3} Since $D_g$ is diameter of $(\mathcal{M},\delta_g)$, we have
$\delta(x)\leq \sqrt{\varrho}\,\Upsilon(D_g)$ for any $x\in \mathcal{M}$, and so for any $\ep>\sqrt{\varrho}\,\Upsilon(D_g)$ we have $N_\delta(\mathcal{M},\ep)=1$. That fact leads to the well-known remark stating that the entropy condition \eqref{lem-ent} is merely relevant for $\ep\to 0$.

In other words, for a fixed element $x\in \mathcal{M}$, the inequalities \eqref{min-maj-nb} show that the entropy condition \eqref{lem-ent} is equivalent to the finiteness of
\begin{equation*}
\int_{0}^{+\infty} \sqrt{\ln\Big(\frac{1}{\mu(B_{\delta_g}(x,\Upsilon^{-1}(\ep)))} \Big)} \,d\ep=\int_{0}^{+\infty} \sqrt{\ln\Big(\frac{1}{\mu(B_\delta(x,\ep))} \Big)} \,d\ep.
\end{equation*}
Note that the two last integrals are indeed supported in $[0,\Upsilon(D_g)]$ and their finiteness is also relevant at $\ep\to 0$. By using \eqref{ln-hom}, their finiteness is then equivalent to the finiteness of the following integrals (the second one is in the sense of Stieltjes):
\begin{equation}\label{alm2-ccc2}
\int_{0}^{\Upsilon(D_g)} \sqrt{\ln(\sfrac{D_g}{\Upsilon^{-1}(\ep)})} \,d\ep =\int_{(0,D_g]}
\sqrt{\ln(\sfrac{D_{g}}{t})}\,d\Upsilon(t).
\end{equation}
For any $a\in (0,D_g)$, an integration by parts in the sense of Stieltjes and the condition $\Upsilon(0)=0$ give\vspace*{-3pt}
\[
\frac{1}{2}\int_{a}^{D_g}\frac{\Upsilon(t)}{\sqrt{\ln(\sfrac{D_g}{t})}}\,\frac{dt}{t}=\sqrt{\ln(\sfrac{D_{g}}{a})}\underbrace{\Upsilon(a)}_{=\int_{(0,a]} d\Upsilon(t)}+\int_{(a,D_g]}
\sqrt{\ln(\sfrac{D_{g}}{t})}\,d\Upsilon(t).
\]
By using the last computations and the inequality $\sqrt{\ln(\sfrac{D_g}{a})}\leq \sqrt{\ln(\sfrac{D_g}{t})}$ for any $t\in (0,a]$, we obtain\vspace*{-3pt}
\[
\int_{(a,D_g]}\sqrt{\ln(\sfrac{D_{g}}{t})}\,d\Upsilon(t)\leq \frac{1}{2}\int_{a}^{D_g}\frac{\Upsilon(t)}{\sqrt{\ln(\sfrac{D_g}{t})}}\,\frac{dt}{t} \leq \int_{(0,D_g]}\sqrt{\ln(\sfrac{D_{g}}{t})}\,d\Upsilon(t).
\]
By making $a$ tend to $0^+$ and looking at \eqref{alm2-ccc2}, we get
\begin{equation*}
\int_{0}^{\Upsilon(D_{g})} \sqrt{\ln(\sfrac{D_{g}}{\Upsilon^{-1}(\ep)})}\,d\ep=\frac{1}{2}\int_{0}^1 \frac{\Upsilon(D_g t)}{t\sqrt{\ln(1/t)}}\,dt.
\end{equation*}

\subsubsection*{Step 4} We now explain how to get rid of the parameter $D_g$ in the last integral.
For any $s>0$, we denote by $\lceil s\rceil$ the least integer greater than or equal to a given number $s>0$.
By using the fact that $\Upsilon$ is subadditive and non-decreasing, one remarks the inequality $\Upsilon(st)\leq \lceil s \rceil \Upsilon(t)$ for any $t\geq 0$. By replacing $(s,t)$ with $(\sfrac{1}{s},st)$, the previous bound can be reversed as follows\vspace*{-1pt}
\begin{equation}\label{subb}
\frac{1}{\lceil 1/s\rceil} \Upsilon(t)\leq \Upsilon(st)\leq \lceil s \rceil \Upsilon(t).
\end{equation}
In other words, one may replace $\Upsilon(D_{g} t)$ with $\Upsilon(t)$ provided we authorize a multiplicative loss merely depending on $D_{g}$. The proof is finished.
\end{proof}

Note that the above proof indeed shows the quantitative following version:
\begin{coro}\label{lem-quant}
Under the assumptions of Proposition \ref{lem-ref}, we have the equivalence\vspace*{-3pt}
\begin{equation*}
\int_{0}^{+\infty} \sqrt{\ln(N_\delta(\mathcal{M},\ep))}\,d\ep
\simeq\int_{0}^1 \frac{\Upsilon(t)}{t\sqrt{\ln(1/t)}} \,dt
\end{equation*}
up to multiplicative constants which are independent of $\Upsilon$ (but possibly dependent on $D_g$, $\sfrac{H_2}{H_1}$ and $\varrho$).
\end{coro}

\begin{proof}
By looking at the proof of Proposition \ref{lem-ref}, we have to check in Step 3 that $\Upsilon(D_g)$ is controlled by any of the two integrals with a constant which is independent of $\Upsilon$.
For the entropy integral, let us denote by $D$ the diameter of $(\mathcal{M},\delta)$. We have $\Upsilon(D_g)\leq \sqrt{\varrho}\, D$ thanks to \eqref{equi}. Since no ball of radius $\sfrac{D}{3}$ can cover $(\mathcal{M},\delta)$, we get
\begin{equation*}
\int_{0}^{+\infty} \sqrt{\ln(N(\mathcal{M},\varepsilon))} d\varepsilon \geq
\int_{0}^{\sfrac{D}{3}} \sqrt{\ln(N(\mathcal{M},\varepsilon))}\, d\varepsilon \geq
\sqrt{\ln(2)}\, \frac{D}{3} \geq \frac{\sqrt{\ln(2)}}{3\varrho}\, \Upsilon(D_g).
\end{equation*}
For the integral $\int_{0}^1 \Psfrac{\Upsilon(t)}{t\sqrt{\ln(1/t)}}\, dt$, we use \eqref{subb} as follows:
\begin{equation*}
\int_{\sfrac{1}{2}}^{1} \frac{\Upsilon(t)}{t\sqrt{\ln(1/t)}} \,dt \geq 2\sqrt{\ln(2)}\,\Upsilon(\sfrac{1}{2})\geq \frac{2\sqrt{\ln(2)}}{\lceil 2 D_g \rceil}\, \Upsilon(D_g).\qedhere
\end{equation*}
\end{proof}

\section{Partial pseudo-distances on spheres}\label{partial1}

For any integers $d\geq 2$ and $n\geq 1$, we denote by $\delta_n$ the pseudo-distance on the sphere $\bbS^d$ given by\vspace*{-3pt}
\begin{equation}\label{defipseu}
\delta_n(x,y)^2:= \frac{1}{\dim(E_n)}\sum_{\lambda_k^2=n(n+d-1)} |\phi_{k}(x)-\phi_{k}(y)|^2,
\end{equation}
for any $(x,y)\in \bbS^d\times \bbS^d$ and where $E_n$ stands for the eigenspace $\ker(\Deltag+n(n+d-1))$.
It is well-known that $E_n$ is the space of the restrictions to $\bbS^{d}$ of the $n$-homogeneous harmonic polynomials of $\R^{d+1}$ (see for instance \cite[Prop.\,9.3.5]{faraut}). In particular, for $n$ even, the relation $\phi_k(x)=\phi_k(-x)$ holds for any $\phi_k\in E_n$ and $x\in \bbS^{d}$. And thus we get $\delta_n(x,-x)=0$ for $n$ even.
In other words, if one wants to compare $\delta_n$ to the Riemannian distance $\delta_{g}$ of $\bbS^d$, one must unavoidably restrict $\delta_n$ to a reasonable subset of $\bbS^d$ that does not contain antipodal points. These considerations lead us to the following statement (in which $\overline{B_{\delta_g}}(p,\vartheta)$ is a Riemannian closed ball of $\bbS^d$).

\begin{theo}\label{pro-s+}
For any $d\geq 2$ and for any $\vartheta\in (0,\frac{\pi}{2})$, there is $C=C(d,\vartheta)\geq 1$ such that for any $n\in \N^\star$ and any $p\in \bbS^d$, the pseudo-distance $\delta_n$ is equivalent on the closed ball $\overline{B_{\delta_g}}(p,\vartheta)$ to the distance $\min(1,n\delta_{g})$ as follows\vspace*{-3pt}
\begin{equation*}
\frac{1}{C} \min(1,n\delta_{g})\leq \delta_n \leq C\min(1,n\delta_{g}).
\end{equation*}
\end{theo}
\begin{proof} Let us denote by $e_n(x,y)$ the reproducing kernel of the eigenspace \hbox{$\ker(\Deltag+n(n+d-1))$}:\vspace*{-3pt}
\begin{equation*}
e_n(x,y):=\sum_{\lambda_k^2=n(n+d-1)} \phi_{k}(x)\phi_{k}(y),
\end{equation*}
where each eigenfunction $\phi_{k}$ is assumed to be real-valued. The pseudo-distance $\delta_n$ then takes the form\vspace*{-3pt}
\begin{equation}\label{partn}
\delta_{n}(x,y)^2=\frac{e_n(x,x)+e_n(y,y)-2e_n(x,y)}{\dim(E_n)}.
\end{equation}
Let us recall how $e_n(x,y)$ may be expressed thanks to orthogonal polynomials (for instance Gegenbauer polynomials).
We prefer here Jacobi polynomials (which are directly related to Gegenbauer polynomials).
Let $(P_n^{(\sfrac{d}{2}-1,\sfrac{d}{2}-1)})_{n\in \N}$ be the family of Jacobi polynomials associated to the weight $t\in [-1,1]\mto (1-t^2)^{\sfrac{d}{2}-1}$ and satisfying (see \cite[p.\,58]{szeg}):\vspace*{-3pt}
\begin{equation}\label{jaco1}
P_n^{(\sfrac{d}{2}-1,\sfrac{d}{2}-1)}(1)=\binom{n + \sfrac{d}{2}-1}{n}\eq{n\to +\infty} \frac{n^{\sfrac{d}{2}-1}}{\Gamma(\sfrac{d}{2})} .
\end{equation}
The so-called additional formula (see \cite[Lem.\,2.8 (p.\,143) \& Th.\,2.14 (p.\,149)]{stein71}) ensures the existence of a constant
$c_{d,n}$ satisfying
\begin{equation}\label{KN}
e_n(x,y) =c_{d,n} P_n^{(\sfrac{d}{2}-1,\sfrac{d}{2}-1)} (\langle x,y\rangle).
\end{equation}
Note that $e_n(x,x)=c_{d,n} P_n^{\left(\sfrac{d}{2}-1,\sfrac{d}{2}-1 \right)} (1)$ does not depend on $x$ and so equals
\begin{equation*}
\int_{\bbS^d}e_n(x,x)\,\frac{dx}{\operatorname{vol}_{\bbS^d}(\bbS^d)}=\frac{\dim(E_n)}{\operatorname{vol}_{\bbS^d}(\bbS^d)}.
\end{equation*}
As a consequence, one has
\begin{equation}\label{equivcn}
\frac{c_{d,n}}{\dim(E_n)}=\frac{1}{\operatorname{vol}_{\bbS^d}(\bbS^d) P_n^{(\sfrac{d}{2}-1,\sfrac{d}{2}-1)}(1)}\simeq \frac{1}{n^{\sfrac{d}{2}-1}}.
\end{equation}
We now simplify \eqref{partn} as $\Psfrac{2}{\dim(E_n)}(e_n(x,x)-e_n(x,y))$ and we get the following closed form
\begin{equation}\label{jaco2}
\delta_n(x,y)^2 =\frac{2c_{d,n}}{\dim(E_n)}\Bigl[P_n^{\left(\sfrac{d}{2}-1,\sfrac{d}{2}-1 \right)} (1)-P_n^{\left(\sfrac{d}{2}-1,\sfrac{d}{2}-1 \right)} (\langle x,y\rangle)\Bigr].
	\end{equation}
One may write $\langle x,y\rangle =\cos(\delta_g(x,y))$ where $\delta_g(x,y)\in [0,\pi]$ is the Riemannian distance between $x$ and $y$ on the sphere $\bbS^d$. Hence, we have the equivalence
\begin{equation}\label{ede}
1-\langle x,y\rangle\simeq \delta_{g}(x,y)^2.\end{equation}
Note that the Riemannian distance between two elements $x$ and $y$ on the spherical cap $\overline{B_{\delta_g}}(p,\vartheta)$ is less or equal to $2\vartheta$. Hence we have
\begin{equation*}
\langle x,y\rangle=\cos(\delta_g(x,y))\geq \cos(2\vartheta)>-1.
\end{equation*}
So the conclusion is a consequence of Proposition \ref{jacobi} below and of \eqref{equivcn}, \eqref{jaco2} and~\eqref{ede}.
\end{proof}

\begin{prop}\label{jacobi}
For any integer $d\geq 2$ and any real number $\vartheta \in (0,\sfrac{\pi}{2})$, there is a constant $C_{d,\vartheta}\geq 1$ such that the following inequalities hold true for any $n\geq 1$ and any $\alpha \in [\cos(2\vartheta),1)$
\begin{equation}\label{jacobi2}
\frac{1}{C_{d,\vartheta}} \leq \frac{P_{n}^{(\sfrac{d}{2}-1,\sfrac{d}{2}-1)}(1)-P_{n}^{(\sfrac{d}{2}-1,\sfrac{d}{2}-1)}(\alpha)}{n^{\sfrac{d}{2}-1}\min(1,n^2(1-\alpha))}\leq C_{d,\vartheta}.
\end{equation}
\end{prop}

\skpt
\begin{proof}
\subsubsection*{Step 1}
We first prove a weak version of \eqref{jacobi2} in which the constant $C_{d,\vartheta}$ may depend on $n$.
We set
\begin{equation*}
Q_{n,d}(\alpha):=\frac{P_{n}^{(\sfrac{d}{2}-1,\sfrac{d}{2}-1)}(1)-P_{n}^{(\sfrac{d}{2}-1,\sfrac{d}{2}-1)}(\alpha)}{1-\alpha},
\end{equation*}
which is a polynomial with respect to $\alpha$.
It turns out that $P_n^{(\sfrac{d}{2}-1,\sfrac{d}{2}-1)}$ reaches its maximum in $[\cos(2\vartheta),1]$ at the mere point $1$ (see \cite[p.\,168]{szeg}) and so the polynomial $Q_{n,d}$ is positive on $[\cos(2\vartheta),1)$.
Moreover, it also does not vanish for $\alpha=1$ thanks to the following formula
(see \eqref{jaco1} and \cite[p.\,63, (4.21.7)]{szeg}):
\begin{equation}\label{deri-j}
\bigl(P_n^{(\sfrac{d}{2}-1,\sfrac{d}{2}-1)}\bigr)'(1)=\frac{n+d-1}{2}\,P_{n-1}^{(\sfrac{d}{2},\sfrac{d}{2})}(1)>0.
\end{equation}
By compactness of $[\cos(2\vartheta),1]$, there is a constant $C_{d,\vartheta,n}\geq 1$ and such that the following holds for any $\alpha \in [\cos(2\vartheta),1]$
\begin{equation*}
\frac{1}{C_{d,\vartheta,n}}\leq Q_{n,d}(\alpha)\leq C_{d,\vartheta,n}.
\end{equation*}
Since $1-\alpha$ belongs to $[0,2]$, we have
\begin{equation*}
\frac{\min(1,n^2(1-\alpha))}{n^2}\leq (1-\alpha) \leq 2 \min(1,n^2(1-\alpha))
\end{equation*}
and hence
\begin{equation*}
\frac{1}{n^2 C_{d,\vartheta,n}}\leq \frac{P_{n}^{(\sfrac{d}{2}-1,\sfrac{d}{2}-1)}(1)-P_{n}^{(\sfrac{d}{2}-1,\sfrac{d}{2}-1)}(\alpha)}{\min(1,n^2(1-\alpha))}\leq 2C_{d,\vartheta,n}.
\end{equation*}
Such estimates are very far from the expected inequalities but show that it is sufficient to prove \eqref{jacobi2} for $n\geq n_{d,\vartheta}$ (for a suitable positive integer $n_{d,\vartheta}$ merely depending on $(d,\vartheta)$).

\subsubsection*{Step 2}
The mean value theorem and the formula giving the derivative of a Jacobi polynomial (as in \eqref{deri-j}) ensure the existence of $z\in (\alpha,1)$ such that the previous computations can be continued:
\begin{equation*}
Q_{n,d}(\alpha)=\frac{d}{dz} P_n^{(\sfrac{d}{2}-1,\sfrac{d}{2}-1)} (z)=\frac{n+d-1}{2}\,P_{n-1}^{(\sfrac{d}{2},\sfrac{d}{2})}(z).
\end{equation*}
Classical estimates on Jacobi polynomials (see the proof of \cite[Lem.\,3.9]{imek-crusep} or \eqref{cvu} below) ensure that there is $c_1>0$, merely depending on the dimension $d$, such that the equivalence $P_{n-1}^{(\sfrac{d}{2},\sfrac{d}{2})}(\cos(\theta))\simeq n^{\sfrac{d}{2}}$ uniformly holds true with respect to $\theta\in \left[0,\sfrac{c_1}{n}\right]$.
As a consequence, we get
\begin{equation*}
\forall n\gg 1,\ \forall \theta \in \left(0,\sfrac{c_1}{n}\right], \quad \frac{Q_{d,n}(\cos(\theta))}{n^2}\simeq n^{\sfrac{d}{2}-1}.
\end{equation*}
In the last regime, we remark the equivalence $n^2(1-\cos(\theta))\simeq \min(1,n^2(1-\cos(\theta)))$. We thus have proved
\begin{equation}\label{step3}
\forall n\gg 1,\, \forall \theta \!\in\!(0,\sfrac{c_1}{n}], \quad
\frac{P_{n}^{(\sfrac{d}{2}-1,\sfrac{d}{2}-1)}(1)-P_{n}^{(\sfrac{d}{2}-1,\sfrac{d}{2}-1)}(\cos(\theta))}{\min(1,n^2(1-\cos(\theta)))}\simeq n^{\sfrac{d}{2}-1}.
\end{equation}

\subsubsection*{Step 3} Thanks to \cite[Th.\,7.32.2, p.\,169]{szeg}, we know that for any constant $c_2>0$ one may find $C_2\geq 1$ such that
\begin{equation*}
\forall n\gg 1,\ \forall \theta \in[\sfrac{c_2}{n},\sfrac{\pi}{2}], \quad |P_n^{(\sfrac{d}{2}-1,\sfrac{d}{2}-1)} (\cos(\theta))|\leq \frac{C_2}{\sqrt{n}\, \theta^{\psfrac{d-1}{2}}}.
\end{equation*}
The formula
$P_n^{(\sfrac{d}{2}-1,\sfrac{d}{2}-1)}(-x)=(-1)^n P_n^{(\sfrac{d}{2}-1,\sfrac{d}{2}-1)}(x)$ (see \cite[p.\,59]{szeg}) and the previous inequality implies the following one for any constant $c_3$ larger than $c_2$:
\begin{equation*}
\forall n\gg 1,\ \forall \theta \in[\sfrac{c_3}{n},\pi-\sfrac{c_3}{n}], \quad |P_n^{(\sfrac{d}{2}-1,\sfrac{d}{2}-1)} (\cos(\theta))|\leq C_2\frac{n^{\sfrac{d}{2}-1}}{c_3^{\psfrac{d-1}{2}}}.
\end{equation*}
Remembering \eqref{jaco1} and noting that $C_2$ is independent of $c_3$, one sees that one may choose $c_3$ large enough satisfying for all $\theta\in \left[\sfrac{c_3}{n},\pi-\sfrac{c_3}{n}\right]$:
\begin{align}\nonumber
|P_n^{(\sfrac{d}{2}-1,\sfrac{d}{2}-1)} (\cos(\theta))|& \leq \frac{1}{2}\, P_n^{(\sfrac{d}{2}-1,\sfrac{d}{2}-1)} (1), \\ \label{jaco-cc3}
P_n^{(\sfrac{d}{2}-1,\sfrac{d}{2}-1)}(1)-P_n^{(\sfrac{d}{2}-1,\sfrac{d}{2}-1)}(\cos(\theta)) & \eqsn n^{\sfrac{d}{2}-1}.
\end{align}
Since we assumed $\vartheta<\sfrac{\pi}{2}$, the inequality $2\vartheta \leq \pi-\sfrac{c_3}{n}$ holds true for $n\geq n_{d,\vartheta}$ (with a suitable positive integer $n_{d,\vartheta}$).
In contrast with \eqref{step3}, we have the equivalence $\min(1,n^2(1-\cos(\theta))) \simeq 1$ in such a regime. As a consequence, we have
\begin{multline}\label{step4}
\forall n\geq n_{d,\vartheta},\ \forall \theta \in[\sfrac{c_3}{n},2\vartheta], \\ \frac{P_n^{(\sfrac{d}{2}-1,\sfrac{d}{2}-1)} (1)-P_n^{(\sfrac{d}{2}-1,\sfrac{d}{2}-1)} (\cos(\theta))}{\min(1,n^2(1-\cos(\theta)))}\simeq n^{\sfrac{d}{2}-1}.
\end{multline}

\subsubsection*{Step 4} Comparing \eqref{step3} and \eqref{step4}, one cannot exclude the possibility that $c_3$ may have been chosen too large in Step 3.
To complete our proof, we need to understand the case $c_1<c_3$.
Here we invoke the following uniform limit with respect to $t\in [0,c_3]$ (see \cite[p.\,192, (8.1.1)]{szeg}) that makes a connexion between the Jacobi polynomials $P_n^{(\sfrac{d}{2}-1,\sfrac{d}{2}-1)}$ and the Bessel function $J_{\sfrac{d}{2}-1}$:
\begin{multline}\label{cvu}
\lim_{n\to +\infty} \frac{P_n^{(\sfrac{d}{2}-1,\sfrac{d}{2}-1)}\left(\cos\left(\sfrac{t}{n}\right) \right)}{n^{\sfrac{d}{2}-1}}=2^{\sfrac{d}{2}-1}\, \frac{J_{\sfrac{d}{2}-1}(t)}{t^{\sfrac{d}{2}-1}}\\ \text{in the Banach space } \CC_t^0([0,c_3],\R).
\end{multline}
The Poisson representation of the Bessel function $J_{\sfrac{d}{2}-1}$ (see \cite[p.\,15, (1.71.6)]{szeg}) suggests introducing the following real-valued function
\begin{equation}\label{Poisson} \widehat{\sigma_{d-1}}(t):=\frac{J_{\sfrac{d}{2}-1}(t)}{t^{\sfrac{d}{2}-1}}=\frac{1}{2^{\sfrac{d}{2}-1}\Gamma(\psfrac{d-1}{2})\sqrt{\pi}}\int_{-1}^{1} e^{its} (1-s^2)^{\psfrac{d-3}{2}}ds.\end{equation}
The notation $\widehat{\sigma_{d-1}}$ is justified since that function may be interpreted as the radial part of the Fourier transform of the spherical measure $\sigma_{d-1}$ on $\bbS^{d-1}$ (see \eqref{defi-bes}).
One directly checks that the function $\widehat{\sigma_{d-1}}$ is real-valued, admits a maximum at the unique point $t=0$ and satisfies $\lim_{t\to +\infty}\widehat{\sigma}_{d-1}(t)=0$ (thanks to the Riemann-Lebesgue lemma).
As a consequence, we get the strict inequality $\sup_{t\geq c_1} \widehat{\sigma}_{d-1}(t)<\widehat{\sigma}_{d-1}(0)$ which in turn implies the equivalent one
\begin{equation}\label{bess}
\exists \rho\in (0,1),\ \forall t\geq c_1,\quad \widehat{\sigma}_{d-1}(t) < \rho \widehat{\sigma}_{d-1}(0).
\end{equation}
Coming back to \eqref{cvu}, we see that the sequence of functions
\begin{equation*}
t\mto \frac{1}{n^{\sfrac{d}{2}-1}}\left[P_n^{(\sfrac{d}{2}-1,\sfrac{d}{2}-1)}\left(\cos\left(\sfrac{t}{n}\right) \right)-\rho
P_n^{(\sfrac{d}{2}-1,\sfrac{d}{2}-1)}\left(1 \right)\right]
\end{equation*}
uniformly converges on $[c_1,c_3]$ to the continuous negative function
\begin{equation*}
t\mto 2^{\sfrac{d}{2}-1}[\widehat{\sigma}_{d-1}(t)-\rho \widehat{\sigma}_{d-1}(0)].
\end{equation*}
Hence, we infer that the following holds true
\begin{equation*}
\forall n\gg 1,\ \forall \theta \in[\sfrac{c_1}{n},\sfrac{c_3}{n}],\quad
P_n^{(\sfrac{d}{2}-1,\sfrac{d}{2}-1)}(\cos(\theta)) \leq\rho P_n^{(\sfrac{d}{2}-1,\sfrac{d}{2}-1)}(1).
\end{equation*}
Such new estimates can now be combined to
\[
-P_n^{(\sfrac{d}{2}-1,\sfrac{d}{2}-1)}(\cos(\theta)) \leq C_d P_n^{(\sfrac{d}{2}-1,\sfrac{d}{2}-1)}(1)
\]
(as a consequence of \eqref{cvu}). Hence
\eqref{jaco-cc3} and \eqref{step4} hold with $c_1$ (instead of $c_3$).
Taking account of \eqref{step3}, we finally obtain the wanted uniform equivalence \eqref{jacobi2}.\end{proof}

\section{Partial pseudo-distances on manifolds}\label{partial2}

In this part, $\mathcal{M}$ is a general boundaryless compact Riemannian manifold (and we use the notations of the introduction about the spectral analysis of the Laplace-Beltrami operator). For any interval $\mathcal{I}\subset [0,+\infty)$ we define the spectral function of $\mathcal{I}$ (with respect to $\sqrt{-\Deltag}$) as follows: for any $(x,y)\in \mathcal{M}^2$ we set
\begin{equation*}
e_{\mathcal{I}}(x,y)=\sum_{\lambda_k\in \mathcal{I}} \phi_k(x)\phi_k(y),
\end{equation*}
where we recall that each eigenfunction $\phi_k$ of $\Deltag$ is assumed to be real-valued.
Let us recall an idea we attribute to Burq and Lebeau in \cite[p.\,923]{burq-lebeau} (we also refer to the end of the proof of \cite[Lem.\,8.1]{imek-aif} for details): one may find two numbers $K_0\geq 1$ and $C\geq 1$ that merely depend on the Riemannian manifold $\mathcal{M}$ such that the following on-diagonal estimates hold
\begin{equation}\label{enxx}
\forall K \geq K_0,\ \forall n\geq 1,\ \forall x\in \mathcal{M}, \quad \frac{K^{d} n^{d-1}}{C}\leq e_{(Kn-K,Kn]}(x,x)\leq
CK^{d} n^{d-1}.
\end{equation}
Upon modifying $C$, we note that the same estimates hold for $\dim(E_{(Kn-K,Kn]})$ by~integration on the compact manifold $\mathcal{M}$:
\begin{equation}\label{enxx2}
\forall K \geq K_0,\ \forall n\geq 1,\ \forall x\in \mathcal{M}, \quad \frac{K^{d} n^{d-1}}{C}\leq
\dim(E_{(Kn-K,Kn]})\leq CK^{d} n^{d-1}.
\end{equation}
We now introduce a pseudo-distance $\delta_n$ on the manifold $\mathcal{M}$ which is analogue to \eqref{defipseu} (we keep the same notation $\delta_n$ for simplicity).
For any pair $(x,y)\in \mathcal{M}^2$ we set
\begin{align}\label{defideltan}
\delta_n(x,y)^2&:=\frac{1}{\dim(E_{(Kn-K,Kn]})}\sum_{\lambda_k\in (Kn-K,Kn]}|\phi_{k}(x)-\phi_{k}(y)|^2 \\ \label{defideltan2}
& \hphantom{:}= \frac{e_{(Kn-K,Kn]}(x,x)+e_{(Kn-K,Kn]}(y,y)-2e_{(Kn-K,Kn]}(x,y) }{\dim(E_{(Kn-K,Kn]})}.
\end{align}
By bounding $|\phi_k(x)-\phi_k(y)|^2 \leq 2(|\phi_k(x)|^2+|\phi_k(y)|^2)$ and using \eqref{enxx} and \eqref{enxx2}, we~remark the uniform estimate
\begin{equation}\label{uniff}
\sup_{n\geq 1} \sup_{(x,y)\in \mathcal{M}^2} \delta_n(x,y)<+\infty.
\end{equation}
In the sequel, we will prove the following analogue result of Theorem \ref{pro-s+} (but here the equivalence \eqref{var-pdis} holds in the whole manifold $\mathcal{M}$).

\begin{theo}\label{pseudo-dis}
There are two constants $K_0\geq 1$ and $C\geq 1$ that merely depend on the Riemannian manifold $\mathcal{M}$ such that the following equivalence\footnote{We note that $\delta_n(x,y)$ implicitly depends on $K$.} holds true for any $K\geq K_0$, any pair $(x,y)\in \mathcal{M}^2$ and any $n\in \N^\star$:
\begin{equation}\label{var-pdis}
\frac{1}{C} \min\big(1,Kn\delta_g(x,y) \big)\leq \delta_n(x,y)\leq C \min\big(1,Kn\delta_g(x,y) \big).
\end{equation}
\end{theo}

The main interest of the previous result is that it holds true without any geometric assumption on the boundaryless compact Riemannian manifold $\mathcal{M}$. It indeed relies on the freedom to set $K$ large enough.
Under specific geometric assumptions, a similar upper bound to that of \eqref{var-pdis} appears in \cite[Th.\,1]{canzani2014high}. Moreover, a much more precise than \eqref{var-pdis} is given in \cite[Th.\,9]{canzani2015scaling} by assuming a ``geometric mutually nonfocal hypothesis'' on the manifold $\mathcal{M}$.

We now recall the relation linking the Bessel function $J_{\sfrac{d}{2}-1}$ and the Fourier transform of the spherical measure $\sigma_{d-1}$ on $\bbS^{d-1}$, for any $\nu\in \R^d\setminus\{0\}$ we set
\begin{equation*}
\frac{J_{\psfrac{d-2}{2}}(|\nu|)}{|\nu|^{\psfrac{d-2}{2}} }=\frac{1}{(2\pi)^{\sfrac{d}{2}}}\int_{\bbS^{d-1}}e^{i\langle \nu,w\rangle} d\sigma_{d-1}(w)=\frac{1}{(2\pi)^{\sfrac{d}{2}}}\int_{\bbS^{d-1}}\exp(i|\nu|w_1) d\sigma_{d-1}(w).
\end{equation*}
Since the previous Fourier transform is radial, we make the following abuse of notation:%
\begin{equation}\label{defi-bes}
\forall \nu>0, \quad \widehat{\sigma_{d-1}}(\nu):=\frac{J_{\psfrac{d-2}{2}}(\nu)}{\nu^{\psfrac{d-2}{2}} }=\frac{1}{(2\pi)^{\sfrac{d}{2}}}\int_{\bbS^{d-1}}\exp(i \nu w_1) d\sigma_{d-1}(w).
\end{equation}
We will use some material developed by Canzani and Hanin. More precisely we need a suitable asymptotic of the spectral function given in the next theorem (proved by Hörmander in \cite[Th.\,4.4]{horm68} for the case $|I|=|J|=0$). But the choice of the
explicit principal term has been enlightened by Canzani and Hanin in \cite{canzani2015scaling}.
The asymptotic formula, involving derivatives, is now considered as known (see the introduction of
\cite{canzani2017c}). We give here a reformulation with coordinate patches.

\begin{theo}\label{hohor}
There is $\alpha_0>0$ such that for any coordinate patch $\tau:U\subset \R^d\to V\subset \mathcal{M}$ with $\operatorname{diam}(V)<\alpha_0$, for any multi-indexes $I\in \N^d$ and $J\in \N^d$, the following asymptotic holds true for any $(x,y)\in V^2$ and $\lambda \geq 0$ (where the derivatives are seen in the coordinate patch):
\begin{equation*}\partial_x^{I}\partial_y^J \sum_{\lambda_j\leq \lambda} \phi_j(x)\phi_j(y)=\frac{\partial_x^{I}\partial_y^J}{(2\pi)^{\sfrac{d}{2}}}\int_{0}^{\lambda }\nu^{d-1} \widehat{\sigma_{d-1}}(\nu \delta_g(x,y))
d\nu+\mathcal{O}\bigl((1+\lambda)^{d-1+|I|+|J|} \bigr).
\end{equation*}
Finally, the remainder is uniform provided that $x$ and $y$ run over a compact subset of the open set $V$.
\end{theo}
\begin{proof}
Since we do not know a published reference where Theorem \ref{hohor} is stated with a proof, we explain in Appendix \ref{annex-weyl} how it can be easily recovered as a consequence of the analysis of \cite{canzani2017c}, a Bernstein-type inequality proved in \cite{bin2004} and a Tauberian theorem as in \cite[Ch.\,4]{sogge2017fourier}.
\end{proof}

The reason of restricting the estimates of Theorem \ref{hohor} to compact subsets of the open subset $V$ is due to the definition of the $\CC^\infty(\mathcal{M})$-topology (see \cite[Ch.\,XVII.1 \& XVII.2]{dieud3}). Obviously for $|I|=|J|=0$, coordinate patches are useless and the estimates hold on the whole compact manifold.

If the two points $x$ and $y$ are far from each other, we shall need to forget the principal term of the asymptotic of the spectral function (at least for $|I|=|J|=0$).
In~\cite[line (79)]{canzani2015scaling}, a geometric mutually nonfocal hypothesis on the manifold~$\mathcal{M}$ is assumed, and thus Canzani and Hanin make use of an asymptotic of Safarov \cite{safarov1988} stating $e_{[0,\lambda]}(x,y)=o(\lambda^{d-1})$.
Since our issue is posed without any geometric assumption, we replace their argument with the off-diagonal Hörmander estimate $e_{[0,\lambda]}(x,y)=\mathcal{O}(\lambda^{d-1})$.

\begin{prop}\label{far}
For any fixed $\alpha>0$, one may find two numbers $K_0\geq 1$ and $c\geq 1$ (depending on the Riemannian manifold $\mathcal{M}$ and $\alpha>0$) satisfying the following property: for any $K\geq K_0$ the following implication holds true for any $n\in \N^\star$ and any pair $(x,y)\in \mathcal{M}^2$:
\begin{equation*}
\delta_g(x,y)\geq \alpha \implies \frac{1}{c}\leq \delta_n(x,y)\leq c.
\end{equation*}
\end{prop}
\begin{proof}
The upper bound $\delta_n(x,y)\leq c$ has already been proved in \eqref{uniff}.
From \cite[line (4.11)]{horm68} (see also \cite[p.\,162, Th.\,21.1, point\,2)]{shub-pseudo}), we know that for any $\alpha>0$ there is $C>0$ such that, for any $(x,y)\in \mathcal{M}^2$, the following uniform off-diagonal estimate holds true
\begin{equation*}
\delta_g(x,y)\geq \alpha \implies \forall \lambda \geq 0,\ \left\vert e_{[0,\lambda]}(x,y)\right\vert\leq C(1+\lambda)^{d-1},
\end{equation*}
where $C$ depends on $\alpha$ and on the Riemannian structure of $\mathcal{M}$.
Thus we obtain
\begin{align*}
|e_{(Kn-K,Kn]}(x,y)| & = \left\vert e_{[0,Kn+K]}(x,y)-e_{[0,Kn]}(x,y) \right\vert \\
& \leq C(1+K n+K)^{d-1}+C(1+K n)^{d-1} \\
& \leq C(3^{d-1}+2^{d-1}) (Kn)^{d-1}.
\end{align*}
We now recall the estimate $e_{(Kn-K,Kn]}(x,x)\simeq K^d n^{d-1}$ (see \eqref{enxx}), which also implies $\dim(E_{(Kn-K,Kn]})\simeq K^d n^{d-1}$ by integration over $\mathcal{M}$.
Remembering the formula \eqref{defideltan2}, we infer that there are three positive constants $C_1,C_2,C_3$ (independent of $n$ and $K$) such that
\begin{equation}\label{mimin}
\begin{aligned}
\delta_n(x,y)^2&\geq \frac{n^{d-1}}{\dim(E_{(Kn-K,Kn]})}\left[C_1K^{d}-C_2 K^{d-1}\right] \\
& \geq \frac{C_1K^{d}-C_2 K^{d-1}}{C_3 K^d}.
\end{aligned}
\end{equation}
We conclude by making $K$ tend to $+\infty$.
\end{proof}

We now need to improve the previous result if $n\delta_g(x,y)$ is bounded from below or from above.
By comparison with \cite[p.\,1728--1729]{canzani2015scaling} that involves geometric assumptions on the manifold $\mathcal{M}$
and relies on several tools obtained by Zelditch, Potash and Xu, our proofs of the next two results merely use the asymptotic of the spectral function given by Theorem \ref{hohor}.

\begin{prop}\label{minn}
For any fixed $\beta>0$, one may find two numbers $K_0\geq 1$ and $c\geq 1$ (depending on the Riemannian manifold $\mathcal{M}$ and $\beta>0$) satisfying the following property: for any $K\geq K_0$, for any $n\in \N^\star$ and any pair $(x,y)\in \mathcal{M}^2$ the following implication holds true
\begin{equation*}
\delta_g(x,y)\geq \frac{\beta}{Kn} \implies \frac{1}{c}\leq \delta_n(x,y)\leq c.
\end{equation*}
\end{prop}
\begin{proof}
As for Proposition \ref{far}, the upper bound $\delta_n\leq c$ is given in \eqref{uniff}.
To get the lower bound, we apply Proposition \ref{far} with the constant $\alpha=\alpha_0$ appearing in the statement of Theorem \ref{hohor}. Hence, the case $\delta_g(x,y)\geq \alpha_0$ is already done.
We now assume $\alpha_0>\delta_g(x,y)\geq \sfrac{\beta}{Kn}$ and we have the following asymptotic\vspace*{-3pt}
\begin{equation}\label{dd2}
e_{(Kn-K,Kn]}(x,y)=\int_{K n-K}^{K n}\nu^{d-1}\widehat{\sigma_{d-1}}(\nu \delta_g(x,y))\,\frac{d\nu}{(2\pi)^{\sfrac{d}{2}}}+(Kn)^{d-1}\mathcal{O}(1),
\end{equation}
where the remainder $\mathcal{O}(1)$ is uniformly bounded with respect to $(x,y,K,n)$.
By using the definition \eqref{defideltan2} of the pseudo-distance $\delta_n$, we get the relation\vspace*{-3pt}\enlargethispage{.5\baselineskip}%
\begin{multline*}
\dim(E_{(Kn-K,Kn]}) \delta_n(x,y)^2 \\[-3pt]
= 2K^{d}\int_{n-1}^{n} \nu^{d-1} \left[\widehat{\sigma}_{d-1}(0)-\widehat{\sigma}_{d-1}(K \nu \delta_g(x,y)) \right]\frac{d\nu}{(2\pi)^{d/2}}+(Kn)^{d-1}\mathcal{O}(1).
\end{multline*}
Remembering that the function $|\widehat{\sigma}_{d-1}|$ is bounded by $\widehat{\sigma}_{d-1}(0)$ (see \eqref{defi-bes}), we can restrict our analysis on $\big[n-\sfrac{1}{2},n\big]$ so that $\delta_n(x,y)^2$ is bounded from below by\vspace*{-7pt}
\begin{multline}\label{minonono}
\frac{2K^{d}}{\dim(E_{(Kn-K,Kn]}) }\int_{n-\sfrac{1}{2}}^{n} \nu^{d-1} \left[\widehat{\sigma}_{d-1}(0)-\widehat{\sigma}_{d-1}(K \nu \delta_g(x,y)) \right]\frac{d\nu}{(2\pi)^{d/2}}\\[-3pt]
+\frac{(Kn)^{d-1}\mathcal{O}(1)}{\dim(E_{(Kn-K,Kn]}) }.
\end{multline}
By looking the last integral and our assumptions, one notices\vspace*{-3pt}
\begin{equation*}
\forall \nu\in[n-\sfrac{1}{2},n], \quad K \nu \delta_g(x,y)\geq \left(n-\sfrac{1}{2}\right)\sfrac{\beta}{n}\geq \sfrac{\beta}{2}.
\end{equation*}
The inequality \eqref{bess} (proved in the case $\mathcal{M}=\bbS^d$) may be used to ensure the existence of $\rho\in (0,1)$ (merely depending on the dimension $d$ and on $\beta$) such that the following holds true\vspace*{-3pt}
\begin{equation*}
\forall \nu\in[n-\sfrac{1}{2},n], \quad \widehat{\sigma}_{d-1}(K \nu \delta_g(x,y))\leq \rho \widehat{\sigma}_{d-1}(0).
\end{equation*}
As at the end of the proof of Proposition \ref{far}, we obtain an inequality similar to \eqref{mimin} which in turn gives suitable constants $K_0$ and $c$.
\end{proof}

\begin{prop}\label{mino2}
There are three numbers $\beta_0\in (0,1), K_0\geq 1$ and $c\geq 1$ (depending on the Riemannian manifold $\mathcal{M}$) satisfying the following property: for any $K\geq K_0$, for any $n\in \N^\star$ and any pair $(x,y)\in \mathcal{M}^2$, the following implication holds true\vspace*{-3pt}
\begin{equation*}
\delta_g(x,y)<\frac{\beta_0}{Kn} \implies \frac{K n\delta_g(x,y)}{c}\leq \delta_n(x,y)\leq cKn\delta_g(x,y).
\end{equation*}
\end{prop}
\begin{proof}
The main issue is that an inequality like \eqref{minonono} seems to give an unavoidable $\mathcal{O}(1)$ remainder. Forgetting technicalities, one may consider that the main idea relies on the formula \eqref{doubl} that will need derivatives of order $2$ of the spectral function.
In~the next steps, we will make several restrictions so that any small enough number~$\beta_0$ will be convenient.\vspace*{-3pt}\mbox{}

\subsubsection*{Step 1}
We first recall that if $\beta_0$ is small enough, then any two points $x$ and $y$ of~$\mathcal{M}$, satisfying $\delta_g(x,y)< \beta_0$, belong to a same geodesically convex Riemannian ball.
To see that point, we invoke the Whitehead theorem (see \cite[Th.\,5.14]{cheeger}):
there is a continuous function $c:\mathcal{M}\to (0,+\infty)$, called the convexity radius, such that each Riemannian ball $B_{\delta_g}(z,r)$, for $z\in \mathcal{M}$, is geodesically convex for $r<c(z)$.
We~finish by using the compactness of $\mathcal{M}$ and choosing $\beta_0$ less than the positive number $\min_{z\in \mathcal{M}}c(z)$.\vspace*{-3pt}\mbox{}

\subsubsection*{Step 2} We now claim that if $\beta_0$ is small enough, then for any $z\in \mathcal{M}$ the open Riemannian ball $B_{\delta_g}(z,\beta_0)$ is relatively compact in the domain of a suitable local chart around $z$.
We again give a compactness argument.
Let $r(z)>0$ be such that there is a coordinate patch around $z$ (for some open subset $U_z$ of $\R^d$)
\begin{equation*}
\begin{array}{rccl}
& \R^d & & \mathcal{M} \\
& \cup & & \cup \\
\tau_z:&U_z& \longrightarrow & B_{\delta_g}(z,r(z)).
\end{array}
\end{equation*}
By compactness, one may cover the compact manifold $\mathcal{M}$ with a finite atlas as follows
\begin{equation*}
\mathcal{M}=B_{\delta_g}(z_1,r(z_1))\cup \dots \cup B_{\delta_g}(z_\ell,r(z_\ell)).
\end{equation*}
In other words, any $z\in \mathcal{M}$ satisfies $
\max_{1\leq k\leq \ell}(r(z_k)-\delta_g(z,z_k))>0$. By continuity and compactness, one may enforce the previous inequality as follows for some positive constant $R$:
\begin{equation*}
\max_{1\leq k\leq \ell}\bigl(r(z_k)-\delta_g(z,z_k)\bigr)\geq 2R.
\end{equation*}
Hence, by choosing $\beta_0\leq R$, for any $z\in \mathcal{M}$, there is $k\in \{1,\dots,\ell\}$ such that we have the inclusion
\begin{equation*}
\overline{B_{\delta_g}(z,\beta_0)}\subset \overline{B_{\delta_g}(z_k,r(z_k)-R)}.
\end{equation*}
In other words, $B_{\delta_g}(z,\beta_0)$ is relatively compact in the domain $B_{\delta_g}(z_k,r(z_k))$.

\subsubsection*{Step 3} Now fix $x$ and $y$ satisfying $\delta_g(x,y)< \sfrac{\beta_0}{Kn}$ as in the statement (other restrictions on $\beta_0$ will be given in Step 4 and Step 5).
Since we are looking for $K\geq 1$, we also have $\delta_g(x,y)<\beta_0$.
Thanks to Step 2, we can consider a coordinate patch $\tau_{k}:U_{z_k}\subset \R^d\to B_{\delta_g}(z_k,r(z_k))\subset \mathcal{M}$ so that the ball $B_{\delta_g}(x,\beta_0)$ is included in the compact set
$\overline{B_{\delta_g}(z_k,r(z_k)-R)}$.
Thanks to Step 1, there is a geodesic
\[
\gamma:[0,1]\to \overline{B_{\delta_g}(z_k,r(z_k)-R)}
\]
starting from $x$ and stopping at $y$.
By using the symmetry of $(x,y)\!\mto\! e_{(Kn-K,Kn]}(x,y)$, we obtain the following integral formulas
\begin{equation}\label{doubl}
\begin{aligned}
e_{(Kn-K,Kn]}(x,x)&+e_{(Kn-K,Kn]}(y,y)-2e_{(Kn-K,Kn]}(x,y) \\
& = e_{(Kn-K,Kn]}(y,y)-e_{(Kn-K,Kn]}(y,x)\\
&\hspace*{3.03cm}
-[e_{(Kn-K,Kn]}(x,y)-e_{(Kn-K,Kn]}(x,x)] \\
& = \int_{0}^1 \frac{\partial }{\partial t_1}\{ e_{(Kn-K,Kn]} (\gamma(t_1),y)-e_{(Kn-K,Kn]} (\gamma(t_1),x)\} dt_1 \\ 
& = \int_{0}^1 \int_{0}^1 \frac{\partial^2 }{\partial t_1 \partial t_2}\Bigl\{ e_{(Kn-K,Kn]}(\gamma(t_1),\gamma(t_2))\Bigr\} dt_1 dt_2.
\end{aligned}
\end{equation}

\subsubsection*{Step 4} Assuming that $\beta_0$ is smaller than the constant $\alpha_0$ of Theorem \ref{hohor}, it is natural to compare the previous integral formula by replacing $e_{(Kn-K,Kn]}$ with its principal term appearing in \eqref{dd2}.
More precisely, let us introduce
\begin{equation}\label{dd3}
e_{(Kn-K,Kn]}^{\wp}(x,y)=K^{d} \int_{n-1}^{n}\nu^{d-1}\widehat{\sigma_{d-1}}(K \nu \delta_g(x,y))\,
\frac{d\nu}{(2\pi)^{\sfrac{d}{2}}}.
\end{equation}
We now claim that the following inequality holds true for any $(t_1,t_2)\in [0,1]^2$:
\begin{multline}\label{remem}
\left\vert \frac{\partial^2}{\partial t_1 \partial t_2} \Big\{ e_{(Kn-K,Kn]}(\gamma(t_1),\gamma(t_2))- e_{(Kn-K,Kn]}^{\wp}(\gamma(t_1),\gamma(t_2)) \Big\} \right\vert\\
\leq C(Kn)^{d+1} \delta_g(x,y)^2,
\end{multline}
where $C$ merely depend on the Riemannian structure of $\mathcal{M}$.
To obtain that bound, we work with the coordinate patch $\tau_{k}:U_{z_k}\subset \R^d\to B_{\delta_g}(z_k,r(z_k))\subset \mathcal{M}$ around $x$ and $y$ as above:
\begin{equation*}
e_{(Kn-K,Kn]}(\gamma(t_1),\gamma(t_2)) = \overbrace{e_{(Kn-K,Kn]}\circ \bigl(\tau_k\oplus \tau_k \bigr)}^{U_{z_k}\times U_{z_k}\to \R} \overbrace{\bigl(\tau_k^{-1}\circ \gamma(t_1),\tau_k^{-1} \circ \gamma(t_2)\bigr)}^{\in U_{z_k}\times U_{z_k}}.
\end{equation*}
By denoting $(p,q)$ a generic point of $U_{z_k}\times U_{z_k}$, the previous formula allows us to write the double derivative
\begin{equation*}
\frac{\partial^2 }{\partial t_1 \partial t_2}\Big\{ e_{(Kn-K,Kn]}(\gamma(t_1),\gamma(t_2))\Big\}
\end{equation*} as follows (in which we denote by $(\tau_k^{-1})_i$ the $i$-th coordinate of $\tau_k^{-1}$):
\begin{multline*}
\sum_{i=1}^{d} \sum_{j=1}^{d}
\frac{d}{dt_1} \bigl((\tau_k^{-1})_i \circ \gamma(t_1) \bigr)\times
\frac{d}{dt_2} \bigl((\tau_k^{-1})_j \circ \gamma(t_2) \bigr) \times \Lambda_{ij}\\
\text{with }\Lambda_{ij}=\frac{\partial^2 }{\partial p_i \partial q_j} \{ e_{(Kn-K,Kn]}\circ (\tau_k\oplus \tau_k)\}\Big\vert_{(\tau_k^{-1}\circ \gamma(t_1),\tau_k^{-1} \circ \gamma(t_2))}.
\end{multline*}
A similar computation holds true for $e_{(Kn-K,Kn]}^{\wp}$.
Here is the point where we really need the asymptotic given by Theorem \ref{hohor} of the spectral function with derivatives and holding uniformly on any compact subset of the domain $B_{\delta_g}(z_k,r(z_k))$ of the local chart $\tau_k^{-1}$ given by the coordinate patch $\tau_k$.
More precisely we choose the compact subset $\overline{B_{\delta_g}(z_k,r(z_k)-R)}$ that turns out to contain the geodesic $\gamma$ (see Step 3).
Theorem \ref{hohor} then shows the bound (uniformly in $(t_1,t_2)\in [0,1]^2$):
\begin{equation*}
\frac{\partial^2 }{\partial p_i \partial q_j} \Big\{ (e_{(Kn-K,Kn]}-e_{(Kn-K,Kn]}^{\wp})\circ (\tau_k\oplus \tau_k)\Big\}\Big\vert_{\big(\tau_k^{-1}\circ \gamma(t_1),\tau_k^{-1} \circ \gamma(t_2)\big)}=\mathcal{O}\bigl((Kn)^{d+1}\bigr).
\end{equation*}
And hence we get
\begin{multline*}
\left\vert \frac{\partial^2 }{\partial t_1 \partial t_2}\Big\{ e_{(Kn-K,Kn]}(\gamma(t_1),\gamma(t_2))-e_{(Kn-K,Kn]}^{\wp}(\gamma(t_1),\gamma(t_2))\Big\} \right\vert \\
\leq C(Kn)^{d+1} \sum_{i=1}^{d} \sum_{j=1}^d \Big\vert \frac{d}{dt_1} \Big((\tau_k^{-1})_i \circ \gamma(t_1)\Big) \Big\vert\times \Big\vert
\frac{d}{dt_2}\Big((\tau_k^{-1})_j \circ \gamma(t_2)\Big) \Big\vert.
\end{multline*}
where $C$ merely depends on the Riemannian structure of $\mathcal{M}$ (we recall that Step 2 allows us to work in a finite atlas and with a finite collection of compact subsets $\overline{B_{\delta_g}(z_k,r(z_k)-R)}$).
Remembering now that $\gamma:[0,1]\to \mathcal{M}$ is a geodesic between $x$ and $y$, the speed of $\gamma$ is constant and must equal $\delta_g(x,y)$ (since we have parametrized the curve with $[0,1]$). In other words, we have
$\left\Vert \gamma'(t_1)\right\Vert_{T_{\gamma(t_1)} \mathcal{M}}=\delta_g(x,y)$ in the tangent space $T_{\gamma(t_1)} \mathcal{M}$ endowed with its Riemannian inner product.
Still using that we are working with a finite atlas, we clearly obtain, for a suitable uniform constant $C>0$, the following bound\vspace*{-3pt}
\begin{equation*}
\forall t_1 \in [0,1],\quad \Bigl\vert \frac{d}{dt_1} \bigl((\tau_k^{-1})_i\circ \gamma(t_1)\bigr) \Bigr\vert\leq C\left\Vert \gamma'(t_1)\right\Vert_{T_{\gamma(t_1)} \mathcal{M}}=C\delta_g(x,y).
\end{equation*}
A similar reasoning may be done for $t_2$ and \eqref{remem} is proved.

\subsubsection*{Step 5}
We now claim that if $\beta_0$ is small enough, then the contribution of the principal term $e_{(Kn-K,Kn]}^\wp$ defined in \eqref{dd3} is given by the following equivalence (uniformly in $(t_1,t_2)\in [0,1]^2$):\vspace*{-3pt}
\begin{equation}\label{remem2}
\frac{\partial^2}{\partial t_1 \partial t_2}\, e_{(Kn-K,Kn]}^\wp(\gamma(t_1),\gamma(t_2)) \simeq K(Kn)^{d+1} \delta_g(x,y)^2 , \end{equation}
where the constants of equivalence merely depend on the dimension $d=\dim(\mathcal{M})$.
To prove such estimates, we shall use the equality $\delta_g(\gamma(t_1),\gamma(t_2))=\delta_g(x,y)|t_1-t_2|$ coming from the fact that geodesic $\gamma:[0,1]\to \mathcal{M}$ has a speed equaling $\delta_g(x,y)$.
Since the function $\widehat{\sigma_{d-1}}:(0,+\infty)\to \R$ admits an even smooth extension to $\R$ (see \eqref{Poisson} and~\eqref{defi-bes}), one may write
\begin{align*}
\widehat{\sigma_{d-1}}\bigl(K\nu\delta_g(\gamma(t_1),\gamma(t_2))\bigr)&=\widehat{\sigma_{d-1}}\bigl(K\nu \delta_g(x,y) |t_1-t_2|\bigr)\\
&= \widehat{\sigma_{d-1}}\bigl(K\nu \delta_g(x,y) (t_1-t_2)\bigr)
\end{align*}
and hence
\begin{align*}
\frac{\partial^2}{\partial t_1 \partial t_2}\, e_{(Kn-K,Kn]}^\wp&(\gamma(t_1),\gamma(t_2)) \\
&= K^d \int_{n-1}^{n} \nu^{d-1} \frac{\partial^2}{\partial t_1 \partial t_2}\widehat{\sigma_{d-1}}\bigl(K\nu \delta_g(x,y) (t_1-t_2)\bigr) \,\frac{d\nu}{(2\pi)^{\sfrac{d}{2}}}\\
&= -K^{d+2} \delta_g(x,y)^2\int_{n-1}^{n} \nu^{d+1} \widehat{\sigma_{d-1}}''\bigl(K\nu \delta_g(x,y) (t_1-t_2)\bigr) \,\frac{d\nu}{(2\pi)^{\sfrac{d}{2}}}.
\end{align*}
Remember now that $x$ and $y$ satisfy the inequality $\delta_g(x,y)<\sfrac{\beta_0}{Kn}$.
In the last integral, we see that the term $K\nu\delta_g(x,y)(t_1-t_2)$ belongs to the interval $(-\beta_0,\beta_0)$.
Differentiating and taking the real part of the Poisson formula \eqref{Poisson} give\vspace*{-3pt}
\begin{equation*}
-\widehat{\sigma_{d-1}}''(t)=\frac{1}{2^{\sfrac{d}{2}-1}\Gamma(\psfrac{d-1}{2})\sqrt{\pi}}\int_{-1}^{1} \cos(ts)s^2 (1-s^2)^{\psfrac{d-3}{2}} ds.
\end{equation*}
Here is our last restriction for $\beta_0$: due to the continuity of $\widehat{\sigma_{d-1}}''$ at $0$, if $\beta_0$ is small enough then one obviously has
\begin{equation*}
\exists c>0,\ \exists C>0,\ \forall t\in (-\beta_0,\beta_0), \quad c\leq -\widehat{\sigma_{d-1}}''(t)\leq C.
\end{equation*}
These inequalities lead to \eqref{remem2}.

\subsubsection*{Step 6} By plugging \eqref{remem} and \eqref{remem2} in \eqref{doubl}, we get the following inequality (for suitable positive constants $C_1$ and $C_2$):
\begin{multline*}
e_{(Kn-K,Kn]}(x,x)+e_{(Kn-K,Kn]}(y,y)-2e_{(Kn-K,Kn]}(x,y)\\
\geq (C_1 K^{d+2}-C_2 K^{d+1})n^{d+1}\delta_g(x,y)^2.
\end{multline*}
Coming back to \eqref{defideltan2} and \eqref{enxx2}, there is also another positive constant $C_3>0$ satisfying
\begin{equation*}
\delta_n(x,y)^2 \geq \Bigl[\frac{C_1 K^{d+2}-C_2 K^{d+1}}{C_3 K^d} \Bigr]n^2 \delta_g(x,y)^2.
\end{equation*}
A similar argument would show
\begin{equation*}
\delta_n(x,y)^2\leq \Bigl[\frac{C_1 K^{d+2}+C_2 K^{d+1}}{C_3 K^d} \Bigr]n^2 \delta_g(x,y)^2.
\end{equation*}
We conclude by making $K$ tend to $+\infty$.
\end{proof}

Theorem \ref{pseudo-dis} follows by an application of Proposition \ref{minn} for the constant $\beta=\beta_0\in (0,1)$ of Proposition \ref{mino2}. Upon increasing the constants $K_0$ and $c$, we assume that those numbers have the same meaning in
Proposition \ref{minn} and Proposition \ref{mino2}. In the zone $\delta_g\geq \sfrac{\beta_0}{Kn}$, we have
$\sfrac{1}{c} \leq \delta_n \leq c$ and hence
\[
\frac{\beta_0}{c}\,\min(1,Kn\delta_g)\leq \delta_n \leq\frac{c}{\beta_0}\,\min(1,Kn\delta_g).
\]
Similarly, in the zone $\delta_g<\sfrac{\beta_0}{Kn}$, the inequalities
$\Psfrac{1}{c}Kn\delta_g \leq \delta_n\leq c Kn\delta_g$ hold true and so do the following ones
\[
\frac{1}{c}\,\min(1,Kn\delta_g)\leq \delta_n \leq c \min(1,Kn\delta_g).
\]
Hence we get \eqref{var-pdis} in which one may choose $C=\sfrac{c}{\beta_0}$.

\section{Preliminaries for proofs of Theorems \ref{sol} and \ref{sol-q}}\label{s-p2}

For our purpose, a Gaussian process $(f^{G,\omega}(x))_{x\in \mathcal{M}}$ on $(\mathcal{M},\delta_g)$ is defined as a family of real random variables satisfying the following two properties
\begin{enumeratei}
\item for any finite subset $\{x_1,\dots,x_k\}\subset \mathcal{M}$, the random vector
\begin{equation*}\omega\in\Omega\mto (f^{G,\omega}(x_1),\dots,f^{G,\omega}(x_k))\in \R^k \end{equation*}
is a centered Gaussian vector. In other words, any linear combination
\[
\omega\mto \sum_{i=1}^k\alpha_i f^{G,\omega}(x_i),
\]
with $\alpha_i\in \R$, is centered and Gaussian.
\item the function $(x,y)\in \mathcal{M}^2\mto \E_{\omega}[f^{G,\omega}(x)f^{G,\omega}(y)]$, called the covariance structure, is continuous. An equivalent condition is the continuity of the function $x\mto f^{G,\omega}(x)$ from $\mathcal{M}$ to $L^2(\Omega)$.
\end{enumeratei}

Here is the simple example of a Gaussian process we are interested in: for any
continuous function $c:\mathcal{M}\to \ell^2(\N)$ the process $(f^{G,\omega}(x))_{x\in \mathcal{M}}$ given by $f^{G,\omega}(x)=\sum_{\ell\in \N} g_\ell(\omega) c_\ell(x)$ is Gaussian.
As recalled in the introduction, we shall need the following pseudo-distance $\delta$ given by
\begin{equation*}
\forall (x,y)\in \mathcal{M}^2, \quad \delta(x,y):=\sqrt{\E_{\omega}[|f^{G,\omega}(x)-f^{G,\omega}(y)|^2]}.
\end{equation*}

We note that if the Gaussian random series $\sum f_n^{G,\omega}$ defined in \eqref{ran-fouG2} almost surely converges in $\CC^0(\mathcal{M}_s)$ (for a submanifold $\mathcal{M}_s$ of $\mathcal{M}$) then the random series $\sum f_n^{G,\omega}(x)$ almost surely converges for any $x\in \mathcal{M}_s$.
For the sake of clarity, we state the following result.

\begin{prop}\label{con-nec}
There is a constant $K_0\geq 1$ merely depending on the Riemannian manifold $\mathcal{M}$ such that for any $K \geq K_0$ and for any sequence $(f_n)_{n\geq 1}$ satisfying $f_n\in E_{(Kn-K,Kn]}$ for $n\geq 1$ (see \eqref{defiEn}), then the following assertions are equivalent:
\begin{enumeratei}
\item\label{con-neci} the series $\sum \left\Vert f_n \right\Vert_{L^2(\mathcal{M})}^2$ is convergent,
\item\label{con-necii} for any $x\in \mathcal{M}$, the random series $\sum f_n^{G,\omega}(x)$ converges in $L^2(\Omega)$ to a Gaussian random variable,
\item\label{con-neciii} for any $x\in \mathcal{M}$, the random series $\sum f_n^{G,\omega}(x)$ almost surely converges in $\R$,
\item\label{con-neciv} there is $x\in \mathcal{M}$ such that the random series $\sum f_n^{G,\omega}(x)$ almost surely converges in $\R$.
\end{enumeratei}
\noindent
Moreover, the previous statements imply the following one:
\begin{enumeratei}\setcounter{enumi}{4}
\item\label{con-necv}
the family of random variables
$f^{G,\omega}(x)=\sum_{n\geq 1} f_n^{G,\omega}(x)$ is a Gaussian process on $\mathcal{M}$ and its Dudley pseudo-distance $\delta:\mathcal{M}^2\!\to\! [0,+\infty)$ is given, for any $(x,y)\!\in\!\mathcal{M}^2$,~by%
\begin{equation}\label{ddu}
\delta(x,y)^2=\E_{\omega}\bigl[\left\vert f^{G,\omega}(x)-f^{G,\omega}(y) \right\vert^2\bigr]=\sum_{n\geq 1} \left\Vert f_n \right\Vert_{L^2(\mathcal{M})}^2 \delta_n(x,y)^2,
\end{equation}
where the partial pseudo-distance $\delta_n$ is defined in \eqref{defideltan}.
\end{enumeratei}
Finally, in the specific case $\mathcal{M}=\bbS^d$, a similar statement holds true if each $f_n$ belongs to the eigenspace $\ker(\Deltag+n(n+d-1))$ provided that we replace $f_n^{G,\omega}$ with~\eqref{ran-fouG} and $\delta_n$ with \eqref{defipseu}.
\end{prop}

\skpt
\begin{proof}\vspace*{-5pt}
\subsubsection*{$\eqref{con-neci}\Leftrightarrow\eqref{con-necii}$}
For any $x\in \mathcal{M}$, the random variable $f_n^{G,\omega}(x)$ is centered and Gaussian with
variance equaling
\begin{equation*}
\E_{\omega}\bigl[|f_n^{G,\omega}(x)|^2\bigr]=\frac{\|f_n\|_{L^2(\mathcal{M})}^2}{\dim(E_{(Kn-K,Kn]})}\sum_{\lambda_k\in (Kn-K,Kn]} \phi_k(x)^2.
\end{equation*}
By orthogonality of the random variables $(f_n^{G,\omega}(x))_{n\geq 1}$ in $L^2(\Omega)$, it is clear and well-known that the Gaussian series $\sum f_n^{G,\omega}(x)$ converges in $L^2(\Omega)$ if and only if the series $\sum \E_{\omega}[|f_n^{G,\omega}(x)|^2]$ converges and we have
\begin{equation*}
\E_{\omega}\biggl[\Big| \sum_{n\geq 1}f_n^{G,\omega}(x) \Big|^2 \biggr]=\sum_{n\geq 1} \E_{\omega}[|f_n^{G,\omega}(x)|^2].
\end{equation*}
Thanks to \eqref{enxx} and \eqref{enxx2}, we get the equivalence $\eqref{con-neci}\Leftrightarrow\eqref{con-necii}$ and
there are two constants $C\geq 1$ and $K_0\geq 1$ that merely depend on the Riemannian structure of $\mathcal{M}$ such that for any $K\geq K_0$ we have
\begin{equation}\label{gaau2}
\frac{1}{C}\sum_{n\geq 1} \left\Vert f_n\right\Vert_{L^2(\mathcal{M})}^2\leq \E_{\omega}
\biggl[\Big|\sum_{n\geq 1} f_n^{G,\omega}(x) \Big|^2 \biggr]
\leq C \sum_{n\geq 1} \left\Vert f_n\right\Vert_{L^2(\mathcal{M})}^2.
\end{equation}

\subsubsection*{$\eqref{con-neciii}\Leftrightarrow\eqref{con-neciv}$}
Obvious.

\subsubsection*{$\eqref{con-neci}\Leftrightarrow\eqref{con-neciii}$ and $\eqref{con-neciv}\Leftrightarrow\eqref{con-neci}$} See for instance \cite[Cor.\,III.6, p.\,26]{queff2} and the previous computations.

\subsubsection*{Proof of $\eqref{con-necv}$} The finite linear combination $\sum_{i=1}^{p} \alpha_i \bigl(\sum_{n\geq 1}f_n^{G,\omega}(x_i)\bigr)$, with $\alpha_i\in \R$, is Gaussian because it is the limit in $L^2(\Omega)$ of the Gaussian random variables
\begin{equation*}
\sum_{n=1}^N \biggl(\sum_{i=1}^p \alpha_i f_n^{G,\omega}(x_i) \biggr)
\end{equation*} as $N\to +\infty$ (with the same argument as in $\eqref{con-neci}\Leftrightarrow\eqref{con-neciii}$).

We now have to check the continuity of the covariance structure.
By using the notation \eqref{defideltan} and orthogonality arguments, the definition \eqref{ran-fouG2} of $f_n^{G,\omega}$ directly proves~\eqref{ddu}.
We recall that the expected continuity is equivalent to the continuity of the function $x\in \mathcal{M}\mto f^{G,\omega}(x)\in L^2(\Omega)$.
But the partial pseudo-distance $\delta_n$ are continuous on $\mathcal{M}^2$ and uniformly bounded on $\mathcal{M}^2$ (see \eqref{uniff}).
Hence, the right-hand side of \eqref{ddu} absolutely converges in $\CC^0(\mathcal{M}^2)$ and we obtain the expected continuity.

\subsubsection*{Towards the case $\mathcal{M}=\bbS^d$} For the sphere $\bbS^d$, we work with \eqref{ran-fouG}. In the same spirit as \eqref{KN}, we have
\begin{equation*}
\sum_{\lambda_k^2=n(n+d-1)}\phi_k(x)^2=\frac{\dim(E_n)}{\operatorname{vol}_{\bbS^d}(\bbS^d)}.
\end{equation*}
From such an identity, we see that the equivalence \eqref{gaau2} becomes
\begin{equation*}
\E_{\omega}\biggl[\Big| \sum_{n\geq 1}f_{n}^{G,\omega}(x) \Big|^2 \biggr]=\frac{1}{\operatorname{vol}_{\bbS^d}(\bbS^d)} \sum_{n\geq 1} \|f_n\|_{L^2(\bbS^d)}^2.
\end{equation*}
We also deduce that the partial pseudo-distances $\delta_n$, associated to the sequence of the eigenspaces of $\Deltag$ and defined in \eqref{defipseu}, are uniformly bounded on $\bbS^d\times \bbS^d$ (an information used in the proof of \eqref{con-necv}):
\begin{equation}\label{idd}
\delta_n\leq \frac{2}{\sqrt{\operatorname{vol}_{\bbS^d}(\bbS^d) }}.\qedhere
\end{equation}
\end{proof}

\section{Proof of Theorem \ref{sol} for spheres \texorpdfstring{$\bbS^d$ with $d\geq 2$}{Sd}}\label{preuv-solSd}

Thanks to the convergence of $\sum_{n\geq 1} \left\Vert f_n\right\Vert_{L^2(\bbS^d)}^2$ and Point v) of Proposition \ref{con-nec}, we know that the process $\bigl(\sum_{n\geq 1} f_n^{G,\omega}(x)\bigr)_{x\in \bbS^d}$ is Gaussian and so is its restriction to any compact submanifold $\mathcal{M}_s$.
We begin by the following proposition.

\begin{prop}\label{casSd}
Let us consider a Gaussian process $(F^\omega(x))_{x\in \bbS^d}$ whose covariance structure reads
\begin{equation}\label{covSD}
\forall (x,y)\in \bbS^d\times \bbS^d, \quad \E_\omega[F^\omega(x) F^\omega(y)]=\mathscr{C}(\langle x, y\rangle)
\end{equation}
for some continuous function $\mathscr{C}:[-1,1]\to \R$. Then the following statement are equivalent:
\begin{enumeratei}
\item\label{casSdi} the condition $\int_{0}^{1} \Psfrac{(\mathscr{C}(1)-\mathscr{C}(\cos(t)))^{1/2}}{t\sqrt{\ln(1/t)}}\, dt<+\infty$ is fulfilled,
\item\label{casSdii} the Gaussian process $(F^\omega(x))_{x\in \bbS^d}$ admits a version $(\widetilde{F}^\omega(x))_{x\in \bbS^d}$ which is sample-continuous in the following sense:
\begin{itemize}
\item with probability one, the function $x\mto \widetilde{F}^{\omega}(x)$ is continuous from
$\bbS^d$ to $\R$,
\item for any $x\in \mathcal{M}$, the equality $\P[F^{\omega}(x)=\widetilde{F}^{\omega}(x)]=1$ holds true.
\end{itemize}
\item\label{casSdiii} on some fixed geodesic arc $\Gamma$ of $\bbS^d$, the Gaussian process $(F^\omega(x))_{x\in \Gamma}$ process on $\Gamma$ admits a version which is sample-bounded in the following sense:
\begin{itemize}
\item with probability $1$, the function $x \mto \widetilde{F}^{\omega}(x)$ is bounded from $\Gamma$ to $\R$,
\item for any $x\in \mathcal{M}$, the equality $\P[F^{\omega}(x)=\widetilde{F}^{\omega}(x)]=1$ holds true.
\end{itemize}
\end{enumeratei}
\end{prop}

\skpt
\begin{proof}\vspace*{-.5\baselineskip}
\subsubsection*{Preliminary remarks} Note that due to the definition \eqref{covSD} of the function $\mathscr{C}$ and to the Cauchy-Schwarz inequality, we easily check the inequality $|\mathscr{C}(u)|\leq \mathscr{C}(1)$ for any $u\in [-1,1]$ (so the square root in the integral in Point i) is well defined). Then one may compute the Dudley pseudo-distance $\delta$ of the Gaussian process $(F^\omega(x))_{x\in \bbS^d}$:
\begin{equation*}
\delta(x,y):=\sqrt{\E_\omega[|F^\omega(x)-F^\omega(y)|^2]}=\sqrt{2\mathscr{C}(1)-2\mathscr{C}(\langle x,y\rangle)}.
\end{equation*}
We now note that for any two points $x$ and $y$ on the sphere $\bbS^d$ we have $\langle x,y\rangle=\cos(\delta_g(x,y))$ where $\delta_g(x,y)$ stands for the Riemannian distance on $\bbS^d$.
In other words, we have
\begin{equation}\label{delt-scr}
\delta(x,y)=\sqrt{2\mathscr{C}(1)-2\mathscr{C}(\cos(\delta_g(x,y)))}.
\end{equation}
Now from the classical result \cite[p.\,101, Th.\,1]{schoenberg1942} by Schoenberg, there exists a sequence of non-negative coefficients $(u_n)_{n\geq 0}$ satisfying
\begin{equation*}
\forall \theta\in \R,\quad \mathscr{C}(\cos(\theta))=\sum_{n=0}^{+\infty} u_n
P_n^{(\sfrac{d}{2}-1,\sfrac{d}{2}-1)}(\cos(\theta))
\end{equation*}
and such that series converges at $\theta=0$ (in other words $\sum_{n=0}^{+\infty} n^{\sfrac{d}{2}-1}u_n <+\infty$).

\subsubsection*{$\eqref{casSdi}\Rightarrow\eqref{casSdii}$}
We may cover $\bbS^d$ by a finite set of closed Riemannian balls:
\begin{equation*}
\bbS^d= \bigcup_{1\leq k\leq N} \overline{B_k} \quad \text{with $B_k:=B_{\delta_g}
(a_k,\sfrac{\pi}{4})$ and $a_k\in \bbS^d$.}
\end{equation*}
In order to show \eqref{casSdii}, it is clear that it is sufficient to prove that each Gaussian process $(F^{\omega}(x))_{x\in \overline{B_k}}$, for $1\leq k\leq N$, admits a sample-continuous version.
Thanks to the Dudley theorem, it is sufficient to show the entropy conditions of the Dudley pseudo-distances of the $N$ Gaussian processes $(F^{\omega}(x))_{x\in \overline{B_k}}$.
Since each closed ball $\overline{B_k}$ has a diameter equaling $\sfrac{\pi}{2}<\pi$,
we may use Proposition \ref{jacobi} (with $\vartheta=\sfrac{\pi}{4}$) to get the equivalence (up to multiplicative constants independent of $x$ and $y$):
\begin{align}\notag
\delta(x,y)^2 &=2 \sum_{n=1}^{+\infty} u_n
\left[
P_n^{(\sfrac{d}{2}-1,\sfrac{d}{2}-1)}(1)
-P_n^{(\sfrac{d}{2}-1,\sfrac{d}{2}-1)}(\cos(\delta_g(x,y)))
\right]\\
\delta(x,y)^2 & \simeq \sum_{n=1}^{+\infty} u_n n^{\sfrac{d}{2}-1} \min(1,n \delta_g(x,y))^2.\label{delt-scr2}
\end{align}
The function $\Upsilon:t\in [0,+\infty)\mto \sqrt{\sum_{n\geq 1} u_n n^{\sfrac{d}{2}-1}\min(1,n t)^2}$ fulfills the assumption of Proposition \ref{lem-ref}. Note that the probability Riemannian volume $\mu$ of $\bbS^d$ satisfies \eqref{ln-hom} (with suitable possibly different constants) on the ball $\overline{B_k}$ since $\overline{B_k}$ is a submanifold with smooth boundary. Then the conclusion of Proposition \ref{lem-ref} combined with \eqref{delt-scr} and~\eqref{delt-scr2} shows that the entropy condition of each $(F^\omega(x))_{x\in \overline{B}_k}$ is equivalent to Point~\eqref{casSdi}.

\subsubsection*{$\eqref{casSdii}\Rightarrow\eqref{casSdiii}$}
Obvious.

\subsubsection*{$\eqref{casSdiii}\Rightarrow\eqref{casSdi}$}
The main point is that \eqref{covSD} implies that the restriction of the Gaussian process $(F^\omega(x))_{x\in\bbS^d}$ is stationary on any geodesic closed curve (identified as the torus~$\mathbb{T}$).
By applying a finite set of translations to the Gaussian process $(F^\omega(x))_{x\in \Gamma}$, one obtains the almost sure boundedness on the whole geodesic curve containing $\Gamma$.
So there is no loss of generality to assume that $\Gamma$ is the whole geodesic curve of the points $x_\theta=(\cos(\theta),\sin(\theta),0,\dots,0)$ for $\theta \in [0,2\pi]$.
But the Dudley-Fernique theorem ensures that the almost sure boundedness is equivalent to the entropy condition on $\Gamma$.
In order to use Proposition \ref{lem-ref} and the equivalence \eqref{delt-scr2} of the Dudley pseudo-distance, we need to restrict to a subcurve of $\Gamma$ on which the geodesic distance is strictly less that $\pi$, for instance
\begin{equation}\label{gammapl}
\Gamma^{+}:=\bigl\{(\cos(\theta),\sin(\theta),0,\dots,0),\quad \theta \in[0,\sfrac{\pi}{2}] \bigr\}.
\end{equation}
It is easy to see that the entropy condition on $\Gamma^{+}$ is still satisfied.\footnote{We indeed have the inequality $N_{\delta}(\Gamma^+,2\ep)\leq N_{\delta}(\Gamma,\ep)$. For the proof, we consider a finite cover $\Gamma=\bigcup B_{\delta}(x_i,\ep)$ with $x_i\in \Gamma$.
And so we have $\Gamma^+=\bigcup (B_{\delta}(x_i,\ep)\cap \Gamma^+)$.
If $B_{\delta}(x_i,\ep)$ intersects $\Gamma^+$ then any element $x_i'\in B_{\delta}(x_i,\ep)\cap \Gamma^{+}$ satisfies $B_{\delta}(x_i,\ep)\subset B_{\delta}(x_i',2\ep)$. Hence we obtain a cover $\Gamma^+\subset \bigcup B_{\delta}(x_i',2\ep)$.}
Let us now denote by $\delta_{\Gamma^+}$ the Dudley pseudo-distance on $\Gamma^+$ of the Gaussian process $(F^\omega(x))_{x\in \Gamma^+}$. Due to the equivalence \eqref{delt-scr2}, we still have $\delta_{\Gamma^+}\simeq \Upsilon(\delta_g)$ (where $\delta_g$ equals both the geodesic distance on $\Gamma^{+}$ and of course the geodesic distance on $\bbS^d$ due to the definition of $\Gamma^{+}$). Moreover, the length measure of $\Gamma^+$ obviously satisfies inequalities like \eqref{ln-hom} (with $\sigma=1$). As above, the conclusion of Proposition \ref{lem-ref} shows Point \eqref{casSdi}.
\end{proof}

In the sequel of this part, we shall consider the case of the Gaussian process $f^{G,\omega}=\sum_{n\geq 1}f_n^{G,\omega}$ defined in \eqref{ran-fouG} with $\sum_{n\geq 1} \|f_n\|_{L^2(\bbS^d)}^2<+\infty$. The proof of Proposition \ref{con-nec} easily gives the formula
\begin{equation*}
\forall (x,y)\in \bbS^d\times \bbS^d, \quad \E[f^{G,\omega}(x)f^{G,\omega}(y)]=\sum_{n\geq 1} \frac{\|f_n\|_{L^2(\bbS^d)}^2}{\dim(E_n)} \sum_{\lambda_k^2=n(n+d-1)} \phi_k(x)\phi_k(y).
\end{equation*}
Hence, \eqref{KN} shows that the previous term equals $\mathscr{C}(\langle x,y\rangle)$ with
\begin{equation*}
\mathscr{C}(\theta)=\sum_{n\geq 1}\frac{\|f_n\|_{L^2(\bbS^d)}^2}{\dim(E_n)}\, c_{d,n} P_n^{(\sfrac{d}{2}-1,\sfrac{d}{2}-1)}(\cos(\theta))
\end{equation*}
(where $c_{d,n}$ is given by \eqref{equivcn} and satisfies $c_{d,n}\eqsn n^{\sfrac{d}{2}}$ and $\dim(E_n)\eqsn n^{d-1}$) and the function $\Upsilon:[0,+\infty[\to [0,+\infty[$ of the previous proof can be chosen as
\begin{equation*}
\Upsilon(t)=\sqrt{\sum_{n\geq 1} \|f_n\|_{L^2(\bbS^d)}^2 \min(1,nt)^2}.
\end{equation*}
Moreover, the Dudley pseudo-distance of $(f^{G,\omega}(x))_{x\in \bbS^d}$ satisfies
\begin{equation}\label{dudd}
\delta(x,y)\simeq \Upsilon(\delta_g(x,y))
\end{equation}
uniformly in $x$ and $y$ provided that $\delta_g(x,y)\leq\sfrac{\pi}{2}$ (that is a consequence of Theorem~\ref{pro-s+} and \eqref{ddu}).
We now prove Theorem \ref{sol} and so we are interested in studying the almost sure continuity of $x\in \mathcal{M}_s\mto \sum_{n\geq 1}f_n^{G,\omega}(x)\in \R$.

\begin{prop}\label{ieii}
The equivalence $\eqref{soli}\Leftrightarrow\eqref{solii}$ of Theorem \ref{sol} is true.
\end{prop}

\skpt
\begin{proof}\vspace*{-.5\baselineskip}
\subsubsection*{Preliminary remarks}
Let us recall why the existence of a sample-continuous version of the Gaussian process $\bigl(\sum_{n\geq 1}f_n^{G,\omega}(x)\bigr)_{x\in \mathcal{M}_s}$ on $\mathcal{M}_s$ is equivalent to the almost sure convergence of the Gaussian random series $\sum f_n^{G,\omega}$ in $\CC^0(\mathcal{M}_s)$, namely Point \eqref{soli} of Theorem \ref{sol}. That is the same argument than that the one used in the classical theory for compact groups. Actually, we apply the Itô-Nisio theorem (see \cite[p.\,238]{queff2}) that merely needs to check the additional assumption ensuring that, for any $x\in \mathcal{M}_s$, each random numerical series $\sum f_n^{G,\omega}(x)$ is almost surely convergent. Such a fact is actually proved in Point iii) of Proposition \ref{con-nec}.

\Subsubsection*{Step 1}
 For $\eqref{soli}\Rightarrow\eqref{solii}$, we may clearly assume that $\mathcal{M}_s$ is a curve with ends. We can also assume that the diameter of $\mathcal{M}_s$ is less than $\sfrac{\pi}{2}$. Let $m_s: [0,\sfrac{\pi}{2}]\to \mathcal{M}_s$ be a bi-Lipschitz parametrization and let us consider the Gaussian process $(f^{G,\omega}(m_s(\theta)))_{\theta \in [0,\sfrac{\pi}{2}]}$. We~clearly have the equivalence
\begin{equation*}
\E\bigl[|f^{G,\omega}(m_s(\theta))-f^{G,\omega}(m_s(\theta'))|^2\bigr]\simeq \Upsilon(\delta_g(m_s(\theta),m_s(\theta')))^2\simeq \Upsilon(|\theta-\theta'|)^2.
\end{equation*}
Now let $\gamma:[0,\sfrac{\pi}{2}]\to \Gamma^{+}$ be the natural parametrization of $\Gamma^{+}$ defined in \eqref{gammapl}. We~also have (see \eqref{dudd}):
\begin{equation}\label{defSS}
\E\bigl[|f^{G,\omega}(\gamma(\theta))-f^{G,\omega}(\gamma(\theta'))|^2\bigr]\simeq \Upsilon(|\theta-\theta'|)^2.
\end{equation}
The classical Slepian comparison theorem ensures that the almost sure boundedness of the Gaussian processes $(f^{G,\omega}(m_s(\theta)))_{\theta \in [0,\sfrac{\pi}{2}]}$ is equivalent to that of
$(f^{G,\omega}(\gamma(\theta)))_{\theta\in [0,\sfrac{\pi}{2}]}$ (see \cite[Cor.\,3.14]{ledoux} or \cite[p.\,73]{queff2}), and so the same is obviously true for $(f^{G,\omega}(x))_{x\in \mathcal{M}_s}$ and $(f^{G,\omega}(x))_{x\in \Gamma^+}$. Then Proposition \ref{casSd} (or its proof) implies the following convergence
\begin{equation}\label{uup}
\int_{0}^1 \frac{\Upsilon(t)}{t\sqrt{\ln(1/t)}}dt <+\infty.
\end{equation}
For the conclusion, we need to admit for a moment Step 2.

\subsubsection*{$\eqref{soli}\Rightarrow\eqref{solii}$} That is a direct consequence of Step 2.

\subsubsection*{$\eqref{solii}\Rightarrow\eqref{soli}$} By Step 2 we obtain \eqref{uup}, which is equivalent to Point i) of Proposition \ref{casSd} (via \eqref{dudd}) which in turns proves the almost sure continuity of $x\mto f^{G,\omega}(x)$ (see the preliminary remarks above) on the whole sphere $\bbS^d$ thanks to Point \eqref{casSdii} of Proposition~\ref{casSd}.

\subsubsection*{Step 2} We shall prove that \eqref{uup} is equivalent to the Salem-Zygmund condition~\eqref{pz-con3}. We begin by writing
\begin{align*}
\int_{0}^1 \sqrt{\sum_{n\geq 1} \left\Vert f_n\right\Vert_{L^2(\bbS^d)}^2 \min(1,n^2 t^2)}&\,\frac{dt}{t\sqrt{-\ln(t)}} \\
&\hspace*{-1cm}=
\sum_{p\geq 1} \int_{\spfrac{1}{p+1}}^{\sfrac{1}{p}}\sqrt{\sum_{n\geq 1} \left\Vert f_n\right\Vert_{L^2(S^d)}^2 \min(1,n^2 t^2)}\,\frac{dt}{t\sqrt{-\ln(t)}} \\
&\hspace*{-1cm}\simeq
\sum_{p\geq 1} \sqrt{\sum_{n\geq 1} \left\Vert f_n\right\Vert_{L^2(\bbS^d)}^2 \min(1,\sfrac{n^2}{p^2})} \,\frac{1}{p\sqrt{\ln(p+1)}}.
\end{align*}
We now write
\begin{equation*}
\sqrt{\sum_{n\geq 1} \left\Vert f_n\right\Vert_{L^2(\bbS^d)}^2 \min(1,\sfrac{n^2}{p^2})}=\sqrt{\frac{1}{p^2}\,U_p+V_p}\simeq \frac{\sqrt{U_p}}{p}+\sqrt{V_p},
\end{equation*}
where $U_p$ and $V_p$ are defined as follows
\begin{equation}\label{defupvp}
U_p=\sum_{n=1}^{p-1} n^2 \left\Vert f_n\right\Vert_{L^2(\bbS^d)}^2 \quad \text{and}\quad V_p=\sum_{n=p}^{+\infty} \left\Vert f_n\right\Vert_{L^2(\bbS^d)}^2, \end{equation}
with the convention $U_1=0$. Hence we get
\begin{equation*}
\int_{0}^1 \frac{\Upsilon(t)}{t\sqrt{-\ln(t)}}\,dt
\simeq \sum_{p\geq 1} \frac{\sqrt{U_p}}{p^2\sqrt{\ln(p+1)}}+\frac{\sqrt{V_p}}{p \sqrt{\ln(p+1)}} .
\end{equation*}
The $V_p$ part is exactly the Salem-Zygmund type term in \eqref{pz-con3}. So the proof will be finished provided that we show that the contribution of the $U_p$ part is controlled by that of the $V_p$ part:
\begin{equation}\label{ddom}
\sum_{p\geq 2} \frac{\sqrt{U_p}}{p^2\sqrt{\ln(p+1)}}\lesssim \sum_{p\geq 1} \frac{\sqrt{V_p}}{p\sqrt{\ln(p+1)}}.
\end{equation}
Thanks to the Cauchy-Schwarz inequality and the definition of $U_p$ in \eqref{defupvp}, we have
\begin{align*}
\sum_{p\geq 2} \frac{\sqrt{U_p}}{p^2\sqrt{\ln(p+1)}} & = \sum_{p\geq 2} \frac{1}{\sqrt{p}\ln(p+1)}\times \frac{\sqrt{U_p} \sqrt{\ln(p+1)}}{p^{3/2}} \\
& \leq C \biggl(\sum_{p\geq 2} \frac{U_p \ln(p+1)}{p^{3}} \biggr)^{\sfrac{1}{2}}\\
& \leq C\biggl(\sum_{n\geq 1} \left\Vert f_n \right\Vert_{L^2(\bbS^d)}^2 n^2 \sum_{p\geq n+1} \frac{\ln(p+1)}{p^3} \biggr)^{\sfrac{1}{2}}\\
& \leq C\biggl(\sum_{n=1}^{+\infty} \left\Vert f_n\right\Vert_{L^2(\mathcal{M)}^2 }^2\ln(n+1) \biggr)^{1/2}.
\end{align*}
By using the equivalence $\sqrt{\ln(n+1)}\simeq \sum_{p=1}^{n} \sfrac{1}{p\sqrt{\ln(p+1)}}$ and the definition of $V_p$ in \eqref{defupvp}, we obtain
\begin{align*}
\sum_{p\geq 1} \frac{\sqrt{U_p}}{p^2\sqrt{\ln(p+1)}} & \leq C\biggl(\sum_{n=1}^{+\infty} \|f_n\|_{L^2(\bbS^d)}^2 \biggl(\sum_{p=1}^n \frac{1}{p\sqrt{\ln(p+1)}} \biggr)^2 \biggr)^{1/2} \\
& \leq C\biggl(\sum_{n=1}^{+\infty} \|f_n\|_{L^2(\bbS^d)}^2 \sum_{p=1}^n \sum_{q=1}^n \frac{1}{p\sqrt{\ln(p+1)}q\sqrt{\ln(q+1)}} \biggr)^{1/2} \\
& \leq C\biggl(\sum_{p=1}^{+\infty} \sum_{q=1}^{+\infty} \frac{V_{\max(p,q)}}{p\sqrt{\ln(p+1)}q\sqrt{\ln(q+1)}} \biggr)^{1/2}.
\end{align*}
By bounding $V_{\max(p,q)}\leq \sqrt{V_p V_q}$, we finally prove \eqref{ddom}.
\end{proof}

The following proposition finishes the proof of Theorem \ref{sol} for spheres.

\begin{prop}\label{ieii2}
The equivalence $\eqref{soli}\Leftrightarrow\eqref{soliii}$ is true in Theorem \ref{sol}.
\end{prop}
\begin{proof}
The implication $\eqref{soliii}\Rightarrow\eqref{soli}$ is nothing other than the Dudley theorem.
Let us prove the converse implication$\eqref{soli}\Rightarrow\eqref{soliii}$. Let us set $\sigma=\dim(\mathcal{M}_s)$. By compactness, we can decompose $\mathcal{M}_s$ as a union of $\sigma$-dimensional compact submanifolds $\mathcal{M}_{s,1},\dots,\mathcal{M}_{s,N}$ with smooth boundary and with diameter less than $\sfrac{\pi}{2}$:
\begin{equation}\label{cover-Ms}
\mathcal{M}_s=\mathcal{M}_{s,1}\cup \dots \cup \mathcal{M}_{s,N}, \quad \max_{\substack{x\in \mathcal{M}_{s,k}\\ y\in \mathcal{M}_{s,k}}} \delta_g(x,y)\leq \frac{\pi}{2}, \quad \forall k\in \{1,\dots,N\}.
\end{equation}
Point \eqref{soli} of Theorem \ref{sol} implies the almost sure continuity of each $x\in \mathcal{M}_{s,k}\mto f^{G,\omega}(x)$ for any $k\in \{1,\dots,N\}$.
Proposition \ref{ieii} (or its proof) shows that the following condition holds
\begin{equation*}
\int_{0}^1 \frac{\Upsilon(t)}{t\sqrt{\ln(1/t)}}\,dt<+\infty.
\end{equation*}
Note that the probability Riemannian volume of each $\mathcal{M}_{s,k}$ satisfies the assumption~\eqref{ln-hom}.
As a consequence of Proposition \ref{lem-ref}, we obtain the entropy condition for each $\mathcal{M}_{s,k}$:
\begin{equation*}
\int_{0}^{+\infty} \sqrt{\ln(N_{\delta} (\mathcal{M}_{s,k},\ep))}\,d\ep <+\infty.
\end{equation*}
The definition of the covering number and \eqref{cover-Ms} lead to the following inequality
\begin{equation}\label{min-NN}
N_{\delta}(\mathcal{M}_s,\ep) \leq \sum_{1\leq k\leq N}N_{\delta}(\mathcal{M}_{s,k},\ep),
\end{equation}
which in turn implies
\begin{equation}\label{min-NN2}
N_{\delta}(\mathcal{M}_{s},\ep) \leq N\times \prod_{1\leq k\leq N}N_{\delta}(\mathcal{M}_{s,k} ,\ep).
\end{equation}
We then easily deduce Point \eqref{soliii} of Theorem \ref{sol}.
\end{proof}

\section{Proof of Theorem \ref{sol-q} for \texorpdfstring{$\bbS^d$ with $d\geq 2$}{Sd2}}\label{preuv-solq}

The proof developed in the last part already shows that the following two numbers are equivalent (up to multiplicative constants):
\begin{enumeratei}
\item the number $ \int_{0}^1 \Psfrac{\Upsilon(t)}{t\sqrt{\ln(1/t)}}\, dt$,
\item the Salem-Zygmund expression $\sum_{p\geq 1} \Psfrac{1}{p\sqrt{\ln(p+1)}} \bigl(\sum_{n\geq p} \|f_n\|_{L^2(\bbS^d)}^2 \bigr)^{1/2}$.
\end{enumeratei}
In order to add the the entropy integral $\int_{0}^{+\infty} \sqrt{\ln(N_{\delta}(\mathcal{M}_s,\ep))}\,d\ep$ to that list, we~merely need the following two arguments.

\begin{enumerate}\renewcommand{\theenumi}{\alph{enumi}}
\item
The entropy integral on $\mathcal{M}_s$ is equivalent to $\sup_{1\leq k\leq N}
\int_{0}^{+\infty} \hspace*{-2mm}\sqrt{\ln(N_{\delta}(\mathcal{M}_{s,k},\ep))}\,d\ep$. To see that equivalence, one firstly precises \eqref{min-NN} as follows
\begin{equation*}
\max_{1\leq k\leq N} N_{\delta}(\mathcal{M}_{s,k},2\ep) \leq N_{\delta}(\mathcal{M}_s,\ep),
\end{equation*}
which can be proved as explained after \eqref{gammapl}. Then, we easily can bound from below the
entropy integral $
\int_{0}^{+\infty} \sqrt{\ln(N_{\delta}(\mathcal{M}_s,\ep))}\,d\ep$. To bound from above the entropy integral on $\mathcal{M}_s$ by $\sup_{1\leq k\leq N}
\int_{0}^{+\infty} \sqrt{\ln(N_{\delta}(\mathcal{M}_{s,k},\ep))}\,d\ep$, we may use \eqref{min-NN2} by remembering that~$N$ in
\eqref{min-NN2} can be chosen merely depending on the submanifold $\mathcal{M}_s$ and that the diameter of $(\mathcal{M}_s,\delta)$ is bounded by $C(d)\sqrt{\lim_{n\geq 1} \|f_n\|_{L^2(\bbS^d)}^2}$ (see \eqref{idd}).

\item
Thanks to Corollary \ref{lem-quant}, each entropy integral
$\int_{0}^{+\infty} \sqrt{\ln(N_{\delta}(\mathcal{M}_{s,k},\ep))}\,d\ep$ of $\mathcal{M}_{s,k}$ is equivalent to the same number $ \int_{0}^1 \Psfrac{\Upsilon(t)}{t\sqrt{\ln(1/t)}} \,dt$.
\end{enumerate}

The last thing to explain is why we can equivalently control the expectation
\begin{equation*}
\E_{\omega}\biggl[\sup_{x\in \mathcal{M}_s}\Big\vert\sum_{n\geq 1} f_n^{G,\omega}(x)\Big\vert \biggr].
\end{equation*}
The quantitative version of the Dudley theorem will allow to bound from above the last expectation by the entropy integral (see \cite[Th.\,11.17]{ledoux}): we fix $x_0\in \mathcal{M}_s$ and write (thanks to the symmetry of the Gaussian process):
\begin{align*}
\E_{\omega}\biggl[\sup_{x\in \mathcal{M}_s}\!\Bigl\vert\sum_{n\geq 1} f_n^{G,\omega}(x)\Bigr\vert \biggr]&\leq
\E_{\omega}\biggl[\Bigl\vert\sum_{n\geq 1} f_n^{G,\omega}(x_0)\Bigr\vert \biggr]+
\E_{\omega}\biggl[\sup_{x,y\in \mathcal{M}_s}\!\Bigl\vert\sum_{n\geq 1} f_n^{G,\omega}(x)-f_n^{G,\omega}(y)\Bigr\vert \biggr] \\
&\leq
\E_{\omega}\biggl[\Big\vert\sum_{n\geq 1} f_n^{G,\omega}(x_0)\Big\vert \biggr]+2
\E_{\omega}\biggl[\sup_{x\in \mathcal{M}_s}\sum_{n\geq 1} f_n^{G,\omega}(x) \biggr].
\end{align*}
The last expectation is controlled by the entropy integral (\cite[Th.\,11.17]{ledoux}) whereas the first is controlled as follows (see the proof of Proposition \ref{con-nec}):
\begin{align*}
\E_{\omega}\biggl[\Big\vert\sum_{n\geq 1} f_n^{G,\omega}(x_0)\Big\vert \biggr] & \leq \E_{\omega}\biggl[\Big\vert\sum_{n\geq 1} f_n^{G,\omega}(x_0)\Big\vert \biggr]^{\sfrac{1}{2}} \leq C \biggl(\sum_{n\geq 1} \|f_n\|_{L^2(\bbS^d)}^2 \biggr)^{1/2} \\
& \leq C\sqrt{\ln(2)}\sum_{p\geq 1} \frac{1}{p\sqrt{\ln(p+1)}}\, \biggl(\sum_{n\geq p} \|f_n\|_{L^2(\bbS^d)}^2 \biggr)^{1/2} \quad \text{(see $p=1$)}
\end{align*}
and then we note that we already proved that the Salem-Zygmund term is controlled by the entropy integral at the beginning of this part.

For bounding $\E_{\omega}\bigl[\sup_{x\in \mathcal{M}_s}\bigl\vert\sum_{n\geq 1} f_n^{G,\omega}(x)\bigr\vert_{\mathcal{C}^{0}(\mathcal{M}_s)} \bigr]$ from below, we may use the quantitative version of the Fernique theorem (see \cite[Th.\,13.3]{ledoux}) and the Slepian comparison theorem (see the proofs of $\eqref{casSdiii}\Rightarrow\eqref{casSdi}$ in Proposition \ref{casSd} and Step 1 in Proposition \ref{ieii}).

\section{Proof of Theorem \ref{sol} and \ref{sol-q} for compact manifolds}\label{preuv-solM}

Here $f^{G,\omega}$ is defined in \eqref{ran-fouG2} and so the Dudley pseudo-distance $\delta$ implicitly depends on $K$ since $f_n^{G,\omega}$ is a random wave with values in $E_{(Kn-K,Kn]}$.
We assume that $K$ is large enough so that the conclusion of Theorem \ref{pseudo-dis} holds true.

\begin{itemize}
\item
For the proof of the analogue of Proposition \ref{ieii}, the mere thing to clarify is the proof of Step 1. Actually, we can repeat some arguments with
\begin{equation*}
\Upsilon(t)=\sqrt{\sum_{n\geq 1} \|f_n\|_{L^2(\mathcal{M})}^2 \min(1,Knt)^2} \simeq \sqrt{\sum_{n\geq 1} \|f_n\|_{L^2(\mathcal{M})}^2 \min(1,nt)^2}.
\end{equation*}
Note the equivalence $\delta\simeq \Upsilon(\delta_g)$.
For $\eqref{soli}\Rightarrow\eqref{solii}$, the idea\footnote{The author is grateful to a referee for that argument.} is that the Slepian theorem shows that the following two Gaussian processes are simultaneously sample-continuous or simultaneously unbounded:
\begin{itemize}
\item $(f^{G,\omega}(x))_{x\in \mathcal{M}_s}$ on a curve $\mathcal{M}_s$ of a compact manifold $ \mathcal{M}$ and $f^{G,\omega}(x)$ defined in \eqref{ran-fouG2},
\item $(f^{G,\omega}(\gamma(\theta))_{\theta\in [0,\sfrac{\pi}{2}]}$ with $\gamma:[0,\sfrac{\pi}{2}]\to \Gamma^+$ as in \eqref{defSS} but $(f^{G,\omega}(x))_{x\in \bbS^d}$ defined on the sphere $\bbS^d$ (as in \eqref{ran-fouG} with $\|f_n\|_{L^2(\mathcal{M})}$ instead of $\|f_n\|_{L^2(\bbS^d)}$) with the following form:
\begin{equation*}
f^{G,\omega}(x):=\sum_{n\geq 1} \frac{\|f_n\|_{L^2(\mathcal{M})}}{\sqrt{\dim(E_n)}}\sum_{\lambda_k^2=n(n+d-1)} g_{n,k}(\omega) \phi_k(x),\quad x\in \bbS^d.
\end{equation*}
\end{itemize}
Repeating the end of argument of Step 1 of Proposition \ref{ieii} or applying Proposition~\ref{ieii} itself, we get the Salem-Zygmund condition.

For $\eqref{solii}\Rightarrow\eqref{soli}$, the idea is that ii) is equivalent to the condition
\[
\int_{0}^{1} \frac{\Upsilon(t)}{t\sqrt{\ln(1/t)}}\,dt<+\infty,
\]
which in turn is equivalent to the entropy condition on the whole manifold $\mathcal{M}$ thanks to Proposition \ref{lem-ref} (note that the Riemannian probability volume $\mu$ of $\mathcal{M}$ satisfies \eqref{ln-hom}). We then get \eqref{soli} thanks to the Dudley theorem.

\item
For the proof of the analogue of Proposition \ref{ieii2}, the argument is much simpler since Theorem \ref{pseudo-dis} gives an equivalence holding true on the whole manifold $\mathcal{M}$.

\item
Finally, the quantitative counterpart given by Theorem \ref{sol-q} can be proved completely similarly once we use the inequalities given by the Slepian theorem (see \cite[Cor.\,3.14]{ledoux} or \cite[p.\,12]{queff2}).
\end{itemize}

\section{Proof of \texorpdfstring{\eqref{sol-u-eq} in Theorem \ref{sol-u}}{Pfth}, quantitative version}\label{sec-marc}

In the following proof, the implicit constants merely depend on $\mathcal{M},\mathcal{M}_s$ and $K$ (except for the setting $\mathcal{M}=\bbS^d$ and $f_n\in \ker(\Deltag+n(n+d-1))$ for which $K$ is obviously irrelevant). Moreover, in the manifold framework, the spectral parameter~$K$ is assumed to be large enough so that the conclusions of Theorem \ref{sol} and Theorem~\ref{sol-q} hold.

The quantitative part of Theorem \ref{sol-u}, namely the equivalence \eqref{sol-u-eq}, is a direct consequence of Proposition \ref{mapi}, Proposition \ref{mapi2} below and the Cauchy-Schwarz inequality:%
\begin{equation}\label{mp-cs}
\E_{\omega}\biggl[\sup_{x\in \mathcal{M}_s}\Big| \sum_{n=1}^N X_n(\omega) f_n^{\omega}(x) \Big| \biggr]\leq \E_{\omega}\biggl[\sup_{x\in \mathcal{M}_s}\Big| \sum_{n=1}^N X_n(\omega) f_n^{\omega}(x) \Big|^2 \biggr]^{\sfrac{1}{2}}.
\end{equation}

\begin{prop}\label{mapi}
There is a constant $C>0$ such that for any $N\in \N^\star$, any sequence~$(f_n)$ as in Theorem \ref{sol-u} and random variables $(X_n)$ belonging to $L^2(\Omega)$, the following holds true
\begin{multline}\label{mapi-b}
\E_{\omega}\biggl[\sup_{x\in \mathcal{M}_s}\Big| \sum_{n=1}^N X_n(\omega) f_n^{\omega}(x) \Big|^2 \biggr]^{\sfrac{1}{2}}\\ \leq C \sqrt{ \sup_{1\leq n\leq N} \E[|X_n|^2]}
\times \sum_{p=1}^N \frac{1}{p\sqrt{\ln(p+1)}}\,\biggl(\sum_{n=p}^{N} \|f_n\|_{L^2(\mathcal{M})}^2 \biggr)^{\sfrac{1}{2}},
\end{multline}
where all the random variables $X_1(\omega),f_1^\omega,\dots,X_N(\omega),f_N^\omega$ are assumed to be independent.
\end{prop}

The last result is analogue to \cite[p.\,53, Lem.\,1.1]{pisier1981} or \cite[p.\,249, Th.\,III.6]{queff2} but the proof of Proposition \ref{mapi} is somehow simpler and relies, among other arguments, on fact that the Salem-Zygmund condition \eqref{pz-con3} is very explicit.
We begin by recalling a form of the \textit{contraction principle} that we will use for the symmetric random functions $F_n:\omega\mto f_n^\omega$ (see the definitions in \eqref{ran-fou4} and \eqref{ran-fouu}).
\begin{prop}\label{p-con-1}
Let $B$ be a Banach space and $F_1,\dots,F_N$ be $N$ independent and symmetric random variables belonging to $L^p(\Omega,B)$ for some $p\in [1,+\infty)$. Then for any tuple $(a_1,\dots,a_N)\in \R^N$, the following inequalities hold true
\begin{equation*}
\min_{1\leq n \leq N}|a_n|\cdot \E\biggl[\Bigl\| \sum_{n=1}^N F_n \Bigr\|_{B}^p \biggr]^{\sfrac{1}{p}}\!\!\leq \E\biggl[\Bigl\| \sum_{n=1}^N a_n F_n \Bigr\|_{B}^p \biggr]^{\sfrac{1}{p}}\!\!\leq
\max_{1\leq n\leq N} |a_n|\cdot \E\biggl[\Bigl\| \sum_{n=1}^N F_n \Bigr\|_{B}^p \biggr]^{\sfrac{1}{p}}.
\end{equation*}
\end{prop}
\begin{proof}
The first inequality is a consequence of the second one (upon assuming that each $a_n$ is not zero and upon replacing $a_n$ with $\sfrac{1}{a_n}$).
For the second one, we refer to \cite[p.\,136]{queff1} or the proof of \cite[Th.\,4.4]{ledoux}.
\end{proof}

It is known that the previous result has a reformulation if each coefficient $a_n$ is random (see
\cite[Lem.\,4.5]{ledoux}, \cite[Th.\,4.9, p.\,45]{pisier1981} or \cite[p.\,137, Th.\,IV.4]{queff1}).
For the convenience of the reader, we write the proof of the following result (but all ideas are included in the last references).

\begin{coro}
Let $B$ be a Banach space, $F_1,\dots,F_N$ be $N$ symmetric random variables belonging to $L^p(\Omega,B)$ for some $p\in [1,+\infty)$ and $A_1,\dots,A_N$ be $N$ real random variables belonging to $L^p(\Omega,\R)$.
We assume that the $2N$ variables $A_1$,$F_1$,...,$A_N$,$F_N$ are mutually independent.
Then the following inequalities hold true
\begin{align}\label{p-con-maj}
\E\biggl[\Big\| \sum_{n=1}^N A_n F_n \Big\|_{B}^p \biggr]^{\sfrac{1}{p}}
& \leq \E\biggl[\max_{1\leq n \leq N} |A_n|^p\biggr]^{\sfrac{1}{p}}\cdot \E\biggl[\Big\| \sum_{n=1}^N F_n \Big\|_{B}^p \biggr]^{\sfrac{1}{p}},\\ \label{p-con-min}
\min_{1\leq n\leq N}\E[|A_n|]\cdot \E\biggl[\Big\| \sum_{n=1}^N F_n \Big\|_{B}^p \biggr]^{\sfrac{1}{p}}& \leq \E\biggl[\Big\| \sum_{n=1}^N A_n F_n \Big\|_{B}^p \biggr]^{\sfrac{1}{p}}.
\end{align}
\end{coro}

\begin{proof}
By independence (see \cite[App.\,F]{imek-aif} for more details), we may write
\begin{equation}\label{p-con-2a}
\E\biggl[\Big\| \sum_{n=1}^N A_n F_n \Big\|_{B}^p \biggr]=
\E_{\omega_1} \E_{\omega_2}\biggl[\Big\| \sum_{n=1}^N A_n(\omega_1) F_n(\omega_2) \Big\|_{B}^p \biggr].
\end{equation}
We now freeze $\omega_1$ and apply Proposition \ref{p-con-1} with respect to the expectation in $\omega_2$. Hence we get
\begin{equation*}
\E_{\omega_2}\biggl[\Big\| \sum_{n=1}^N A_n(\omega_1) F_n(\omega_2) \Big\|_{B}^p \biggr] \leq \max_{1\leq n\leq N} |A_n(\omega_1)|^p\times
\E_{\omega_2}\biggl[\Big\| \sum_{n=1}^N F_n(\omega_2) \Big\|_{B}^p \biggr].
\end{equation*}
We then obtain the expected inequality by integrating in $\omega_1$.

In order to get the lower bound, we first need to modify the right-hand side of~\eqref{p-con-2a}.
For $\omega_1$ fixed, let us consider $\theta_n(\omega_1)=\pm 1$ such that $A_n(\omega_1)\theta_n(\omega_1)=|A_n(\omega_1)|$.
By~\hbox{symmetry} and independence, one can replace each $F_n(\omega_2)$ with $\theta_n(\omega_1) F_n(\omega_2)$ in the expectation $\E_{\omega_2}$ in \eqref{p-con-2a}. So we get\vspace*{-3pt}
\begin{equation}\label{symsym}
\E\biggl[\Big\| \sum_{n=1}^N A_n F_n \Big\|_{B}^p \biggr]=
\E_{\omega_1} \E_{\omega_2}\biggl[\Big\| \sum_{n=1}^N |A_n(\omega_1)| F_n(\omega_2) \Big\|_{B}^p \biggr].
\end{equation}
We now bound from below with the Hölder inequality and the triangular inequality in $\omega_1$, namely $\E_{\omega_1}[\|\cdot\|]\geq \|\E_{\omega_1}[\hspace{1 mm}\cdot\hspace{1 mm}]\|$, as follows\vspace*{-3pt}
\begin{align*}
\E\biggl[\Big\| \sum_{n=1}^N A_n F_n \Big\|_{B}^p \biggr]&=
\E_{\omega_2}\biggl[\E_{\omega_1}\biggl[\Big\| \sum_{n=1}^N |A_n(\omega_1)| F_n(\omega_2) \Big\|_{B}^p \biggr]\biggr]\\[-3pt]
&\geq \E_{\omega_2}\biggl[\E_{\omega_1}\biggl[\Big\|\sum_{n=1}^N |A_n(\omega_1)|F_n(\omega_2) \Big\|_B \biggr]^p \biggr]\\[-3pt]
&\geq \E_{\omega_2}\biggl[\Big\|\sum_{n=1}^N \E_{\omega_1}\big[|A_n(\omega_1)|F_n(\omega_2) \big]\Big\|_B^p \biggr]=\E\biggl[\Big\|\sum_{n=1}^N \E[|A_n|]F_n \Big\|_B^p \biggr].
\end{align*}
The second inequality of the statement is then a consequence of Proposition \ref{p-con-1}.
\end{proof}

\skpt
\begin{proof}[Proof of Proposition \ref{mapi}]
\vspace*{-.5\baselineskip}
\subsubsection*{Step 1} We first need to compare $f_n^\omega$ with $f_n^{G,\omega}$.
More precisely, we shall prove
\begin{equation}\label{mapi-b1}
\E_{\omega}\biggl[\sup_{x\in \mathcal{M}_s}\Big|\sum_{n=1}^N X_n(\omega)f_n^{\omega}(x)\Big|^{2} \biggr]\leq C\E_{\omega'}\E_{\omega}\biggl[\sup_{x\in \mathcal{M}_s}\Big| \sum_{n=1}^N |X_n(\omega')| f_n^{G,\omega}(x) \Big|^2 \biggr].
\end{equation}
The distribution equivalence \eqref{chi2a} for the manifold framework suggests introducing the following random variable (collinear to a so-called chi random variable)
\begin{equation}\label{chi}
\forall \omega\in \Omega,\quad \chi_n(\omega):=\frac{1}{\sqrt{\dim(E_{(Kn-K,Kn]})}}\biggl(\sum_{Kn-K<\lambda_k\leq Kn} g_{n,k}^2(\omega)\biggr)^{\sfrac{1}{2}}.\end{equation}
The case $\mathcal{M}=\bbS^d$ with $f_n\in \ker(\Deltag+n(n+d-1))$ is completely similar and merely needs to change the definition \eqref{chi}.
Having assumed the independence of all the random variables, we may write\vspace*{-3pt}
\begin{multline}\label{chi-g}
\E_{\omega'}\E_{\omega}\biggl[\sup_{x\in \mathcal{M}_s}\Bigl| \sum_{n=1}^N |X_n(\omega')| f_n^{G,\omega}(x) \Bigr|^2 \biggr]\\[-10pt]
=\E_{\omega'}\E_\omega\biggl[\sup_{x\in \mathcal{M}_s}\Bigl|\sum_{n=1}^N |X_n(\omega')| \chi_n(\omega) f_n^{\omega}(x)\Bigr|^2 \biggr].
\end{multline}
For almost any $\omega'\in \Omega$, one may see $|X_n(\omega')|$ as a constant in the expectation $\E_{\omega}$.
Hence, the contraction principle given by \eqref{p-con-min} in $L^2(\Omega,\CC^0(\mathcal{M}_s))$, with $A_n(\omega)=\chi_n(\omega)$ and $F_n(\omega)=|X_n(\omega')|f_n^{\omega}(x)$, gives the bound from below\vspace*{-3pt}
\begin{multline*}
\E_{\omega'}\E_\omega\biggl[\sup_{x\in \mathcal{M}_s}\Big|\sum_{n=1}^N |X_n(\omega')| \chi_n(\omega) f_n^{\omega}(x)\Big|^2 \biggr]\\[-8pt]
\geq \Bigl(\inf_{1\leq n\leq N} \E[\chi_n]^2 \Bigr)
\E_{\omega'}\E_\omega\biggl[\sup_{x\in \mathcal{M}_s}\Big|\sum_{n=1}^N |X_n(\omega')| f_n^{\omega}(x)\Big|^2 \biggr].
\end{multline*}
By independence and symmetry of each $\omega\mto f_n^\omega$ (as in \eqref{symsym}), the last lower bound equals\vspace*{-5pt}
\begin{equation*}
\Big(\inf_{1\leq n\leq N} \E[\chi_n]^2 \Big)
\E_\omega\biggl[\sup_{x\in \mathcal{M}_s}\Big|\sum_{n=1}^N X_n(\omega) f_n^{\omega}(x)\Big|^2 \biggr].
\end{equation*}
Now we use the inequality $\E[\chi_n]\geq \sfrac{\sqrt{2}}{\sqrt{\pi}}$ (see Appendix \ref{chi2}) to get \eqref{mapi-b1}.

\subsubsection*{Step 2}
For $\omega'$ fixed, we now recall that $\sum_{n=1}^N |X_n(\omega')| f_n^{G,\omega}$ can be seen, with respect to $\omega$, as a Gaussian random finite sum in the Banach space $\CC^0(\mathcal{M}_s)$.
We now invoke the important property stating that all the moments of such a Gaussian sum are universally equivalent, that is the Gaussian version of the Kahane-Khintchine inequalities (see \cite[Cor.\,3.2]{ledoux}, \cite[p.\,44]{pisier1981} or \cite[p.\,256, Cor.\,V.27]{queff1}).
Hence, \eqref{mapi-b1} implies the following bound
\begin{equation*}
\E_{\omega}\biggl[\sup_{x\in \mathcal{M}_s}\Big|\sum_{n=1}^N X_n(\omega)f_n^{\omega}(x)\Big|^{2} \biggr]\leq C\E_{\omega'}\biggl[\E_{\omega}\biggl[\sup_{x\in \mathcal{M}_s}\Big| \sum_{n=1}^N |X_n(\omega')| f_n^{G,\omega}(x) \Big| \biggr]^2 \biggr].
\end{equation*}
Note now that changing $f_n$ by $X_n(\omega')f_n$ in \eqref{ran-fouG2} leads to change $f_n^{G,\omega}$ by $|X_n(\omega')| f_n^{G,\omega}$.
Consequently, Theorem \ref{sol-q} (for $f_n$ replaced with $X_n(\omega')f_n$) allows us to replace the last upper bound with
\begin{equation}\label{mapi-b3}
C\E_{\omega'}\biggl[\biggl(\sum_{p=1}^N \frac{1}{p\sqrt{\ln(p+1)}}\Big(\sum_{n=p}^{N} X_n(\omega')^2\|f_n\|_{L^2(\mathcal{M})}^2 \Big)^{\sfrac{1}{2}}\biggr)^2\biggr].
\end{equation}
Step 3 below will show the concavity of the following function:
\begin{equation}\label{concav}
\begin{split}
\Psi: [0,+\infty)^N & \to\R \\
(t_1,\dots,t_N) & \mto \biggl(\sum_{p=1}^N \frac{1}{p\sqrt{\ln(p+1)}}\Bigl(\sum_{n=p}^{N} t_n \|f_n\|_{L^2(\mathcal{M})}^2 \Bigr)^{\sfrac{1}{2}}\biggr)^2.
\end{split}
\end{equation}
Note that the number in \eqref{mapi-b3} equals $C \E[\Psi(X_1^2,\dots,X_N^2)]$.
Hence, the multidimensional Jensen inequality allows to bound it by
\begin{align*}
C \Psi(\E[X_1^2],\dots,&\E[X_N^2])
=C\biggl(\sum_{p=1}^N \frac{1}{p\sqrt{\ln(p+1)}}\biggl(\sum_{n=p}^{N} \E[| X_n|^2]\|f_n\|_{L^2(\mathcal{M})}^2 \biggr)^{\sfrac{1}{2}}\biggr)^2\\
&\leq C \Bigl(\sup_{1\leq n\leq N} \E[|X_n|^2] \Bigr)
\biggl(\sum_{p=1}^N \frac{1}{p\sqrt{\ln(p+1)}}\biggl(\sum_{n=p}^{N} \|f_n\|_{L^2(\mathcal{M})}^2 \biggr)^{\sfrac{1}{2}}\biggr)^2.
\end{align*}
The expected inequality \eqref{mapi-b} is proved.

\subsubsection*{Step 3} It remains to check the concavity of \eqref{concav}. By developing the square in \eqref{concav}, we see that is is sufficient to prove the concavity of any continuous function of the form $\sqrt{\Phi_1 \Phi_2}$ on the set $\Lambda:=\{\Phi_1\geq 0\}\cap\{\Phi_2\geq 0\}$ in which $\Phi_1:\R^N\to \R$ and $\Phi_2:\R^N\to \R$ are two linear functionals. Actually, such a fact is a straightforward consequence of the Cauchy-Schwarz inequality applied to the
two vectors $(\sqrt{\Phi_1(s)},\sqrt{\Phi_1(t)})$ and $(\sqrt{\Phi_2(s)},\sqrt{\Phi_2(t)})$
for any $(s,t)\in \Lambda^2$:
\begin{align*}
\sqrt{\Phi_1\bigl(\psfrac{s+t}{2}\bigr)\Phi_2\bigl(\psfrac{s+t}{2}\bigr)}&=\frac{\sqrt{\Phi_1(s)+\Phi_1(t)}\sqrt{\Phi_2(s)+\Phi_2(t)}}{2} \\
& \geq \frac{\sqrt{\Phi_1(s)}\sqrt{\Phi_2(s)}+\sqrt{\Phi_1(t)}\sqrt{\Phi_2(t)}}{2}.\qedhere
\end{align*}
\end{proof}

In in order to finish the proof of the numerical equivalence \eqref{sol-u-eq}, we have to reverse the inequality \eqref{mapi-b}.
Actually, we will use a suitable truncation argument first used by Marcus and Pisier (see \cite[p.\, 53-54, proof of Lem.\,1.1 \& p.\,99, Lem.\,3.7]{pisier1981}).
In~our specific context, the truncation argument is based on the uniform estimate \eqref{chi-M} of the sequence of random variables $(\chi_n)_{n\geq 1}$ defined in \eqref{chi}.

\begin{prop}\label{mapi2}
There is a constant $C>0$ such that for any $N\in \N^\star$ and any sequence $(f_n)$ and random variables $(X_n)$ as in Theorem \ref{sol-u}, the following holds true\vspace*{-.5\baselineskip}
\begin{multline*}
\E_{\omega}\biggl[\sup_{x\in \mathcal{M}_s}\Big| \sum_{n=1}^N X_n(\omega) f_n^{\omega}(x) \Big| \biggr]\\[-5pt]
\geq C\Big(\inf_{1\leq n\leq N} \E[|X_n|]\Big)\times \sum_{p=1}^N \frac{1}{p\sqrt{\ln(p+1)}}\biggl(\sum_{n=p}^{N} \|f_n\|_{L^2(\mathcal{M})}^2 \biggr)^{\sfrac{1}{2}}.
\end{multline*}
\end{prop}
\begin{proof}

The contraction principle \eqref{p-con-min} and Theorem \ref{sol-q} ensure that we merely have to prove the following inequality:\vspace*{-3pt}
\begin{equation}\label{mapi2-b2}
\E_{\omega}\biggl[\sup_{x\in \mathcal{M}_s} \Big| \sum_{n=1}^N f_n^{G,\omega}(x)\Big| \biggr]\leq C\E_{\omega}\biggl[\sup_{x\in \mathcal{M}_s} \Big| \sum_{n=1}^N f_n^{\omega}(x)\Big| \biggr].
\end{equation}
As in \eqref{chi-g}, we have\vspace*{-3pt}
\begin{equation*}
\E_{\omega}\biggl[\sup_{x\in \mathcal{M}_s} \Big| \sum_{n=1}^N f_n^{G,\omega}(x)\Big| \biggr]=\E_{\omega}\biggl[\sup_{x\in \mathcal{M}_s} \Big| \sum_{n=1}^N \chi_n(\omega)f_n^{\omega}(x)\Big| \biggr],
\end{equation*}
where all the random variables are assumed to be mutually independent.
Let us now consider a number $M>0$ (that will be chosen below) and we bound\vspace*{-.5\baselineskip}
\begin{multline*}
\E_{\omega}\biggl[\sup_{x\in \mathcal{M}_s} \Big| \sum_{n=1}^N f_n^{G,\omega}(x)\Big| \biggr]\leq \E_{\omega}\biggl[\sup_{x\in \mathcal{M}_s} \Big| \sum_{n=1}^N \chi_n(\omega)\pun_{\chi_n(\omega)> M}f_n^{\omega}(x)\Big| \biggr] \\[-5pt]
+ \E_{\omega}\biggl[\sup_{x\in \mathcal{M}_s} \Big| \sum_{n=1}^N \chi_n(\omega)\pun_{\chi_n(\omega)\leq M}f_n^{\omega}(x)\Big| \biggr].
\end{multline*}
By using \eqref{mp-cs} and \eqref{mapi-b} with $X_n=\chi_n \pun_{\chi_n>M}$ and Theorem \ref{sol-q} on the one hand, and then the contraction principle \eqref{p-con-maj} on the other hand, we obtain\vspace*{-3pt}
\begin{multline*}
\E_{\omega}\biggl[\sup_{x\in \mathcal{M}_s} \Big| \sum_{n=1}^N f_n^{G,\omega}(x)\Big| \biggr]\\[-5pt]
\leq C\Big(\sup_{n\in \N^\star} \sqrt{\E[|\chi_n \pun_{\chi_n> M}|^2]} \Big)\E_{\omega}\biggl[\sup_{x\in \mathcal{M}_s} \Big| \sum_{n=1}^N f_n^{G,\omega}(x)\Big| \biggr]
+ M\E_{\omega}\biggl[\sup_{x\in \mathcal{M}_s} \Big| \sum_{n=1}^N f_n^{\omega}(x)\Big| \biggr].
\end{multline*}
In order to get \eqref{mapi2-b2}, it is sufficient to justify that there is
$M>0$ satisfying
\begin{equation*}\sup_{n\in \N^\star} \sqrt{\E[|\chi_n \pun_{\chi_n> M}|^2]} <\frac{1}{C}.
\end{equation*}
That is actually a consequence of the following more precise bound (See Appendix~\ref{chi2}):%
\begin{equation}\label{chi-M}
\E[|\chi_n \pun_{\chi_n> M}|^2]=\E[\chi_n^2 \pun_{\chi_n> M}] \leq \E[\sfrac{\chi_n^4}{M^2}]\leq \frac{3}{M^2}.\qedhere
\end{equation}
\end{proof}

\section{Proof of Theorem \ref{sol-u}, qualitative version}\label{sec-marc2}

The end of the proof of Theorem \ref{sol-u} is a consequence of the equivalence $\eqref{plu-eqi}\Leftrightarrow\eqref{plu-eqiv}$ of the following result.

\begin{prop}\label{plu-eq}
Let us assume the same assumptions as in Theorem \ref{sol-u}. Then the following assertions are equivalent:
\begin{enumeratei}
\item\label{plu-eqi} the random series $\sum X_n(\omega) f_n^{\omega}$ almost surely converges in $\CC^0(\mathcal{M}_s)$,
\item\label{plu-eqii} with probability $1$, the sequence of partial sums of the random series $\sum X_n(\omega) f_n^{\omega}$ is bounded in $\CC^0(\mathcal{M}_s)$,
\item\label{plu-eqiii} the sequence of partial sums of the random series $\sum X_n(\omega) f_n^{\omega}$ is bounded in $L^1(\Omega,\CC^0(\mathcal{M}_s))$,
\item\label{plu-eqiv} the Salem-Zygmund condition
\[
\sum_{p=1}^{+\infty} \frac{1}{p\sqrt{\ln(p+1)}}\biggl(\sum_{n=p}^{+\infty} \left\Vert f_n\right\Vert_{L^2(\mathcal{M})}^2 \biggr)^{\sfrac{1}{2}}<+\infty
\]
holds,
\item\label{plu-eqv} the random series $\sum X_n(\omega) f_n^\omega$ converges in $L^2(\Omega,\CC^0(\mathcal{M}_s))$.
\end{enumeratei}
\end{prop}

\begin{rema} The Gaussian random series $\sum f_n^{G,\omega}$ can be written as $\sum \chi_n(\omega) f_n^{\omega}$ (see \eqref{chi-g}) in which the sequence $(X_n)=(\chi_n)$ satisfies the assumptions \eqref{hyp-Xn} (see Appendix \ref{chi2}). In other words, Proposition \ref{plu-eq} holds true for the Gaussian random series $\sum f_n^{G,\omega}$.
\end{rema}

Let us prove Proposition \ref{plu-eq}.

\subsubsection*{$\eqref{plu-eqi}\Rightarrow\eqref{plu-eqii}$}
Obvious.

\subsubsection*{$\eqref{plu-eqii}\Rightarrow\eqref{plu-eqiii}$}
We shall use the Paley-Zygmund inequality in a similar spirit to
\cite[p.\,55, Lem.\,1.2]{pisier1981} and \cite[Prop.\,2.17]{imek-crusep}. For any $N\in \N^\star$ and $\omega\in \Omega$, we set the partial sum
\begin{equation*}
S_N(\omega):=\sum_{n=1}^N X_n(\omega) f_n^{\omega}.
\end{equation*}
We now give a proof by contradiction and we may assume that there is a subsequence of integers $(N_k)_{k\geq 1}$ satisfying
\begin{equation}\label{contradi}
\E[\|S_{N_k}\|_{\CC^0(\mathcal{M})_s}]\geq 2k.
\end{equation}
The classical Paley-Zygmund inequality \cite[p.\,8, Ineq.\,II]{kahane} can be combined to Proposition \ref{mapi} and Proposition \ref{mapi2} so we get
\begin{equation*}
\P\Big(\|S_{N_k}\|_{\CC^0(\mathcal{M}_s)} \geq \frac{1}{2} \E[\|S_{N_k}\|_{\CC^0(\mathcal{M}_s)}]\Big)\geq \frac{\E[\|S_{N_k}\|_{\CC^0(\mathcal{M}_s)}]^2}{4
\E[\|S_{N_k}\|_{\CC^0(\mathcal{M}_s)}^2]} \geq C\,\frac{\inf_{1\leq n\leq N_k}\E[|X_n|]^2}{\sup_{1\leq n \leq N_k} \E[|X_n|^2] }.
\end{equation*}
Therefore \eqref{hyp-Xn} and \eqref{contradi} imply
\begin{equation*}
\inf_{k\geq 1} \P\bigl(\|S_{N_k}\|_{\CC^0(\mathcal{M}_s)} \geq k]\bigr)>0.
\end{equation*}
That inequality contradicts ii) since the dominated convergence theorem proves
\begin{equation*}
\lim_{k\to +\infty} \P(\|S_{N_k} \|_{\CC^0(\mathcal{M}_s)}\geq k)=\lim_{k\to +\infty} \E_{\omega}\Big[\pun_{\{\|S_{N_k}(\omega) \|_{\CC^0(\mathcal{M}_s)}\geq k\}}\Big]=0.
\end{equation*}

\subsubsection*{$\eqref{plu-eqiii}\Rightarrow\eqref{plu-eqiv}$}
Thanks to Proposition \ref{mapi2} and the assumption \eqref{hyp-Xn}, we have
\begin{equation*}
\sup_{N\in \N^\star} \sum_{p=1}^N \frac{1}{p\sqrt{\ln(p+1)}} \biggl(\sum_{n=p}^{N} \|f_n\|_{L^2(\mathcal{M})}^2 \biggr)^{\sfrac{1}{2}}<+\infty,
\end{equation*}
which implies
\begin{equation*}
\sup_{\substack{(N,N')\in (\N^\star)^2 \\ N<N' }} \sum_{p=1}^{N} \frac{1}{p\sqrt{\ln(p+1)}} \biggl(\sum_{n=p}^{N'} \|f_n\|_{L^2(\mathcal{M})}^2 \biggr)^{\sfrac{1}{2}} <+\infty,
\end{equation*}
and which in turn means the convergence of
\[
\sum_{p=1}^{+\infty} \frac{1}{p\sqrt{\ln(p+1)}} \biggl(\sum_{n=p}^{+\infty} \|f_n\|_{L^2(\mathcal{M})}^2 \biggr)^{\sfrac{1}{2}}.
\]

\subsubsection*{$\eqref{plu-eqiv}\Rightarrow\eqref{plu-eqv}$}
Let us prove that $(S_N)_{N\geq 1}$ is a Cauchy sequence in $L^2(\Omega,\CC^0(\mathcal{M}_s))$.
As a consequence of Proposition \ref{mapi}, for any integers $N'>N$ we have
\begin{multline*}
\E \Bigl[\|S_{N'}-S_N\|_{\CC^0(\mathcal{M}_s)}^2 \Bigr]^{\sfrac{1}{2}} \\
\leq
C \sqrt{\sup_{N< n\leq N'}\E[|X_n|^2]} \times \sum_{p=1}^{N'} \frac{1}{p\sqrt{\ln(p+1)}} \biggl(\sum_{n=p}^{N'} \|f_n\|_{L^2(\mathcal{M})}^2 \pun_{[N+1,+\infty)}(n)\biggr)^{\sfrac{1}{2}}.
\end{multline*}
Looking at $p=1$ in the Salem-Zygmund condition, we see that \eqref{plu-eqiv} implies the convergence of the series $\sum \|f_n\|_{L^2(\mathcal{M})}^2$.
We then conclude with the assumption \eqref{hyp-Xn} and the following limit
proved via the discrete dominated convergence theorem:
\begin{equation*}
\lim_{N\to +\infty}\sum_{p=1}^{+\infty} \frac{1}{p\sqrt{\ln(p+1)}} \biggl(\sum_{n=p}^{+\infty} \|f_n\|_{L^2(\mathcal{M})}^2 \pun_{[N+1,+\infty)}(n) \biggr)^{\sfrac{1}{2}}=0.
\end{equation*}

\subsubsection*{$\eqref{plu-eqv}\Rightarrow\eqref{plu-eqi}$}
Due to the Markov inequality, we get the convergence in probability in $\CC^0(\mathcal{M}_s)$. Then \eqref{plu-eqi} is a known consequence (see \cite[Th.\,6.1]{ledoux} or \cite[p.\,130--131, Th.\,III.3]{queff1}).

\section{Proof of Theorem \ref{indep-met}, semi-classical analysis and independence of the\nobreakspace Riemannian metric}\label{s-wrt}

The Salem-Zygmund condition \eqref{pz-con30} depends on the Riemannian structure via the spectral analysis of the Laplace-Beltrami operator $\Deltag$.
Hence, we have to prove that if one considers another Laplace-Beltrami operator $\Deltag$ on $\mathcal{M}$ (defined with another Riemannian structure of $\mathcal{M}$), then the Salem-Zygmund condition does not change.

For any $K\geq K_0$ we set $f=\sum_{n\geq 1} f_n\in L^2(\mathcal{M})$ with $f_n\in E_{(Kn-K,Kn]}$. Noting that $f_1$ does not impact the convergence of the Salem-Zygmund condition \eqref{pz-con30} and remembering that the subspaces $E_{(Kn-K,Kn]}$ are orthogonal, we see that \eqref{pz-con30} is indeed equivalent to
\begin{equation}\label{pz-con4}
\sum_{p=1}^{+\infty} \frac{\|\Pi_{(Kp,+\infty)}(f)\|_{L^2(\mathcal{M})}}{p\sqrt{\ln(p+1)}}<+\infty,
\end{equation}
where $\Pi_{(Kp,+\infty)}:L^2(\mathcal{M})\to L^2(\mathcal{M})$ is the spectral projector on the spectral window $(Kp,+\infty)$ with respect to $\sqrt{-\Deltag}$.
Although the inequality $K\geq K_0$ is essential in the proof of Theorem \ref{sol-u} (see Theorem \ref{pseudo-dis}), it turns out that \eqref{pz-con4} is independent of $K>0$ from a Hilbertian point of view. More precisely, the following lemma shows the first part of the statement of Theorem \ref{indep-met}.
\begin{lemm}\label{wrtK}
For any $f\in L^2(\mathcal{M})$, the condition \eqref{pz-con4} is independent of $K\in (0,+\infty)$.
\end{lemm}
\begin{proof}
Let $K'$ and $K$ be two spectral parameters satisfying $K\leq K'$. Let us moreover consider a positive integer $q$ satisfying $K\leq K'\leq Kq$.
Thus we clearly have
\begin{equation*}
\sum_{p=1}^{+\infty} \frac{\|\Pi_{(Kp,+\infty)}(f)\|_{L^2(\mathcal{M})}}{p\sqrt{\ln(p+1)}}\geq \sum_{p=1}^{+\infty} \frac{\|\Pi_{(K'p,+\infty)}(f)\|_{L^2(\mathcal{M})}}{p\sqrt{\ln(p+1)}}.
\end{equation*}
In order to reverse that condition, we just note that for any $r\in \{0,1,\dots,q-1\}$ we~have $K'p\leq Kpq+Kr$ and hence for $p\geq 1$
\begin{equation*}
\frac{\|\Pi_{(K'p,+\infty)}(f)\|_{L^2(\mathcal{M})}}{p\sqrt{\ln(p+1)}} \geq \frac{\|\Pi_{(Kpq+Kr,+\infty)}(f)\|_{L^2(\mathcal{M})}}{p\sqrt{\ln(p+1)}} \geq \frac{\|\Pi_{(Kpq+Kr,+\infty)}(f)\|_{L^2(\mathcal{M})}}{(pq+r)\sqrt{\ln(pq+r+1)}} .
\end{equation*}
By summing over $r$ and $p$, we get
\begin{multline*}
\sum_{p=1}^{+\infty} \frac{\|\Pi_{(K'p,+\infty)}(f)\|_{L^2(\mathcal{M})}}{p\sqrt{\ln(p+1)}} =\frac{1}{q}\sum_{r=0}^{q-1} \sum_{p=1}^{+\infty} \frac{\|\Pi_{(K'p,+\infty)}(f)\|_{L^2(\mathcal{M})}}{p\sqrt{\ln(p+1)}}\\
\geq \frac{1}{q}\sum_{r=0}^{q-1} \sum_{p=1}^{+\infty}\frac{\|\Pi_{(Kpq+Kr,+\infty)}(f)\|_{L^2(\mathcal{M})}}{(pq+r)\sqrt{\ln(pq+r+1)}}=\frac{1}{q} \sum_{p=q}^{+\infty} \frac{\|\Pi_{(Kp,+\infty)}(f)\|_{L^2(\mathcal{M})}}{p\sqrt{\ln(p+1)}}.\qedhere
\end{multline*}
\end{proof}

The sequel is devoted to the proof of the invariance of \eqref{pz-con4} with respect to the choice of the Laplace-Beltrami operator $\Deltag$.
We first need a semi-classical reformulation of the Salem-Zygmund condition.

\begin{lemm}\label{wrt1}
Let us fix a smooth function $\Psi:[0,+\infty)\to [0,1]$ satisfying
\begin{equation}\label{est-sig}
\begin{split}
t \in [0,1] & \implies \Psi(t)=1, \\
t \geq 2 & \implies \Psi(t)=0.
\end{split}
\end{equation}
Then, for any $f\in L^2(\mathcal{M})$, the Salem-Zygmund condition \eqref{pz-con4} is equivalent to the following semi-classical condition
\begin{equation}\label{pz-con-sc}
\int_{0}^1 \frac{\|f-\Psi(-h^2 \Deltag)f \|_{L^2(\mathcal{M})}}{h\sqrt{-\ln(h)}} \,dh<+\infty.
\end{equation}
\end{lemm}
\begin{proof}
Following \eqref{hilb-bas}, we decompose $f=\sum_{k\in \N} c_k \phi_k$ with $(c_k)\in \ell^{2}(\N)$. Hence we~get
\begin{align*}
f-\Psi(-h^2 \Deltag)f& = \sum_{k=0}^{+\infty} (1-\Psi(h^2 \lambda_k^2)) c_k\phi_k, \\
\|f-\Psi(-h^2 \Deltag)f\|_{L^2(\mathcal{M})}^2 &= \sum_{k=0}^{+\infty}
(1-\Psi(h^2 \lambda_k^2))^2 |c_k|^2.
\end{align*}
The imposed conditions on $\Psi$ allow us to bound as follows
\begin{align*}
\sum_{\lambda_k >\sfrac{\sqrt{2}}{h}} |c_k|^2& \leq \|f-\Psi(-h^2 \Deltag)f\|_{L^2(\mathcal{M})}^2 \leq \sum_{\lambda_k >\sfrac{1}{h}} |c_k|^2,\\
\|\Pi_{(\sfrac{\sqrt{2}}{h},+\infty)}(f)\|_{L^2(\mathcal{M})}& \leq \|f-\Psi(-h^2 \Deltag)f\|_{L^2(\mathcal{M})} \leq \|\Pi_{(\sfrac{1}{h},+\infty)}(f)\|_{L^2(\mathcal{M})}.
\end{align*}
By integrating with respect to $h\in (0,1]$, we get
\begin{align}\label{wrtK2}
\int_{0}^1 \frac{\|\Pi_{(\sfrac{\sqrt{2}}{h},+\infty)}(f)\|_{L^2(\mathcal{M})} }{h\sqrt{-\ln(h)}}\,dh & \leq 
\int_{0}^1 \frac{\|f-\Psi(-h^2 \Deltag)f\|_{L^2(\mathcal{M})} }{h\sqrt{-\ln(h)}}\,dh \\ \label{wrtK2b}
& \leq
\int_{0}^1 \frac{\|\Pi_{(\sfrac{1}{h},+\infty)}(f)\|_{L^2(\mathcal{M})} }{h\sqrt{-\ln(h)}}\,dh.
\end{align}
By using that $h\in (0,1]\mto \|\Pi_{(\sfrac{1}{h},+\infty)}(f)\|_{L^2(\mathcal{M})}$ is non-decreasing, we see that the upper bound in \eqref{wrtK2b} is bounded from above by
\begin{align*}
\sum_{p\geq 1} \int_{\spfrac{1}{p+1}}^{\sfrac{1}{p}}\frac{\|\Pi_{(\sfrac{1}{h},+\infty)}(f)\|_{L^2(\mathcal{M})} }{h\sqrt{-\ln(h)}}\,dh &\leq\sum_{p\geq 1} \|\Pi_{(p,+\infty)}(f)\|_{L^2(\mathcal{M})} \int_{\spfrac{1}{p+1}}^{\sfrac{1}{p}} \frac{dh}{h\sqrt{-\ln(h)}} \\
& \lesssim \sum_{p\geq 1} \frac{\|\Pi_{(p,+\infty)}(f)\|_{L^2(\mathcal{M})}}{p\sqrt{\ln(p+1)}}.
\end{align*}
Similarly, we see that the lower bound in \eqref{wrtK2} is bounded from below by
\[
\sum_{p\geq 1}\int_{\spfrac{1}{p+1}}^{\sfrac{1}{p}}\frac{\|\Pi_{(\sfrac{\sqrt{2}}{h},+\infty)}(f)\|_{L^2(\mathcal{M})} }{h\sqrt{-\ln(h)}}\,dh \gtrsim \sum_{p\geq 2} \frac{\|\Pi_{(\sqrt{2}p,+\infty)}(f)\|_{L^2(\mathcal{M})} }{p\sqrt{\ln(p+1)}}.
\]
We conclude with Lemma \ref{wrtK}.
\end{proof}

Before analyzing the semi-classical Salem-Zygmund condition \eqref{pz-con-sc}, we need to recall a few facts about semi-classical symbols and semi-classical pseudo-differential operators.
The standard relation between a symbol $s:\R^d\times \R^d\times (0,1]\to \R$ and its quantized operator, usually denoted by $s(x,D,h)$, is given by the formula
\begin{equation}\label{defi-opd}
\forall F\in \mathcal{S}(\R^d),\ \forall x\in \R^d, \quad (s(x,D,h)F)(x)=\int_{\R^d}e^{i\langle x,\xi\rangle}
s(x,\xi,h) \widehat{F}(\xi)\, \frac{d\xi}{(2\pi)^d}.
\end{equation}
For the particular case $s(x,h\xi)$, one prefers writing $s(x,hD)$.
Such operators, with $h$ running over $(0,1]$, are called semi-classical pseudo-differential operators. Here is the result we need in the sequel.

\begin{lemm}\label{lem-sc}
For any symbols $s_1$ and $s_2$ belonging to $\CC_b^\infty(\R^d\times \R^d)$ and having disjoint supports, for any $N\in \N^\star$, there are constant $C>0$ and $C_N$ such that the following estimates uniformly hold for any $h\in (0,1]$:
\begin{align}\label{ca1}
\| s_1(x,hD)\|_{L^2(\R^d)\to L^2(\R^d)} &\leq C, \\ \label{ca2}
\| s_1(x,hD)\circ s_2(x,hD)\|_{L^2(\R^d)\to L^2(\R^d)} & \leq C_N h^N.
\end{align}
\end{lemm}
\begin{proof}
We will use the notations of \cite[p.\,22--23]{lerner-pseudo}. For any $m\in \R$, one denotes by $S_{\operatorname{scl}}^m$ the space of symbols $s(x,\xi,h)$ satisfying for any $\alpha \in \N^d$ and $\beta \in \N^d$:
\begin{equation*}
\sup_{x\in \R^d} \sup_{\xi\in \R^d}\sup_{h\in (0,1]}|(\partial_\xi^\alpha \partial_x^\beta s)(x,\xi,h)| h^{m-|\alpha|}<+\infty.
\end{equation*}
With those notations, the symbols $s_1(x,h\xi)$ and $s_2(x,h\xi)$ belong to the class $S_{\operatorname{scl}}^0$.
The first estimate is merely the semi-classical Calderon-Vaillancourt (see \cite[Th.\,1.1.30]{lerner-pseudo}).
For the second estimate, \cite[Th.\,1.1.31]{lerner-pseudo} states that there is a ``composition'' symbol $s(x,\xi,h)$ belonging to the class $S_{\operatorname{scl}}^0$ satisfying
\begin{equation*}
s_1(x,hD)\circ s_2(x,hD)=s(x,D,h).
\end{equation*}
Then we use the semi-classical symbolic calculus at rank $N$ of $s_1(x,h\xi)\in S_{\operatorname{scl}}^0$ and $s_2(x,h\xi)\in S_{\operatorname{scl}}^0$ (see \cite[Th.\,1.1.32]{lerner-pseudo}) that ensures that the ``composition'' symbol $s(x,\xi,h)$ satisfies
\begin{equation*}
s(x,\xi,h) - \sum_{|\alpha|<N} \frac{1}{\alpha! i^\alpha}\, \partial_\xi^{\alpha} \{s_1(x,h\xi)\} \partial_x^{\alpha}\{s_2(x,h\xi)\} \in S_{\operatorname{scl}}^{-N},
\end{equation*}
which means that $s(x,\xi,h)$ belongs to $S_{\operatorname{scl}}^{-N}$ because the partial sum
vanishes thanks to the assumption of disjoints supports.
Finally, a new use of the semi-classical Calderon-Vaillancourt (see \cite[Th.\,1.1.30]{lerner-pseudo}) proves \eqref{ca2}.
\end{proof}

We now recall the local expression of the Laplace-Beltrami operator $\Deltag$ of $\mathcal{M}$.
Let $\tau:U\subset \R^d \to V\subset \mathcal{M}$ be a coordinate patch of $\mathcal{M}$.
For any $x\in U$, we denote by $g_{x}$ the $d\times d$ positive-definite matrix corresponding to the Riemannian metric at $\tau(x)\in \mathcal{M}$ via the coordinate patch $\tau$.
Moreover, $(g_x^{ij})_{i,j}$ is the usual notation for the inverse of the matrix $g_x$.
For any smooth function $f:\mathcal{M}\to \C$ with compact support in $V$, the action of the Laplace-Beltrami operator on $f$ can be seen as that of a differential operator acting on $f\circ \tau:U\to \R$ as follows (see for instance \cite[p.\,167]{shub-pseudo}):
\begin{align*}
(\Deltag f)(\tau(x))& = \frac{1}{\sqrt{\det g_{x}}}\sum_{i=1}^d \sum_{j=1}^d \partial_{x_i} \Big(g_{x}^{ij} \sqrt{\det g_{x}} \partial_{x_j} (f\circ \tau) \Big) \\
& = \sum_{i=1}^d \sum_{j=1}^d g_x^{ij} \partial_{x_i} \partial_{x_j} (f\circ \tau) +\text{differential terms of order }1.
\end{align*}
By introducing the so-called principal symbol of $-\Deltag$ on the coordinate patch $\tau$
\begin{equation}\label{prin-sy}
\wp:(x,\xi)\in \R^d\times \R^d \mto \pun_{U}(x)\sum_{i=1}^d \sum_{j=1}^d g_x^{ij} \xi_i \xi_j\in\R
\end{equation}
and using the pseudo-differential notation \eqref{defi-opd} above (without $h$), the principal term of $(-\Deltag f)(\tau(x))$ reads
\begin{equation*}
-\sum_{i=1}^d \sum_{j=1}^d g_x^{ij}\partial_{x_i} \partial_{x_j} (f\circ \tau) =\wp(x,D)(f\circ\tau),
\end{equation*}
where $f\circ \tau$ belongs to $\mathcal{S}(\R^d)$ after being extended by $0$ outside $U$.
Moreover, on any compact subset of the domain $U$ of the coordinate patch, there is by continuity a constant $\gamma\geq 1$ for which the following uniform estimates hold
\begin{equation}\label{sym-p2}
\forall \xi\in \R^d,\quad \frac{1}{\gamma}\,|\xi|^2\leq \wp(x,\xi)\leq \gamma|\xi|^2.
\end{equation}
These inequalities implies that $-\Deltag$ is a second order elliptic operator.
We stress that~$\wp$ depends on the chosen coordinate patch $\tau:U\subset \R^d\to V\subset \mathcal{M}$ (although we merely write $\wp$ for simplicity).

In order to avoid repeating the same technical assumptions, we set the following definition.
\begin{defi}\label{defi-LM} We denote by $\Lambda(\mathcal{M})$ the set of triplets $(\tau,\psi,\widetilde{\psi})$ as follows
\begin{enumeratei}
\item $\tau:U\to V$ is a coordinate patch of $\mathcal{M}$ from an open subset $U\subset \R^d$ to an open subset $V\subset \mathcal{M}$,
\item $\psi$ belongs to $\CC_c^\infty(V)$,
\item $\widetilde{\psi}$ belongs to $\CC_c^\infty(V)$ and equals $1$ on a neighborhood of the support of $\psi$.
\end{enumeratei}
Furthermore, for any $(\tau,\psi,\widetilde{\psi})\in \Lambda(\mathcal{M})$ and any smooth function $\Psi:[0,+\infty)\to [0,1]$ satisfying \eqref{est-sig}, we define the following two symbols
\begin{align}\label{wrt2-s}
\forall (x,\xi)\in \R^d \times \R^d, \quad
s_{\tau,\psi}^{\wp,\Psi}(x,\xi)&:= \pun_{U}(x)\psi(\tau(x))\times (1-\Psi\circ\wp(x,\xi)), \\ \label{s-bul}
s_{\tau,\psi}^{\sbullet \Psi}(x,\xi)&:= \pun_{U}(x)\psi(\tau(x))\times (1-\Psi(|\xi|^2)),
\end{align}
where $\wp$ in \eqref{wrt2-s} is given in \eqref{prin-sy} (note that $s_{\tau,\psi}^{\wp,\Psi}(x,\xi)=0$ for $x\not \in U$).
\end{defi}

For the sequel, we warn the reader that the symbol $s_{\tau,\psi}^{\wp,\Psi}$ will appear by making a semi-classical expansion of $\Psi(-h^2 \Deltag)$ in \eqref{pz-con-sc} (see the proof of Proposition \ref{wrt2}).
The main drawback of $s_{\tau,\psi}^{\wp,\Psi}$ is its dependence with respect to the Laplace-Beltrami operator $\Deltag$ (more precisely to $\wp$ via the coordinate patch $\tau$).
The next result shows that the family of symbols $s_{\tau,\psi}^{\wp,\Psi}$ is equivalent, in a semi-classical sense, to the family of symbols $s_{\tau,\psi}^{\sbullet \Psi}$ whose expression \eqref{s-bul} is clearly independent of the Laplace-Beltrami operator.

\begin{prop}\label{equiv-eucl}
Let us fix a smooth function $\Psi:[0,+\infty)\to [0,1]$ satisfying \eqref{est-sig}.
For any $(\tau,\psi,\widetilde{\psi})\in \Lambda(\mathcal{M})$, there are $\alpha>0$ and $C>0$ such that
the following inequalities hold true uniformly in $F\in L^2(\R^d)$ and $h\in(0,1]$:
\begin{align}\label{equi-eu1}
\|s_{\tau,\psi}^{\wp,\Psi}(x,hD)F \|_{L_x^2(\R^d)}& \leq C\| s_{\tau,\widetilde{\psi}}^{\sbullet\Psi}(x,\alpha hD)F \|_{L_x^2(\R^d)}+\|F\|_{L^2(\R^d)}\mathcal{O}(h),\\\label{equi-eu2}
\| s_{\tau,\psi}^{\sbullet\Psi}(x,hD)F \|_{L_x^2(\R^d)} & \leq C\|s_{\tau,\widetilde{\psi}}^{\wp,\Psi}(x,\alpha hD)F \|_{L_x^2(\R^d)} +\|F\|_{L^2(\R^d)}\mathcal{O}(h).
\end{align}
\end{prop}
\begin{proof}
Let us explain the main idea for \eqref{equi-eu1}. We will decompose
\begin{equation*}
F=s_{\tau,\widetilde{\psi}}^{\sbullet\Psi}(x,\alpha hD)F+\bigl(1-s_{\tau,\widetilde{\psi}}^{\sbullet\Psi}(x,\alpha h D)\bigr)F,
\end{equation*}
which in turn will imply that $\|s_{\tau,\psi}^{\wp,\Psi}(x,hD)F \|_{L_x^2(\R^d)}$ is bounded by
\begin{equation*}
\bigl\|s_{\tau,\psi}^{\wp,\Psi}(x,hD)s_{\tau,\widetilde{\psi}}^{\sbullet\Psi}(x,\alpha hD)F \bigr\|_{L_x^2(\R^d)}+\bigl\|s_{\tau,\psi}^{\wp,\Psi}(x,hD)(1-s_{\tau,\widetilde{\psi}}^{\sbullet\Psi}(x,\alpha hD))F \bigr\|_{L_x^2(\R^d)}.
\end{equation*}
We will then conclude by applying \eqref{ca1} and \eqref{ca2}
provided that $\alpha>0$ is chosen large enough so that the two symbols
$s_{\tau,\psi}^{\wp,\Psi}(x,\xi)$ and $1-s_{\tau,\widetilde{\psi}}^{\sbullet\Psi}(x,\alpha\xi)$ have disjoint supports.

Let us now go into technical details.
We first remark that the smooth function $\pun_{U}(x)\times (\psi\circ \tau(x))$ has compact support and thus belongs to $\CC_b^\infty(\R^d)$.
Similarly, thanks to \eqref{est-sig} and \eqref{sym-p2}, one also checks
that the two symbols $\pun_{U}(x)(\psi\circ \tau(x))\times \Psi(\wp(x,\xi))$ and $\pun_{U}(x)(\widetilde{\psi}\circ \tau(x))\times \Psi(\alpha^2 |\xi|^2)$
belong to $\CC_b^\infty(\R^d\times \R^d)$.
As a consequence of \eqref{ca1}, the pseudo-differential operators $s_{\tau,\psi}^{\wp,\Psi}(x,hD)$ and $s_{\tau,\widetilde{\psi}}^{\sbullet \Psi}(x,\alpha hD)$ are well-defined and bounded on $L^2(\R^d)$.

We now turn to the choice of $\alpha$. The two symbols $s_{\tau,\psi}^{\wp,\Psi}$ and $\pun_{U}(x)(\widetilde{\psi}\circ \tau(x))\times \Psi(\alpha^2 |\xi|^2)$ are respectively supported in
\begin{equation*} U\times \bigl\{|\xi|^2\geq \sfrac{1}{\gamma}\bigr\} \quad \text{and} \quad U\times \bigl\{ |\xi|^2 \leq \sfrac{2}{\alpha^2}\bigr\}.
\end{equation*}
For $\alpha>\sqrt{2\gamma}$, the two supports are disjoint.
In order to prove \eqref{equi-eu1}, we decompose as follows for any $(x,\xi)\in \R^d\times \R^d$:
\begin{multline*}
1 = \Big(\pun_{U}(x) (\widetilde{\psi}\circ \tau(x)) \times (1-\Psi(\alpha^2 h^2 |\xi|^2))\Big)+\Big(\pun_{U}(x)
(\widetilde{\psi}\circ \tau(x))\times \Psi(\alpha^2 h^2 |\xi|^2)\Big)\\
+\Big(1-\pun_{U}(x)(\widetilde{\psi}\circ \tau(x)) \Big),
\end{multline*}
which leads, after using the quantization formula \eqref{defi-opd}, to
\begin{equation}\label{ca3}
F= s_{\tau,\widetilde{\psi}}^{\sbullet\Psi}(x,\alpha hD)F + \bigl(\pun_{U}\times (\widetilde{\psi}\circ \tau)\times \Psi(-\alpha h^2 \Delta) \bigr)F 
+ \bigl(1-\pun_{U}\times (\widetilde{\psi}\circ\tau) \bigr)F.
\end{equation}
Hence we get
\begin{align}\nonumber
\|s_{\tau,\psi}^{\wp,\Psi}(x,hD)F\|_{L_x^2(\R^d)} &\leq \|s_{\tau,\psi}^{\wp,\Psi}(x,hD) s_{\tau,\widetilde{\psi}}^{\sbullet\Psi}(x,\alpha h D) F\|_{L_x^2(\R^d)} \\ \label{ca4}
&\hspace*{5mm} + \bigl\| s_{\tau,\psi}^{\wp,\Psi}(x,hD)\bigl(\pun_{U}\times (\widetilde{\psi}\circ \tau)\times \Psi(-\alpha h^2 \Delta)\bigr) F \bigr\|_{L_x^2(\R^d)}\\ \label{ca5}
& \hspace*{5mm}+ \bigl\|s_{\tau,\psi}^{\wp,\Psi}(x,hD)\bigl((1-\pun_{U}\times (\widetilde{\psi}\circ \tau))F\bigl)\bigl\|_{L_x^2(\R^d)}.
\end{align}
By using the Calderon-Vaillancourt inequality \eqref{ca1} for $s_1=s_{\tau,\psi}^{\wp,\Psi}$, we obtain
\begin{equation*}
\|s_{\tau,\psi}^{\wp,\Psi}(x,hD) s_{\tau,\widetilde{\psi}}^{\sbullet\Psi}(x,\alpha h D) F\|_{L_x^2(\R^d)}
\leq C\|s_{\tau,\widetilde{\psi}}^{\sbullet\Psi}(x,\alpha h D) F\|_{L_x^2(\R^d)}.
\end{equation*}
The inequality \eqref{equi-eu1} can then be proved by bounding the two terms \eqref{ca4} and \eqref{ca5} by $\|F\|_{L^2(\R^d)}\mathcal{O}(h)$ thanks to \eqref{ca2} and the following two remarks:
\begin{enumeratei}
\item accordingly to the choice of $\alpha>\sqrt{2\gamma}$, the symbols $s_{\tau,\psi}^{\wp,\Psi}(x,\xi)$ and
\begin{equation*}\pun_U(x)\times (\widetilde{\psi}\circ \tau(x))\times \Psi(\alpha^2 |\xi|^2)
\end{equation*}
have disjoint supports because of the frequency variable $\xi$,
\item let $s_2(x,\xi)$ be the smooth symbol $1-\pun_{U}(x)\times (\widetilde{\psi}\circ \tau(x))$ that merely depends on~$x$. The quantization formula \eqref{defi-opd} shows that $s_2(x,hD)$ is the multiplication operator by $1-\pun_{U}(x)\times (\widetilde{\psi}\circ \tau(x))$. We furthermore note that the spatial component of the support of $s_{\tau,\psi}^{\wp,\Psi}(x,\xi)$ (see \eqref{wrt2-s}) is included in $U$ and more precisely in the support of $\psi\circ \tau$.
But the supports of $\psi\circ \tau$ and $1-\widetilde{\psi}\circ \tau$ are disjoint because $1-\widetilde{\psi}$ vanishes in a neighborhood of the support of $\psi$ (see Definition \ref{defi-LM}). In other words, the symbols $s_{\tau,\psi}^{\wp,\Psi}(x,\xi)$ and $s_2(x,\xi)$ have disjoint supports.
\end{enumeratei}

The second inequality \eqref{equi-eu2} can be proved with a similar strategy. Actually, we have to choose $\alpha>0$ so that the two symbols $s_{\tau,\psi}^{\sbullet\Psi}(x,\xi)$ and $\pun_{U}(x)\times (\widetilde{\psi}\circ \tau(x)) \times (\Psi\circ \wp(x,\alpha \xi))$ have disjoint supports.
According to \eqref{est-sig}, \eqref{sym-p2} and \eqref{s-bul}, the supports of those two symbols are respectively included in
\begin{equation*}
U\times \bigl\{|\xi|\geq 1 \bigr\} \quad \text{and}\quad U\times \bigl\{|\xi|^2\leq \sfrac{2\gamma}{\alpha^2}\bigr\}.
\end{equation*}
Here again, for the same choice of $\alpha>\sqrt{2\gamma}$, we may apply \eqref{ca1} and \eqref{ca2} to the following similar decomposition to \eqref{ca3}:
\begin{equation*}
F=s_{\tau,\widetilde{\psi}}^{\wp,\Psi}(x,\alpha h D)F+
\bigl(\pun_{U}\times \widetilde{\psi}\circ \tau \times \Psi(\wp(x,\alpha D)) \bigr)F+\bigl(1-\pun_{U}\times \widetilde{\psi}\circ \tau \bigr)F.\qedhere
\end{equation*}
\end{proof}

We are now ready to give a local but pseudo-differential reformulation of the semi-classical Salem-Zygmund condition \eqref{pz-con-sc}.

\begin{prop}\label{wrt2}
Let us fix a smooth function $\Psi:[0,+\infty)\to [0,1]$ satisfying \eqref{est-sig}.
For any $f\in L^2(\mathcal{M})$, the semi-classical Salem-Zygmund condition \eqref{pz-con-sc} is equivalent to the following condition: for any $(\tau,\psi,\widetilde{\psi})\in \Lambda(\mathcal{M})$ (according to Definition \ref{defi-LM}), we~have
\begin{equation}\label{pz-con-loc}
\int_{0}^1 \frac{\|s_{\tau,\psi}^{\wp,\Psi}(x,hD)\Big((\widetilde{\psi}f)\circ \tau\Big)\|_{L_x^2(\R^d)}}{h\sqrt{-\ln(h)}} \,dh <+\infty,
\end{equation}
where $s_{\tau,\psi}^{\wp,\Psi}$ is defined in \eqref{wrt2-s}.
\end{prop}
\begin{proof}
By compactness of $\mathcal{M}$, there is a finite open cover $\mathcal{M}=V_1\cup \dots \cup V_n$ where each open subset $V_i$ corresponds to a coordinate patch $\tau_i:U_i\subset \R^d \to V_i\subset \mathcal{M}$.
Now consider a smooth partition of unity $1=\psi_1+\dots+\psi_n$ subordinate to the previous open cover of $\mathcal{M}$.
Since the support of $\psi_i$ is a compact subset of $V_i$, one infers that, for any $g\in L^2(\mathcal{M})$, the equivalence $\|\psi_i g\|_{L^2(V_i)}\simeq \|(\psi_i g)\circ \tau_i\|_{L^2(\R^d)}$ holds with constants independent of $g$. We then easily obtain the equivalence
\begin{equation*}
\|g\|_{L^2(\mathcal{M})}\simeq \sum_{i=1}^n \|(\psi_i g)\circ \tau_i\|_{L^2(\R^d)}.
\end{equation*}
As a consequence, it is clear that \eqref{pz-con-sc} is equivalent to the following property: for any coordinate patch $\tau:U\subset \R^d\to V\subset \mathcal{M}$ and for any $\psi\in \CC_c^\infty(V)$, we have
\begin{equation}\label{pz-con-loc2}
\int_{0}^1 \frac{\bigl\|\bigl(\psi \times f-\psi \times (\Psi(-h^2 \Deltag)f) \bigr) \circ \tau \bigr\|_{L_x^2(\R^d)}}{h\sqrt{-\ln(h)}}\, dh<+\infty,
\end{equation}
with the convention $(\psi\times f)\circ \tau(x)=0$ for $x\not\in U$.
We now invoke the semi-classical functional calculus of the Laplace-Beltrami operator as stated in \cite[Prop.\,2.1 with $N=1$ and $\sigma=0$]{burq2004}.
Let $\widetilde{\psi}\in \CC_c^\infty(V)$ be a function equaling $1$ on the support of $\psi$ (in other words, $(\tau,\psi,\widetilde{\psi})$ belongs to $\Lambda(\mathcal{M})$ in Definition \ref{defi-LM}).
The semi-classical functional calculus ensures that there exists an explicit symbol $\Psi_0\in \CC_c^\infty(U\times \R^d)$ satisfying the following properties:
\begin{enumeratei}
\item the pseudo-differential $\Psi_0(x,hD)$ quantizes $\Psi_0$ as in \eqref{defi-opd} and satisfies the uniform estimates
\begin{equation}\label{semi-cl}
\bigl\|\bigl(\psi \Psi(-h^2 \Deltag) f\bigr)\circ\tau -\Psi_0(x,h D) \bigl((\widetilde{\psi} f)\circ \tau\bigr)\bigr\|_{L^2(\R^d)}\leq C h\|f\|_{L^2(\mathcal{M})},
\end{equation}
\item the symbol $\Psi_0$ has the following explicit expression:
\begin{equation}\label{sym-p}
\forall (x,\xi)\in U\times \R^d, \quad \Psi_0(x,\xi)=\psi(\tau(x))\Psi(\wp(x,\xi)).
\end{equation}
\end{enumeratei}
By plugging \eqref{semi-cl} in \eqref{pz-con-loc2} and using the finiteness of $\int_{0}^{1} \sfrac{dh}{\sqrt{-\ln(h)}}=\sqrt{\pi}$, we~see that \eqref{pz-con-loc2} is equivalent to
\begin{equation}\label{pz-con-loc3}
\int_{0}^1 \frac{\bigl\|(\psi f)\circ \tau-\Psi_0(x,hD)\bigl((\widetilde{\psi}f)\circ \tau \bigr)\bigr\|_{L_x^2(\R^d)}}{h\sqrt{-\ln(h)}}\, dh<+\infty.
\end{equation}
Remembering that $\widetilde{\psi}$ equals $1$ on the support of $\psi$, we get $\psi=\psi\widetilde{\psi}$ and so
$(\psi f)\circ \tau=(\psi\circ \tau)\times ((\widetilde{\psi}f)\circ \tau)$.
Since the pseudo-differential operator $(\pun_{U}\times \psi\circ \tau)(x,hD)$ with symbol $\pun_{U}\times \psi\circ \tau$ (independent of $\xi$) is merely the multiplication by $\pun_{U}\times \psi\circ \tau$ (see~\eqref{defi-opd}), one may look at \eqref{wrt2-s} and \eqref{sym-p} to conclude that \eqref{pz-con-loc3} reduces to \eqref{pz-con-loc}.
\end{proof}

The following result achieves the proof of Theorem \ref{indep-met} because it allows us to get rid of the Laplace-Beltrami operator $\Deltag$ and because the vector space $L^2(\mathcal{M})$ is independent of the Riemannian structure of $\mathcal{M}$.

\begin{prop}
Let us fix a smooth function $\Psi:[0,+\infty)\to [0,1]$ satisfying \eqref{est-sig}.
Then for any $f\in L^2(\mathcal{M})$, the Salem-Zygmund condition \eqref{pz-con4} is equivalent to the following condition: for any $(\tau,\psi,\widetilde{\psi})\in \Lambda(\mathcal{M})$ (according to Definition \ref{defi-LM}), we have
\begin{equation}\label{wrt-fin}
\int_{0}^1 \frac{\bigl\|s_{\tau,\psi}^{\sbullet\Psi}(x,hD)\bigl((\widetilde{\psi}f)\circ \tau \bigr) \bigr\|_{L_x^2(\R^d)}}{h\sqrt{-\ln(h)}} \,dh <+\infty,
\end{equation}
where $s_{\tau,\psi}^{\sbullet\Psi}$ is defined in \eqref{s-bul}.
\end{prop}
\begin{proof}
The key point is the following: for any $(\tau,\psi,\widetilde{\psi})\in \Lambda(\mathcal{M})$ there is $\overline{\psi}\in \CC_c^\infty(V)$ such that $(\tau,\psi,\overline{\psi})$ and $(\tau,\overline{\psi},\widetilde{\psi})$ belong to $\Lambda(\mathcal{M})$.
Indeed, there is an open subset $W\subset V$ such that $\operatorname{supp}(\psi)\subset W\subset \{\widetilde{\psi}=1\}$. So we just have to choose $\overline{\psi}\in \CC_c^\infty(W)$
that equals $1$ on a neighborhood of the compact subset $\operatorname{supp}(\psi)$.

If the finiteness conditions \eqref{wrt-fin} hold, then for $(\tau,\overline{\psi},\widetilde{\psi})\in \Lambda(\mathcal{M})$ we get
\begin{equation}\label{pz-con-fin}
\int_{0}^1 \frac{\bigl\|s_{\tau,\overline{\psi}}^{\sbullet\Psi}(x,hD)\bigl((\widetilde{\psi}f)\circ \tau \bigr) \bigr\|_{L_x^2(\R^d)}}{h\sqrt{-\ln(h)}} \,dh<+\infty.
\end{equation}
Thanks to the Calderon-Vaillancourt inequality \eqref{ca1} (as in the proof of Proposition~\ref{equiv-eucl}), we know that $\|s_{\tau,\overline{\psi}}^{\sbullet\Psi}(x, hD)((\widetilde{\psi}f)\circ \tau) \|_{L_x^2(\R^d)}$ is bounded with respect to $h\in (0,1]$ and hence the integrability of \eqref{pz-con-fin} is at $h=0$. For any $\alpha>0$, a simple linear change of variable shows
\begin{equation*}
\int_{0}^1 \frac{\bigl\|s_{\tau,\overline{\psi}}^{\sbullet\Psi}(x,\alpha hD)\bigl((\widetilde{\psi}f)\circ \tau \bigr) \bigr\|_{L_x^2(\R^d)}}{h\sqrt{-\ln(h)}} \,dh<+\infty.
\end{equation*}
By applying \eqref{equi-eu1} of Proposition \ref{equiv-eucl} with $(\tau,\psi,\overline{\psi})\in \Lambda(\mathcal{M})$, we clearly get
the finiteness conditions \eqref{pz-con-loc} of Proposition \ref{wrt2}.
A perfectly similar argument using \eqref{equi-eu2} allows us to reverse the previous implication.
Finally, Lemma \ref{wrt1} and Proposition \ref{wrt2} show the equivalence with the Salem-Zygmund condition
\eqref{pz-con4}.
\end{proof}

\appendix
\section{First moments of chi distributions}\label{chi2}

We give elementary estimates about the first moments of chi distributions.
For any $n\in \N^\star$, we set
\begin{equation*}
\chi_n:=\frac{1}{\sqrt{n}}\biggl(\sum_{k=1}^n g_{k}^2\biggr)^{\sfrac{1}{2}},
\end{equation*}
where $g_1,\dots,g_n$ are i.i.d. Gaussian random variables $\mathcal{N}_\R(0,1)$.
The equality $\E[\chi_n^2]=1$ is obvious. The following uniform inequalities also hold (with sharpness for $n=1$):
\begin{equation*}
\frac{\sqrt{2}}{\sqrt{\pi}} \leq \E[\chi_n]\leq 1 \quad\text{and} \quad 1\leq \E[\chi_n^4]\leq 3.
\end{equation*}

\begin{itemize}
\item
For the bound on $\E[\chi_n^4]$, we write
\[
\E[\chi_n^4]= \frac{1}{n^2} \sum_{\substack{1\leq k \leq n\\
1\leq \ell \leq n }} \E[g_{k}^2 g_{\ell}^2]
\]
and use the Cauchy-Schwarz inequality $\E[g_{k}^2 g_{\ell}^2] \leq \sqrt{\E[g_{k}^4]
\E[g_{\ell}^4]}=3 $. Note that the Cauchy-Schwarz inequality also shows $1=\E[\chi_n^2]^2\leq \E[\chi_n^4]$.

\item
Let us now justify the bounds of $\E[\chi_n]$. The inequality \hbox{$\E[\chi_n]\leq \sqrt{\E[\chi_n^2]}=1$} is direct. For the lower bound $\E[\chi_n]\geq \sfrac{\sqrt{2}}{\sqrt{\pi}}$, one considers 	 random variables $\ep_1,\dots,\ep_{n}$
(which are, as usual, independent of all other random variables), then one has
\begin{align*}
\sum_{k=1}^n g_k(\omega)^2 & = \E_{\omega'}\biggl[\biggl(\sum_{k=1}^n \ep_k(\omega')g_{k}(\omega)\biggr)^2 \biggr] \geq \E_{\omega'}\biggl[\Big|\sum_{k=1}^n \ep_k(\omega')g_{k}(\omega)\Big| \biggr]^2\\
\tag*{$\implies$} \E_{\omega}\left[\chi_n\right] & \geq \frac{1}{\sqrt{n}}\E_{\omega}\E_{\omega'}\biggl[\Bigl| \sum_{k=1}^n \ep_k(\omega') g_{k}(\omega)\Bigr|\biggr].
\end{align*}

We now remark that, for a fixed $\omega'$, the random variable $\omega\mto \sum_{k=1}^n \ep_k(\omega') g_{k}(\omega)$ is Gaussian and more precisely has the same distribution as $\sqrt{n}g_1$. Hence, we get
$\E[\chi_n] \geq \E[|g_1|]=\sfrac{\sqrt{2}}{\sqrt{\pi}}$.
\end{itemize}

\section{Proof of Theorem \ref{hohor} via a result of Canzani-Hanin}\label{annex-weyl}

As written above, Theorem \ref{hohor} is considered as known (see for instance \cite[lines (4),\,(5) \& (8)]{canzani2017c}).
We merely write elements of proofs for the non-specialist reader since we do not know a published reference.
The argument here relies on the published references \cite{canzani2017c,sogge2017fourier} and on a Bernstein-type inequality on the boundaryless compact Riemannian manifold $\mathcal{M}$ proved in
\cite[line (11)]{bin2004}.

We recall a few facts of Riemannian geometry. For any $y\in \mathcal{M}$, the Riemannian structure of $\mathcal{M}$ induces a canonical isomorphism between the tangent space $T_y\mathcal{M}$ and the cotangent space $T_{y}^\star \mathcal{M}$. Hence, we may endow $T_y^\star \mathcal{M}$ with a canonical inner product and a canonical Euclidean norm $\xi\mto |\xi|_{y}$.
Consequently, in the Riemannian setting, the exponential map of $\mathcal{M}$ is canonically defined on the cotangent fiber bundle.
We denote by $\operatorname{inj}(M)$ the injectivity radius of $\mathcal{M}$, namely the supremum of radii $r>0$ us such that, for any $y\in \mathcal{M}$, the restriction of the exponential map $\exp_y:T_y^\star\mathcal{M}\to M$ on the open ball $\{\xi\in T_y\mathcal{M}, |\xi|_{y}< r\}$ is a diffeomorphism.
The injectivity radius $\operatorname{inj}(M)$ is positive thanks to the compactness of $\mathcal{M}$. Moreover one has
\begin{equation*}
\exp_y(\{\xi\in T_y^\star \mathcal{M},\, |\xi|_{y}< r\})=B_{\delta_g}(y,r).
\end{equation*}
For any $\xi\in T_y^\star \mathcal{M}$ in the ball $\{|\xi|_y<\operatorname{inj}(M)\}$, the norm $|\xi|_y$ equals the distance $\delta_g(\exp_y(\xi),y)$.
In particular for any point $x\in \mathcal{M}$ satisfying $\operatorname{dist}(x,y)<\operatorname{inj}(M)$, the element $\xi=\exp_y^{-1}(x)\in T_y^\star \mathcal{M}$ is well defined and moreover satisfies the following norm equality in $T_y^\star \mathcal{M}$:
\begin{equation}\label{norm-exp}
|\exp_y^{-1}(x)|_{y}=\delta_g(x,y).
\end{equation}

In order to keep the same notations as in \cite[line (11)]{canzani2017c}, we also denote by $\rho:\R\to \R$ a Schwartz function satisfying the following properties
\begin{align}\label{rho0}
|t|\leq \frac{1}{2} \operatorname{inj}(M) &\implies \widehat{\rho}(t) = 1, \\
\nonumber
|t|\geq \operatorname{inj}(M)& \implies \widehat{\rho}(t)=0.
\end{align}
For any $(x,y)\in \mathcal{M}^2$ and $\lambda \geq 0$, we also write the spectral function
\begin{equation*}
e_{[0,\lambda]}(x,y)=\sum_{\lambda_j\leq \lambda} \phi_j(x)\phi_j(y).
\end{equation*}
We now decompose the derivatives of $e_{[0,\lambda]}(x,y)$ via a coordinate patch $\tau:U\subset\R^d\to V\subset \mathcal{M}$ as follows:
\begin{equation}\label{decomp}
\partial_x^{I} \partial_y^J e_{[0,\lambda]}(x,y)=\rho \star \partial_x^{I} \partial_y^J e_{[0,\lambda]}(x,y) +\bigl[\partial_x^{I} \partial_y^J e_{[0,\lambda]}(x,y)-\rho \star \partial_x^{I} \partial_y^J e_{[0,\lambda]}(x,y)\bigr],
\end{equation}
where the convolution
has to be understood with respect to $\lambda$.
We now invoke the part of a work of Canzani-Hanin that holds true without any geometric property on the manifold. In local coordinates, Lemma 5 of \cite{canzani2017c} (with $Q=\operatorname{Id}$, $D_0^Q=1$ and $D_{-1}^Q=0$) implies the following statement.

\begin{prop}[Canzani-Hanin]\label{can-han} Let $\mathcal{M}$ be a boundaryless compact Riemannian manifold of dimension $d\geq 2$ and let us consider a coordinate patch $\tau:U\subset\R^d\to V\subset\mathcal{M}$ with $\operatorname{diam}(V)\leq \frac{1}{2}\operatorname{inj}(M)$.
For any multi-indexes $I\in \N^{d}$ and $J\in \N^d$, the following asymptotics holds true for any $(x,y)\in V^2$ and any $\nu\geq 1$ (where the spatial derivatives $\partial_x$ and $\partial_y$ are seen in the coordinate patch and $\partial_\nu$ is seen in $\R$):
\begin{multline*}
\partial_\nu (\rho \star \partial_x^I \partial_y^J e_{[0,\nu]})(x,y)=\frac{\nu^{d-1}}{(2\pi)^d} \partial_x^{I} \partial_y^{J} \biggl\{\int_{|\xi|_{y}=1} e^{i \nu\langle \exp_y^{-1}(x),\xi\rangle }ds(\xi)\biggr\}\\
+\mathcal{O}\bigl(\nu^{d-2+|I|+|J|}\delta_g(x,y)+(1+\nu)^{d-3+|I|+|J|} \bigr),
\end{multline*}
where $ds$ denotes the measure of the unit sphere of
$T_y^\star \mathcal{M}$ (canonically induced by the inner product on the cotangent space $T_y^\star\mathcal{M}$). Finally, the remainder is uniform provided that $x$ and $y$ run over a compact subset of the open set $V$.
\end{prop}
The presence of the distance $\delta_g(x,y)$ in the remainder is of interest in \cite{canzani2017c} and also in \cite[Prop.\,10]{canzani2015scaling} because $\delta_g(x,y)$ may become very small in the previous papers. However, by using the fact that the Riemannian distance is bounded on the compact manifold $\mathcal{M}$, we may integrate on $\nu\in [0,\lambda]$ with polar coordinates to get
\begin{align*}
\rho\star \partial_x^I \partial_y^J e_{[0,\lambda]}(x,y)&=\frac{\partial_x^I \partial_y^J}{(2\pi)^d}\biggl\{ \int_{0}^\lambda\! \biggl(\int_{|\xi|_{y}=1}\hspace*{-5mm} e^{i\langle \exp_y^{-1}(x),\nu \xi\rangle} \nu^{d-1} ds(\xi)\biggr) d\nu \biggr\}+\mathcal{O}(\lambda^{d-1+|I|+|J|}) \\
&=\frac{\partial_x^I \partial_y^J}{(2\pi)^d}\biggl\{ \int_{|\xi|_{y}<\lambda}\hspace*{-5mm} e^{i\langle \exp_y^{-1}(x),\xi\rangle} \frac{ d\xi }{\sqrt{|g_y|}} \biggr\}+\mathcal{O}(\lambda^{d-1+|I|+|J|}),
\end{align*}
where $\sfrac{d\xi}{\sqrt{|g_y|}}$ is the usual notation for the Lebesgue measure associated to the inner product of the cotangent space $T_y^\star \mathcal{M}$. With the formulas \eqref{norm-exp}, \eqref{defi-bes} and \cite[lines (8--9)]{canzani2017c}, we recover the principal term of the statement of Theorem \ref{hohor}:
\begin{equation*}
\rho\star \partial_x^I \partial_y^J e_{[0,\lambda]}(x,y)
= \frac{\partial_x^{I}\partial_y^J}{(2\pi)^{\sfrac{d}{2}}}\biggl\{\int_{0}^{\lambda }\nu^{d-1} \widehat{\sigma_{d-1}}(\nu \delta_g(x,y))
d\nu\biggr\} +\mathcal{O}(\lambda^{d-1+|I|+|J|}).
\end{equation*}

To complete the proof of Theorem \ref{hohor}, we merely need to look at the decomposition \eqref{decomp} and the following result (still written in local coordinates).

\begin{prop} Let $\mathcal{M}$ be a boundaryless compact Riemannian manifold of dimension $d\geq 2$
and let us consider a coordinate patch $\tau:U\subset\R^d\to V\subset\mathcal{M}$.
For any two multi-indexes $I\in \N^{d}$ and $J\in \N^d$, the following inequalities hold true for any $x$ and $y$ belonging to a compact subset of the open set $V$:
\begin{equation}\label{taub}
\forall \lambda\geq 1, \quad \bigl|\rho \star \partial_x^I \partial_y^J e_{[0,\lambda]}(x,y)-\partial_x^I \partial_y^J e_{[0,\lambda]}(x,y)\bigr|\leq C\lambda^{d-1+|I|+|J|}.
\end{equation}
\end{prop}
\begin{proof}
Let us convene that $e_{[0,\lambda]}(x,y)$ equals $0$ for $\lambda<0$.
The Fourier transform of $\lambda\in \R \mto \rho \star \partial_x^I \partial_y^J e_{[0,\lambda]}(x,y)-\partial_x^I \partial_y^J e_{[0,\lambda]}(x,y)$ can be factorized by $\widehat{\rho}(t)-1$ and hence vanishes near $0$ thanks to \eqref{rho0}.
Hence, a standard Tauberian lemma (see \cite[Lem.\,4.2.3, lines (4.2.13) and (4.2.14)]{sogge2017fourier}) ensures that \eqref{taub} will be a consequence of the following estimates for all $\lambda \geq 0$ and $s\in [0,1]$
\begin{equation*}
\bigl|\rho \star \partial_x^I \partial_y^J e_{(\lambda,\lambda+s]}(x,y)-\partial_x^I \partial_y^J e_{(\lambda,\lambda+s]}(x,y)\bigr|\leq C(1+\lambda)^{d-1+|I|+|J|}.
\end{equation*}
Clearly, it is sufficient to prove the following two inequalities
\begin{align}\label{taub1}
|\partial_x^I \partial_y^J e_{(\lambda,\lambda+s]}(x,y)|&\leq C(1+\lambda)^{d-1+|I|+|J|}, \\ \label{taub2}
|\rho \star \partial_x^I \partial_y^J e_{(\lambda,\lambda+s]}(x,y)|&\leq C(1+\lambda)^{d-1+|I|+|J|}.
\end{align}
Let us prove \eqref{taub1}. The Cauchy-Schwarz inequality gives
\begin{align*}
|\partial_x^I \partial_y^J e_{(\lambda,\lambda+s]}(x,y)| & = \biggl\vert \sum_{\lambda<\lambda_j\leq \lambda +s} \partial_x^I \phi_j(x) \partial_y^J \phi_j(y) \biggr\vert \\
& \leq \sqrt{\sum_{\lambda <\lambda_j\leq \lambda+s} |\partial_x^I \phi_j(x)|^2} \sqrt{\sum_{\lambda <\lambda_j\leq \lambda+s} |\partial_y^J \phi_j(y)|^2}.
\end{align*}
We then get \eqref{taub1} by invoking the following Bernstein-type inequality proved by Bin (see \cite[line (11)]{bin2004}):
\begin{equation}\label{taub1a}
\sqrt{\sum_{\lambda \leq \lambda_j\leq \lambda+1} |\partial_x^I \phi_j(x)|^2}\leq C(1+\lambda)^{\psfrac{d-1}{2}+|I|}.
\end{equation}
Note actually that the results of \cite{bin2004} are stated in specific coordinate patches that are usually called normal coordinates. Here are some details explaining why \eqref{taub1a} is still true for other coordinates. Let us consider normal coordinates
\[
\tau_{\mathcal{N}}:U_{\mathcal{N}}\subset \R^d \to V\subset \mathcal{M}
\]
(see \cite{bin2004} for definitions).
Since $\tau_{\mathcal{N}}^{-1}\circ \tau:U\to U_{\mathcal{N}}$ is a diffeomorphism, all its derivatives are bounded on any compact subset of $U$.
Given an eigenfunction \hbox{$\phi_j:\mathcal{M}\to \R$}, it is clear that bounding a finite number of derivatives of $\phi_j\circ \tau:U\to \R$ or $\phi_j\circ \tau_{\mathcal{N}}:U_{\mathcal{N}}\to \R$ are equivalent problems since, by decomposing $\phi_j\circ \tau=(\phi_j\circ \tau_{\mathcal{N}})\circ (\tau_{\mathcal{N}}^{-1}\circ \tau)$ and using the Faà di Bruno's formula, we have for any $k\in \N$ and $x\in V$ an inequality of the form
\begin{equation*}
\sum_{|I|\leq k} \bigl(\partial^I (\phi_j\circ \tau)|_{\tau^{-1}(x)}\bigr)^2 \leq C
\sum_{|I| \leq k} \bigl(\partial^I (\phi_j\circ \tau_{\mathcal{N}})|_{\tau_{\mathcal{N}}^{-1}(x)}\bigr)^2
\end{equation*}
and the last inequality is moreover uniform for $x$ running over a compact subset of~$V$.
Hence the inequality \eqref{taub1a} of \cite{bin2004} still holds for non-normal coordinates.

Let us turn to the proof of \eqref{taub2}. We begin by writing the finite sums
\begin{align*}
\rho \star \partial_x^I \partial_y^J e_{(\lambda,\lambda+s]}(x,y) & = \rho \star \biggl(\sum_{j\in \N} \partial_x^I \phi_j(x) \partial_y^J \phi_j (y) \pun_{[\lambda_j-s,\lambda_j)}(\lambda) \biggr)\\
& = \sum_{j\in \N} \partial_x^I \phi_j(x)\partial_y^J\phi_j(y) \int_{\lambda-\lambda_j}^{\lambda-\lambda_j+s} \rho(\nu)d\nu .
\end{align*}
Using that $\rho$ belongs to the Schwartz space and the fact that the length of the interval $[\lambda-\lambda_j,\lambda-\lambda_j+s]$ is less or equal to $1$, we obtain the following upper bound
\begin{align*}
|\rho \star \partial_x^I \partial_y^J e_{(\lambda,\lambda+s]}(x,y)| & \leq  C\sum_{j\in \N} \frac{|\partial_x^{I} \phi_j(x)||\partial_y^J \phi_j(y)|}{(1+|\lambda-\lambda_j|)^{d+1+|I|+|J|}}\\
& \leq C \sqrt{\sum_{j\in \N} \frac{|\partial_x^{I} \phi_j(x)|^2}{(1+|\lambda-\lambda_j|)^{d+1+2|I|}}}
\sqrt{\sum_{j\in \N} \frac{|\partial_y^{J} \phi_j(y)|^2}{(1+|\lambda-\lambda_j|)^{d+1+
2|J|}}}.
\end{align*}
We finish by bounding the first square root (the other is obviously similar):
\begin{align*}
\sum_{j\in \N} \frac{|\partial_x^{I} \phi_j(x)|^2}{(1+|\lambda-\lambda_j|)^{d+1+2|I|}} &=
\sum_{k\in \N} \sum_{|\lambda_j-\lambda|\in [k,k+1)} \frac{|\partial_x^{I} \phi_j(x)|^2}{(1+|\lambda-\lambda_j|)^{d+1+2|I|}} \\
& \leq \sum_{k\in \N}\frac{1}{(1+k)^{d+1+2|I|}} \sum_{|\lambda_j-\lambda|\in [k,k+1]}|\partial_x^{I} \phi_j(x)|^2.
\end{align*}
In the last sum, $\lambda_j$ belongs to $[\lambda-k-1,\lambda-k]\cup[\lambda+k,\lambda+k+1]$ and we may use \eqref{taub1a} to get
\begin{align*}
\sum_{j\in \N} \frac{|\partial_x^{I} \phi_j(x)|^2}{(1+|\lambda-\lambda_j|)^{d+1+2|I|}} &\leq C \sum_{k\in \N}\frac{(1+\lambda+k)^{d-1+2|I|}}{(1+k)^{d+1+2|I|}}\\
& \leq C (1+\lambda)^{d-1+2|I|}.\qedhere
\end{align*}
\end{proof}

\backmatter
\bibliographystyle{jepalpha+eid}
\bibliography{imekraz}
\end{document}
